\chapter{Merkle commitment scheme}
\label{chapter:merkle-commitment}

We describe the \emph{Merkle commitment scheme}, a commitment scheme in the ROM that produces succinct commitments to long lists of values and enables cheaply opening particular subsets of values in the list. This useful \DoQuote{succinct opening} property is used to construct succinct arguments in the ROM.\footnote{More generally, succinct commitments with succinct opening are known as \emph{vector commitments} in the cryptography literature, and constructions are known from various cryptographic assumptions. They are used to construct succinct arguments in various settings also beyond the (pure) ROM.} The commitment scheme is so-called because it relies on Merkle trees \cite{Merkle89-tree}, a powerful and versatile data structure based on hash functions. A Merkle commitment scheme uses the basic commitment scheme $\CMSymbol$ as a building block, so throughout this chapter we assume familiarity with the material in \Cref{chapter:basic-commitment}.

\parhead{Organization}
In \Cref{section:merkle-commitment-definition} we describe the construction of a Merkle commitment scheme, and discuss its efficiency properties. Then in \Cref{section:merkle-commitment-completeness} we discuss its \emph{completeness property}, that is, how an honest use of the commitment scheme always leads to successful validity checks. After that we discuss security properties. In \Cref{section:merkle-commitment-collision-lemma} we establish a technical property that we use in several security analyses. After that we discuss the main security properties of a Merkle commitment scheme:
\begin{itemize}[noitemsep]
  \item in \Cref{section:merkle-commitment-binding} we prove that $\MTSymbol$ is \emph{binding};
  \item in \Cref{section:merkle-commitment-extractability} we prove that $\MTSymbol$ is \emph{extractable};
  \item in \Cref{section:merkle-commitment-hiding} we prove that $\MTSymbol$ is \emph{hiding};
  \item in \Cref{section:merkle-commitment-equivocation} we prove that $\MTSymbol$ is \emph{equivocable}.
\end{itemize}

%%%%%%%%%%%%%%%%%%%%%%%%%%%%%%%%%%%%%%%%%%%%%%%%%%%%%%%%%%%%%%%%%%%%%%%%%%%%%%%
%%%%%%%%%%%%%%%%%%%%%%%%%%%%%%%%%%%%%%%%%%%%%%%%%%%%%%%%%%%%%%%%%%%%%%%%%%%%%%%
%%%%%%%%%%%%%%%%%%%%%%%%%%%%%%%%%%%%%%%%%%%%%%%%%%%%%%%%%%%%%%%%%%%%%%%%%%%%%%%
\section{Definition}
\label{section:merkle-commitment-definition}

A \emph{Merkle commitment scheme} is a tuple of three algorithms:
\begin{equation*}
\MTSymbol = (\MTCommit,\MTOpen,\MTCheck)
\enspace.
\end{equation*}
The scheme has several parameters: an output size $\ROOutputSize \in \Naturals$ of the random oracle, a message alphabet $\MTAlphabet$, a message length $\MTMessageLength$, and a salt size $\MTSaltSize \in \Naturals$. We use the notation $\MTConstructor{\ROOutputSize}{\MTAlphabet}{\MTMessageLength}{\MTSaltSize}$ when we wish to make these parameters explicit.

The algorithms receive query access to a random oracle $\ROFunction \in \RODistribution{\ROOutputSize}$ and have the following syntax.
\begin{itemize}
  \item \emph{Commit.}
  The algorithm $\MTCommit$ receives as input a message vector $\MTMessageVector \in \MTAlphabet^{\MTMessageLength}$, and computes a Merkle commitment $\MTCommitment \in \Bits^{\ROOutputSize}$ and corresponding opening trapdoor $\MTTrapdoor \in \Bits^{O((\MTSaltSize + \ROOutputSize) \cdot \MTMessageLength)}$. This algorithm is probabilistic if privacy is desired (otherwise it is deterministic).
  \item \emph{Open.}
  The algorithm $\MTOpen$ receives as input opening trapdoor $\MTTrapdoor$ and a subset $\MTMessageIndexSet \subseteq [\MTMessageLength]$, and computes an opening proof $\MTProof$ that authenticates the values at the locations in $\MTMessageIndexSet$.
  \item \emph{Check.}
  The algorithm $\MTCheck$ receives as input a Merkle commitment $\MTCommitment$, subset $\MTMessageIndexSet \subseteq [\MTMessageLength]$, claimed values $\MTMessageSubVector \in \MTAlphabet^{\MTMessageIndexSet}$, and opening proof $\MTProof$, and computes a bit indicating whether the opening proof $\MTProof$ authenticates $\MTMessageSubVector$ as values for the locations in $\MTMessageIndexSet$ with respect to $\MTCommitment$.
\end{itemize}
Henceforth we assume that $\MTMessageLength$ is a power of $2$; see \Cref{remark:mt-any-message-length} for a discussion of how to remove this assumption.

\parhead{Binary tree}
We introduce notation for a (perfect) binary tree on $\MTMessageLength$ leaves,\footnote{A binary tree is \emph{perfect} if all internal vertices have two children and all leaf vertices belong to the same layer.} which then we use to describe the algorithms of the Merkle commitment scheme. Those algorithms can be viewed as labeling (or checking the labels of) vertices of this (perfect) binary tree.

\begin{definition}
\label{definition:mt-graph}
For $\MTMessageLength \in \Naturals$, $\MTGraph_{\MTMessageLength} = (\MTVertexSet_{\MTMessageLength},\MTEdgeSet_{\MTMessageLength})$ is the (perfect) \defemph{binary tree graph} of depth
\begin{equation}
\label{equation:mt-depth-definition}
\MTDepth \DefineEqual \log_{2} \MTMessageLength
\end{equation}
where the vertex set $\MTVertexSet_{\MTMessageLength}$ and edge set $\MTEdgeSet_{\MTMessageLength}$ are defined as follows:
\begin{align*}
\MTVertexSet_{\MTMessageLength} &\DefineEqual \left\{ (\MTLayer,\MTMessageIndex) : \MTLayer \in \{0,1,\dots,\MTDepth\} , \MTMessageIndex \in [2^{\MTLayer}] \right\} \enspace, \\
\MTEdgeSet_{\MTMessageLength} &\DefineEqual \left\{ \big((\MTLayer-1,\MTMessageIndex),(\MTLayer,2\MTMessageIndex-b)\big) : \MTLayer \in \{1,\dots,\MTDepth\} , \MTMessageIndex \in [2^{\MTLayer-1}] , b \in \{1,0\} \right\} \enspace.
\end{align*}
\end{definition}

We use the depth $\MTDepth$ throughout this chapter without explicitly re-defining it as in \Cref{equation:mt-depth-definition}.

\begin{definition}
\label{definition:mt-paths}
We make the following notational definitions:
\begin{itemize}[nolistsep]
  \item The \defemph{root vertex} of $\MTGraph_{\MTMessageLength}$ is the vertex $(0,1)$.
  \item The \defemph{leaf vertices} of $\MTGraph_{\MTMessageLength}$ are the vertices $\{(\MTDepth,\MTMessageIndex)\}_{\MTMessageIndex \in [\MTMessageLength]}$.
  \item A vertex $(\MTLayer,\MTMessageIndex)$ is \defemph{odd} if $\MTMessageIndex$ is odd and it is \defemph{even} if $\MTMessageIndex$ is even.
  \item The \defemph{sibling} of a non-root vertex $(\MTLayer,\MTMessageIndex)$ (i.e., with $\MTLayer > 0$) is $(\MTLayer,\MTMessageIndex+1)$ if $(\MTLayer,\MTMessageIndex)$ is odd and is $(\MTLayer,\MTMessageIndex-1)$ if $(\MTLayer,\MTMessageIndex)$ is even.
  \item The \defemph{path} from the $\MTMessageIndex$-th leaf vertex $(\MTDepth,\MTMessageIndex)$ to the root vertex is denoted by $\MTPath(\MTMessageIndex)$; the vertex in $\MTPath(\MTMessageIndex)$ that is in layer $\MTLayer$ is denoted $\MTPathIndex{\MTMessageIndex}{\MTLayer} \in \{\MTLayer\} \times [2^{\MTLayer}]$.
  \item The \defemph{copath} from the $\MTMessageIndex$-th leaf vertex $(\MTDepth,\MTMessageIndex)$ to the root vertex is denoted by $\MTCoPath(\MTMessageIndex)$, and is the list of siblings of each vertex in $\MTPath(\MTMessageIndex)$, except the root vertex (which has no siblings); the vertex in $\MTCoPath(\MTMessageIndex)$ that is in layer $\MTLayer$ is denoted $\MTPathCoIndex{\MTMessageIndex}{\MTLayer} \in \{\MTLayer\} \times [2^{\MTLayer}]$ (and is the sibling of $\MTPathIndex{\MTMessageIndex}{\MTLayer}$).
  \item For $\MTMessageIndexSet \subseteq [\MTMessageLength]$, $\MTPath(\MTMessageIndexSet) \DefineEqual \idxcup{\MTMessageIndex \in \MTMessageIndexSet} \MTPath(\MTMessageIndex)$ and $\MTCoPath(\MTMessageIndexSet) \DefineEqual \idxcup{\MTMessageIndex \in \MTMessageIndexSet} \MTCoPath(\MTMessageIndex)$ (viewing lists as sets).
\end{itemize}
\end{definition}

Using this notation we can write:
\begin{equation*}
\MTPath(\MTMessageIndex)
=
\Big(
\MTPathIndex{\MTMessageIndex}{\MTLayer}
\Big)_{\MTLayer \in \{0,1,\dots,\MTDepth\}}
\TextAndInMath
\MTCoPath(\MTMessageIndex)
=
\Big(
\MTPathCoIndex{\MTMessageIndex}{\MTLayer}
\Big)_{\MTLayer \in \{1,\dots,\MTDepth\}}
\enspace.
\end{equation*}
For every $\MTMessageIndex \in [\MTMessageLength]$, $\MTPathIndex{\MTMessageIndex}{0} = (0,1)$ is the root vertex and $\MTPathIndex{\MTMessageIndex}{\MTDepth} = (\MTDepth,\MTMessageIndex)$ is the $\MTMessageIndex$-th leaf.

\begin{figure}[htp!]
\centering
\includegraphics[width=0.75\textwidth]{\FigureFolder/tree-diagram}
\caption[The binary tree graph $\MTGraph_{8}$.]{Diagram of the binary tree graph $\MTGraph_{8}$. The vertices highlighted in yellow are $\MTPath(6) = \big( (3,6),(2,3),(1,2),(0,1) \big)$ and the vertices highlighted in blue are $\MTCoPath(6) = \big( (3,5),(2,4),(1,1) \big)$. For example, $\MTPathIndex{6}{1}=(1,2)$ and $\MTPathCoIndex{6}{1}=(1,1)$.}
\label{figure:tree-diagram}
\end{figure}

\parhead{Construction}
We describe each of the algorithms in the Merkle commitment scheme $\MTSymbol$.
\begin{itemize}[leftmargin=*]

  \item $\MTCommit^{\ROFunction}(\MTMessageVector) \to (\MTCommitment,\MTTrapdoor)$
  \begin{enumerate}[nolistsep]
    \item \label{step:mt-sample-salts}
    For each message location $\MTMessageIndex = 1,\dots,\MTMessageLength$:
    \begin{itemize}[nolistsep]
      \item sample a salt $\MTSaltString_{\MTMessageIndex} \in \Bits^{\MTSaltSize}$;
      \item compute the (hiding) commitment $\MTVertexLabelByName{\MTDepth}{\MTMessageIndex} \DefineEqual \ROFunction(\MTMessageVector[\MTMessageIndex],\MTSaltString_{\MTMessageIndex}) \in \Bits^{\ROOutputSize}$.
    \end{itemize}
    \item For each layer $\MTLayer = \MTDepth-1,\dots,0$ and for each $\MTMessageIndex = 1,\dots,2^{\MTLayer}$:
    \begin{itemize}[nolistsep]
      \item compute the (non-hiding) commitment $\MTVertexLabelByName{\MTLayer}{\MTMessageIndex} \DefineEqual \ROFunction(\MTVertexLabelByName{\MTLayer+1}{2\MTMessageIndex-1},\MTVertexLabelByName{\MTLayer+1}{2\MTMessageIndex}) \in \Bits^{\ROOutputSize}$.
    \end{itemize}
    \item Set the Merkle commitment $\MTCommitment \DefineEqual \MTVertexLabelByName{0}{1}$.
    \item Set the salts $\MTSaltStrings \DefineEqual (\MTSaltString_{\MTMessageIndex})_{\MTMessageIndex \in [\MTMessageLength]}$.
    \item Set the commitments $\MTVertexLabels \DefineEqual ((\MTVertexLabelByName{\MTLayer}{\MTMessageIndex})_{\MTMessageIndex \in [2^{\MTLayer}]})_{\MTLayer=0}^{\MTDepth}$, which we can visualize in the following way:
\begin{equation*}
\MTVertexLabels
=
\begin{bmatrix}
    \MTVertexLabelByName{0}{1} & & & & & \\
    \MTVertexLabelByName{1}{1} & \MTVertexLabelByName{1}{2} & & & & \\
    \MTVertexLabelByName{2}{1} & \MTVertexLabelByName{2}{2} & \MTVertexLabelByName{2}{3} & \MTVertexLabelByName{2}{4} & & \\
    \vdots & \vdots & \vdots & \vdots & \ddots & \\
    \MTVertexLabelByName{\MTDepth}{1} & \dots & \dots & \dots & \dots & \MTVertexLabelByName{\MTDepth}{2^{\MTDepth}}
\end{bmatrix}
\enspace.
\end{equation*}
    \item Set the opening trapdoor $\MTTrapdoor \DefineEqual (\MTSaltStrings,\MTVertexLabels)$.
    \item Output $(\MTCommitment,\MTTrapdoor)$.
  \end{enumerate}

  Informally, this algorithm first commits to each message entry via a basic commitment scheme (the one discussed in \Cref{chapter:basic-commitment}) and labels the leaves of the binary tree with these commitments, and then recursively labels every vertex in the tree with the hash of the labels of its two children.

  \item $\MTOpen^{\ROFunction}(\MTTrapdoor, \MTMessageIndexSet) \to \MTProof$
  \begin{enumerate}[nolistsep]
    \item For every $\MTMessageIndex \in \MTMessageIndexSet$, set the authentication path for location $\MTMessageIndex$:
\begin{equation}
\label{equation:mt-authpath-structure}
\MTAuthPath_{\MTMessageIndex}
\DefineEqual
\big(
  \MTSaltString_{\MTMessageIndex}
  ,
  (\MTVertexLabelByCoPath{\MTMessageIndex}{\MTLayer})_{\MTLayer \in \{1,\dots,\MTDepth\}}
\big)
\enspace.
\end{equation}
    \item Output the opening proof $\MTProof \DefineEqual (\MTAuthPath_{\MTMessageIndex})_{\MTMessageIndex \in \MTMessageIndexSet}$.
  \end{enumerate}

  Informally, this algorithm includes an authentication path for each opened location. An authentication for a location includes the salt for that location, as well as the values assigned to siblings of vertices from that leaf location to the root.

  \item $\MTCheck^{\ROFunction}(\MTCommitment,\MTMessageIndexSet,\MTMessageSubVector,\MTProof) \to b$
  \begin{enumerate}[nolistsep]
    \item Parse $\MTProof$ as $(\MTAuthPath_{\MTMessageIndex})_{\MTMessageIndex \in \MTMessageIndexSet}$.
    \item For every $\MTMessageIndex \in \MTMessageIndexSet$, check that $\MTAuthPath_{\MTMessageIndex}$ authenticates the value $\MTMessageSubVector[\MTMessageIndex] \in \MTAlphabet$ as the $\MTMessageIndex$-th opening relative to the Merkle commitment $\MTCommitment \in \Bits^{\ROOutputSize}$ by computing $(\MTVertexLabelByPath{\MTMessageIndex}{\MTLayer})_{\MTLayer \in \{0,1,\dots,\MTDepth\}}$:
    \begin{enumerate}[nolistsep]
      \item \label{step:mtcheck:leaf-query}
      compute the commitment $\MTVertexLabelByName{\MTDepth}{\MTMessageIndex} \DefineEqual \ROFunction(\MTMessageSubVector[\MTMessageIndex],\MTSaltString_{\MTMessageIndex})$;
      \item compute the commitments $(\MTVertexLabelByPath{\MTMessageIndex}{\MTLayer})_{\MTLayer \in \{0,1,\dots,\MTDepth-1\}}$ as follows:\\ for each layer $\MTLayer = \MTDepth-1,\dots,1,0$:
      \begin{itemize}[nolistsep]
        \item if $\MTPathIndex{\MTMessageIndex}{\MTLayer+1}$ is an odd vertex then set $\MTLeftVertexLabel \DefineEqual \MTVertexLabelByPath{\MTMessageIndex}{\MTLayer+1}$ and $\MTRightVertexLabel \DefineEqual \MTVertexLabelByCoPath{\MTMessageIndex}{\MTLayer+1}$;
        \item if $\MTPathIndex{\MTMessageIndex}{\MTLayer+1}$ is an even vertex then set $\MTLeftVertexLabel \DefineEqual \MTVertexLabelByCoPath{\MTMessageIndex}{\MTLayer+1}$ and $\MTRightVertexLabel \DefineEqual \MTVertexLabelByPath{\MTMessageIndex}{\MTLayer+1}$;
        \item \label{step:mtcheck:internal-query}
        compute the commitment $\MTVertexLabelByPath{\MTMessageIndex}{\MTLayer} \DefineEqual \ROFunction(\MTLeftVertexLabel,\MTRightVertexLabel)$.
      \end{itemize}
      \item check that $\MTCommitment = \MTVertexLabelByName{0}{1}$.
    \end{enumerate}
  \end{enumerate}

  Informally, this algorithm checks the authentication path from each location. Checking an authentication path involves re-computing the values assigned to vertices from the leaf vertex of the location to the root vertex, and checking that the resulting root value equals the given Merkle commitment. Note that
\begin{equation}
\label{equation:mt-check-and-for-each-location}
\MTCheck^{\ROFunction}(\MTCommitment,\MTMessageIndexSet,\MTMessageSubVector,\MTProof)
=
\idxland{\MTMessageIndex \in \MTMessageIndexSet}\, \MTCheck^{\ROFunction}(\MTCommitment,\MTMessageIndex,\MTMessageSubVector[\MTMessageIndex],\MTAuthPath_{\MTMessageIndex})
\enspace.
\end{equation}
Above we slightly abuse notation in that we drop the set notation when $\MTCheck$ receives a single location and a single corresponding authentication path.

\end{itemize}

\parhead{Efficiency}
We discuss the efficiency of each algorithm.
\begin{itemize}

  \item The algorithm $\MTCommit$ makes $\MTMessageLength$ queries of size $\log\Cardinality{\MTAlphabet}+\MTSaltSize$, and $\sum_{\MTLayer \in [\MTDepth]} \MTMessageLength/2^{\MTDepth-\MTLayer-1} \leq \MTMessageLength$ queries of size $2\ROOutputSize$. The output Merkle commitment $\MTCommitment$ has size $\ROOutputSize$ and opening trapdoor $\MTTrapdoor$ has size $\MTSaltSize \MTMessageLength + 2\MTMessageLength\ROOutputSize$.

  \item The algorithm $\MTOpen$ makes no oracle calls, and simply includes in the opening proof an authentication path for each index in the set $\MTMessageIndexSet$. An authentication path has size $\MTSaltSize + \MTDepth\ROOutputSize$, since it includes the salt and a hash per layer. Hence an opening proof consists of $\Cardinality{\MTMessageIndexSet} \cdot (\MTSaltSize + \MTDepth\ROOutputSize)$ bits. The path-pruning optimization discussed in \Cref{section:compressing-opening-proof} reduces the size of an opening proof.

  \item The algorithm $\MTCheck$, for each location in the set $\MTMessageIndexSet$, makes one query of size $\log\Cardinality{\MTAlphabet}+\MTSaltSize$ and $\MTDepth$ queries of size $2\ROOutputSize$. The total number of queries is $\Cardinality{\cup_{\MTMessageIndex \in \MTMessageIndexSet}\MTPath(\MTMessageIndex)}$ (there may be duplicates).

\end{itemize}

\begin{figure}[htp!]
\centering
\includegraphics[width=0.7\textwidth]{\FigureFolder/mt-diagram}
\caption[Computation of $\MTCommit$ for $\MTMessageLength=8$.]{Diagram of the computation in $\MTCommit$ to produce a Merkle commitment for a message of length $\MTMessageLength=8$. Moreover, highlighted in blue is the information included in an opening proof $\MTProof$ for the $6$-th entry of the message (which is highlighted in yellow): the salt $\MTSaltString_{6}$ and the commitments $(\MTVertexLabel_{(3,5)},\MTVertexLabel_{(2,4)},\MTVertexLabel_{(1,1)})$ that correspond to the copath $\MTCoPath(6) = \big( (3,5),(2,4),(1,1) \big)$.}
\label{figure:mt-diagram}
\end{figure}

\begin{remark}[no privacy]
\label{remark:mt-no-privacy}
If no privacy is desired then the construction of $\MTSymbol$ simplifies as follows:
\begin{inparaenum}[(i)]
  \item $\MTSaltSize = 0$ (no salts are sampled);
  \item in $\MTCommit$, set $\MTVertexLabelByName{\MTDepth}{\MTMessageIndex} \DefineEqual \MTMessageVector[\MTMessageIndex]$ for each $\MTMessageIndex \in [\MTMessageLength]$ (instead of committing first by setting $\MTVertexLabelByName{\MTDepth}{\MTMessageIndex} \DefineEqual \ROFunction(\MTMessageVector[\MTMessageIndex],\MTSaltString_{\MTMessageIndex})$) and then the opening trapdoor $\MTTrapdoor$ does not include any salts;
  \item in $\MTOpen$, each authentication path $\MTAuthPath_{\MTMessageIndex}$ does not include any salt;
  \item in $\MTCheck$, similarly to $\MTCommit$, simply set $\MTVertexLabelByName{\MTDepth}{\MTMessageIndex} \DefineEqual \MTMessageVector[\MTMessageIndex]$ for each $\MTMessageIndex \in \MTMessageIndexSet$ (in \Cref{step:mtcheck:leaf-query}).
\end{inparaenum}
The efficiency measures are correspondingly reduced by setting $\MTSaltSize=0$.
\end{remark}

\begin{remark}[recomputing the tree]
\label{remark:recomputing-the-tree}
Some information from $\MTTrapdoor$ can be omitted (thereby reducing its size) at the expense of additional computation by $\MTOpen$. Specifically, the commitments $((\MTVertexLabelByName{\MTLayer}{\MTMessageIndex})_{\MTMessageIndex \in [2^{\MTLayer}]})_{\MTLayer=1}^{\MTDepth-1}$ in the opening trapdoor $\MTTrapdoor$ can be derived via at most $\MTMessageLength$ random oracle calls from the commitments $(\MTVertexLabelByName{\MTDepth}{\MTMessageIndex})_{\MTMessageIndex \in [\MTMessageLength]}$, and thus could be omitted from $\MTTrapdoor$. In such a case, $\MTOpen$ would have to query the random oracle to (re)derive the commitments required to construct the relevant authentication paths. In fact, $(\MTVertexLabelByName{\MTDepth}{\MTMessageIndex})_{\MTMessageIndex \in [\MTMessageLength]}$ can themselves be derived from the message $\MTMessageVector$ and salts $\MTSaltStrings$ via $\MTMessageLength$ random oracle calls; so one could even set $\MTTrapdoor \DefineEqual (\MTMessageVector,\MTSaltStrings)$, which may be advantageous if $\log\Cardinality{\MTAlphabet}$ is smaller than $\ROOutputSize$. These can be useful time-memory tradeoffs in practice.
\end{remark}

\begin{remark}[local updates]
\label{remark:locally-updatable}
Merkle commitments are \emph{locally updatable}, which is a useful feature in many applications (though we do not explore it further within the context of succinct arguments). Suppose you are given a Merkle commitment $\MTCommitment$ of a message $\MTMessageVector$, and you want to update it to be a Merkle commitment $\MTCommitment'$ of a message $\MTMessageVector'$ where $\MTMessageVector'$ is equal to $\MTMessageVector$ except for the $\MTMessageIndex$-th entry. There is no need to recompute the tree from scratch: one only needs to recompute the commitments for the vertices on the path $\MTPath(\MTMessageIndex)$, while the other commitments remain unchanged. The work needed for this update is proportional to the depth of the tree (which is logarithmic in the message length) rather than the message length.
\end{remark}

\begin{remark}[$\MTMessageLength=1$]
\label{remark:mt-message-length-one}
The message length $\MTMessageLength=1$ is a power of $2$, and in this case the binary tree $\MTGraph_{\MTMessageLength}$ contains a single vertex, the root vertex. In this degenerate case ($\MTMessageLength=1$) the Merkle commitment $\MTSymbol$ happens to equal the basic commitment $\CMSymbol$.
\end{remark}

\begin{remark}[any message length]
\label{remark:mt-any-message-length}
The description of the Merkle commitment scheme $\MTSymbol$ given above involves a binary tree over $\MTMessageLength$ leaves, where $\MTMessageLength$ is the number of entries in the message vector. Notationally it is far simpler to describe and analyze the case when $\MTMessageLength$ is a power of $2$ (which we assume throughout this chapter) because the internal vertices in the binary tree can be labeled via a simple naming scheme and, more generally, all paths from the root vertex to a leaf vertex \DoQuote{look the same}. But how do Merkle commitment schemes work when $\MTMessageLength$ is not a power of $2$?

One option is to pad the message with dummy values (an arbitrary symbol from the alphabet $\MTAlphabet$) until the padded message length reaches a power of $2$; this, at most, doubles the length relative to the original (unpadded) message. All security analyses of this chapter hold for the padded message, and in particular for the original (unpadded) message.

While the padding is not expensive, there are cheaper solutions that involve binary trees with precisely $\MTMessageLength$ leaves (without using any padding). This optimization is explained in \Cref{section:any-message-length}.
\end{remark}

\begin{remark}[other arities]
\label{remark:other-arities}
The description above involves a perfect \emph{binary} tree. However the construction of a Merkle commitment scheme and the properties that we study (completeness, binding, extractability, hiding) directly extend to perfect trees of any arity. (In fact, see \Cref{section:any-message-length} for a \emph{generalized} Merkle commitment scheme that is defined for \emph{any tree}.)  If the tree has arity $\MTArity \geq 2$ then each authentication path contains $\MTArity-1$ vertices per layer rather than $1$, and the number of layers in the tree is $\log_{\MTArity} \MTMessageLength$ rather than $\log_{2} \MTMessageLength$. Larger arities are usually not advantageous for the size of opening proofs (the tree has fewer layers but there are more sibling vertices for an overall bigger opening proof size), though there may be other reasons to consider trees of higher arity (e.g., fewer calls to the random oracle).
\end{remark}

%%%%%%%%%%%%%%%%%%%%%%%%%%%%%%%%%%%%%%%%%%%%%%%%%%%%%%%%%%%%%%%%%%%%%%%%%%%%%%%
%%%%%%%%%%%%%%%%%%%%%%%%%%%%%%%%%%%%%%%%%%%%%%%%%%%%%%%%%%%%%%%%%%%%%%%%%%%%%%%
%%%%%%%%%%%%%%%%%%%%%%%%%%%%%%%%%%%%%%%%%%%%%%%%%%%%%%%%%%%%%%%%%%%%%%%%%%%%%%%
\section{Completeness}
\label{section:merkle-commitment-completeness}

We prove that $\MTCommit$ has \emph{completeness}, namely, if the algorithms of the Merkle commitment scheme are used as intended then validity checks always pass. In more detail, $\MTCommit$ can be used to commit to any message vector, and $\MTOpen$ can be used to open any subset of locations of the message vector in such a way that the resulting opening proof is accepted by $\MTCheck$.

\begin{lemma}[$\MTSymbol$ is complete]
\label{lemma:mt-completeness}
Let $\MTSymbol \DefineEqual \MTConstructor{\ROOutputSize}{\MTAlphabet}{\MTMessageLength}{\MTSaltSize}$. For every adversary $\MTAdversary$,
\begin{equation*}
\Pr\left[
\MTCheck^{\ROFunction}(\MTCommitment,\MTMessageIndexSet,\MTMessageVector[\MTMessageIndexSet],\MTProof)=1
\GivenExperiment
\StateExperiment{
\ROFunction \gets \RODistribution{\ROOutputSize} \\
(\MTMessageVector,\MTMessageIndexSet) \gets \MTAdversary^{\ROFunction} \\
(\MTCommitment,\MTTrapdoor) \gets \MTCommit^{\ROFunction}(\MTMessageVector) \\
\MTProof \gets \MTOpen^{\ROFunction}(\MTTrapdoor, \MTMessageIndexSet)
}
\right]
= 1
\enspace.
\end{equation*}
\end{lemma}

A formal proof of this lemma would involve analyzing each step of each Merkle commitment algorithm in order to establish that all validity checks pass. This would not provide useful insights. Instead we sketch the proof at high level to provide the main intuition.

Fix an arbitrary choice of oracle $\ROFunction \in \RODistribution{\ROOutputSize}$, message $\MTMessageVector \in \MTAlphabet^{\MTMessageLength}$, and subset $\MTMessageIndexSet \subseteq [\MTMessageLength]$.

The algorithm $\MTCommit$ computes the Merkle commitment for $\MTMessageVector$ by first committing to each entry of $\MTMessageVector$ (using a random salt $\MTSaltString_{\MTMessageIndex}$ for each message entry $\MTMessageVector[\MTMessageIndex]$), then pairwise hashes the resulting list to halve its size, and so on, until it obtains a single output $\MTCommitment$; the trapdoor $\MTTrapdoor$ is set to contain all the salts that were sampled and all the commitments $\MTVertexLabels \DefineEqual ((\MTVertexLabelByName{\MTLayer}{\MTMessageIndex})_{\MTMessageIndex \in [2^{\MTLayer}]})_{\MTLayer=0}^{\MTDepth}$ that were computed.

Subsequently, the algorithm $\MTOpen$ includes in the opening proof $\MTProof$ an authentication path $\MTAuthPath_{\MTMessageIndex} = \big(\MTSaltString_{\MTMessageIndex},(\MTVertexLabelByCoPath{\MTMessageIndex}{\MTLayer})_{\MTLayer \in \{1,\dots,\MTDepth\}}\big)$ for each index $\MTMessageIndex$ in the subset $\MTMessageIndexSet$.

The algorithm $\MTCheck$, for each index $\MTMessageIndex$ in the subset $\MTMessageIndexSet$, checks that the authentication path $\MTAuthPath_{\MTMessageIndex}$ is valid relative to the Merkle commitment $\MTCommitment$, index $\MTMessageIndex$, and value $\MTMessageSubVector[\MTMessageIndex]$. The check for each authentication path succeeds because $\MTOpen$ constructed that path by including the correct salt and the correct commitments, which will lead $\MTCheck$ to recompute the same value for the root vertex as $\MTCommitment$.



%%%%%%%%%%%%%%%%%%%%%%%%%%%%%%%%%%%%%%%%%%%%%%%%%%%%%%%%%%%%%%%%%%%%%%%%%%%%%%%
%%%%%%%%%%%%%%%%%%%%%%%%%%%%%%%%%%%%%%%%%%%%%%%%%%%%%%%%%%%%%%%%%%%%%%%%%%%%%%%
%%%%%%%%%%%%%%%%%%%%%%%%%%%%%%%%%%%%%%%%%%%%%%%%%%%%%%%%%%%%%%%%%%%%%%%%%%%%%%%
\section{Collision lemma}
\label{section:merkle-commitment-collision-lemma}

We formulate and prove a \emph{collision lemma} for Merkle commitments, which is a basic security property that we use in later sections. Informally, the lemma states that authenticating different values for the same leaf in a Merkle commitment leads to a collision in the random oracle. Stating this property precisely, however, requires some care. The discussion below takes an incremental approach by considering simpler types of collision lemmas. First, we describe a collision lemma for the basic commitment scheme $\CMSymbol$ in \Cref{chapter:basic-commitment}. Second, we describe a special case of the collision lemma for Merkle commitment schemes. Third, and finally, we describe the (general) collision lemma for Merkle commitment schemes.

\parhead{Starting simple: collisions in $\CMSymbol$}
Consider the basic commitment scheme $\CMSymbol$ in \Cref{chapter:basic-commitment}, whose checking algorithm $\CMCheck$ simply consists of checking a query-answer pair of the random oracle: $\CMCheck^{\ROFunction}(\CMCommitment,\CMMessage,\CMSaltString)=1$ if and only if $\CMCommitment = \ROFunction(\CMMessage,\CMSaltString)$. Consider a commitment $\CMCommitment$ and two message-salt pairs, $(\CMMessage_{0},\CMSaltString_{0})$ and $(\CMMessage_{1},\CMSaltString_{1})$. Suppose that both pairs are valid openings of the commitment, that is, $\CMCheck^{\ROFunction}(\CMCommitment,\CMMessage_{0},\CMSaltString_{0})=1$ and $\CMCheck^{\ROFunction}(\CMCommitment,\CMMessage_{1},\CMSaltString_{1})=1$. By the definition of $\CMCheck$, these conditions are equivalent to $\CMCommitment = \ROFunction(\CMMessage_{0},\CMSaltString_{0})$ and $\CMCommitment = \ROFunction(\CMMessage_{1},\CMSaltString_{1})$, from which we deduce that $\ROFunction(\CMMessage_{0},\CMSaltString_{0})=\ROFunction(\CMMessage_{1},\CMSaltString_{1})$. This tells us that if either $\CMMessage_{0} \neq \CMMessage_{1}$ (the messages are distinct) or $\CMSaltString_{0} \neq \CMSaltString_{1}$ (the salts are distinct) then the two executions of $\CMCheck$ on the two inputs result in a collision for $\ROFunction$. Thus, in $\CMSymbol$, \emph{opening the same commitment either via two different messages or via two different salts leads to a collision}. We can summarize this property as a \DoQuote{collision lemma for $\CMSymbol$} via the following implication:
\begin{equation*}
\left[
\begin{array}{l}
\big(
\CMMessage_{0} \neq \CMMessage_{1}
\lor
\CMSaltString_{0} \neq \CMSaltString_{1}
\big) \\
\land \;\CMCheck^{\ROFunction}(\CMCommitment,\CMMessage_{0},\CMSaltString_{0})=1 \\
\land \;\CMCheck^{\ROFunction}(\CMCommitment,\CMMessage_{1},\CMSaltString_{1})=1
\end{array}
\right]
\Rightarrow\;
\left[
\begin{array}{c}
\text{$(\CMMessage_{0},\CMSaltString_{0})$ and $(\CMMessage_{1},\CMSaltString_{1})$} \\
\text{are a collision for $\ROFunction$}
\end{array}
\right]
\enspace.
\end{equation*}

\parhead{What about $\MTSymbol$?}
We wish to derive a similar lemma for the Merkle commitment scheme $\MTSymbol$. The challenge is that messages in $\MTSymbol$ can be opened at subsets of locations (rather than only at all locations), which complicates the assumptions under which a collision can be proved to exist. Moreover, the algorithm $\MTCheck$, which validates the corresponding opening proofs, involves a more complex computation. In particular, the computation involves multiple queries to the random oracle (rather than one query as in $\CMSymbol$), and arguing that a collision exists in the relevant query-answer traces will involve a delicate case analysis. We tackle this additional complexity by first analyzing a special case and then analyzing the general case.

\parhead{Special case: opening a single leaf}
As a warmup we consider the special case where we only consider opening proofs for a single leaf. In other words, we consider computations for $\MTCheck^{\ROFunction}(\MTCommitment,\MTMessageIndexSet,\MTMessageSubVector,\MTProof)$ where the set $\MTMessageIndexSet = \{\MTMessageIndex\}$ is a singleton and so:
\begin{inparaenum}[(i)]
  \item $\MTMessageSubVector$ is a vector consisting of a single entry $\MTMessageEntry$ in $\MTAlphabet$; and
  \item $\MTProof$ consists of a single authentication path $\MTAuthPath$.
\end{inparaenum}

Consider a Merkle commitment $\MTCommitment$ and position $\MTMessageIndex \in [\MTMessageLength]$, and two openings: a claimed value $\MTMessageEntry_{0} \in \MTAlphabet$ and corresponding authentication path $\MTAuthPath_{0}$, and another claimed value $\MTMessageEntry_{1} \in \MTAlphabet$ and corresponding authentication path $\MTAuthPath_{1}$. Suppose that both openings are valid, that is, $\MTCheck^{\ROFunction}(\MTCommitment,\MTMessageIndex,\MTMessageEntry_{0},\MTAuthPath_{0})=1$ and $\MTCheck^{\ROFunction}(\MTCommitment,\MTMessageIndex,\MTMessageEntry_{1},\MTAuthPath_{1})=1$. Intuitively, since the two computations share the same Merkle commitment $\MTCommitment$, if $\MTMessageEntry_{0} \neq \MTMessageEntry_{1}$ then the two random oracle computations \DoQuote{collide} at some layer of the binary tree in order to eventually lead to the same Merkle commitment $\MTCommitment$ at the root of the tree. The collision is either at a leaf commitment, or at an internal commitment; see \Cref{figure:mt-collision-lemma} for an example. In fact, we can prove that the same is true even if $\MTMessageEntry_{0} = \MTMessageEntry_{1}$ and $\MTAuthPath_{0} \neq \MTAuthPath_{1}$.

\begin{figure}[htp!]
\centering
\begin{subfigure}[b]{0.4\textwidth}
\centering
  \includegraphics[width=\textwidth]{\FigureFolder/mt-collision-lemma-diagram-1}
  \caption{}
  \label{figure:mt-collision-lemma-1}
\end{subfigure}
\quad
\begin{subfigure}[b]{0.55\textwidth}
  \centering
  \includegraphics[width=\textwidth]{\FigureFolder/mt-collision-lemma-diagram-2}
  \caption{}
  \label{figure:mt-collision-lemma-2}
\end{subfigure}
\caption[Collision between two authentication paths.]{Example of a collision for two paths that authenticate two different values for the $6$-th entry (out of $8$ entries). There are two cases: in the left, a collision at a leaf commitment; and, on the right, a collision at an internal commitment.}
\label{figure:mt-collision-lemma}
\end{figure}

More precisely, the collision involves a query from the query-answer trace $\ROHonestTrace_{0}$ of the computation $\MTCheck^{\ROFunction}(\MTCommitment,\MTMessageIndex,\MTMessageEntry_{0},\MTAuthPath_{0})$ and a query from the query-answer trace $\ROHonestTrace_{1}$ of the computation $\MTCheck^{\ROFunction}(\MTCommitment,\MTMessageIndex,\MTMessageEntry_{1},\MTAuthPath_{1})$.

Note that here it is necessary for the two openings to be about the same position $\MTMessageIndex$. This is because it is possible to open two different positions via two different values (or different authentication paths) without causing collisions. Indeed, this would happen for most choices of random oracle for a message with different values.

We formally state and prove this below.

\begin{lemma}
\label{lemma:simple-mt-colliding-paths}
Let $\MTSymbol \DefineEqual \MTConstructor{\ROOutputSize}{\MTAlphabet}{\MTMessageLength}{\MTSaltSize}$. Fix arbitrary choices of oracle $\ROFunction \in \RODistribution{\ROOutputSize}$ and:
\begin{inparaenum}[\emptylabel]
  \item a Merkle commitment $\MTCommitment \in \Bits^{\ROOutputSize}$;
  \item an index $\MTMessageIndex \in [\MTMessageLength]$;
  \item message entries $\MTMessageEntry_{0},\MTMessageEntry_{1} \in \MTAlphabet$; and
  \item authentication paths $\MTAuthPath_{0},\MTAuthPath_{1}$.
\end{inparaenum}
For $b \in \Bits$, let $\ROHonestTrace_{b}$ be the query-answer trace of $\MTCheck^{\ROFunction}(\MTCommitment,\MTMessageIndex,\MTMessageEntry_{b},\MTAuthPath_{b})$. Then the following implication holds:
\begin{equation}
\label{equation:simple-mt-colliding-paths}
\left[
\begin{array}{l}
\big(
\MTMessageEntry_{0} \neq \MTMessageEntry_{1}
\lor
\MTAuthPath_{0} \neq \MTAuthPath_{1}
\big) \\
\land \;\MTCheck^{\ROFunction}(\MTCommitment,\MTMessageIndex,\MTMessageEntry_{0},\MTAuthPath_{0})=1 \\
\land \;\MTCheck^{\ROFunction}(\MTCommitment,\MTMessageIndex,\MTMessageEntry_{1},\MTAuthPath_{1})=1
\end{array}
\right]
\Rightarrow\;
\left[
\begin{array}{c}
\exists\,
(\ROQuery_{0},\ROAnswer_{0}) \in \ROHonestTrace_{0}
,
(\ROQuery_{1},\ROAnswer_{1}) \in \ROHonestTrace_{1}
\\
\text{with $\ROQuery_{0} \neq \ROQuery_{1}$}
\text{ such that } \ROAnswer_{0}=\ROAnswer_{1}
\end{array}
\right]
\enspace.
\end{equation}
\end{lemma}

\begin{proof}
Fix $b \in \Bits$. The authentication path $\MTAuthPath_{b}$ for index $\MTMessageIndex$ has the following structure:
\begin{equation*}
\MTAuthPath_{b}
=
\Big(
  \MTSaltString_{\MTMessageIndex}^{(b)}
  ,
  (\MTVertexLabelByCoPath{\MTMessageIndex}{\MTLayer}^{(b)})_{\MTLayer \in \{1,\dots,\MTDepth\}}
\Big)
\enspace,
\end{equation*}
where $\MTSaltString_{\MTMessageIndex}^{(b)}$ is a Merkle salt and $(\MTVertexLabelByCoPath{\MTMessageIndex}{\MTLayer}^{(b)})_{\MTLayer \in \{1,\dots,\MTDepth\}}$ are commitments labeling vertices in $\MTCoPath(\MTMessageIndex)$.

Moreover, the computation $\MTCheck^{\ROFunction}(\MTCommitment,\MTMessageIndex,\MTMessageEntry_{b},\MTAuthPath_{b})$ leads to a query-answer trace $\ROHonestTrace_{b}$ consisting of $\MTDepth+1$ query-answer pairs. The answers in $\ROHonestTrace_{b}$ are denoted
\begin{equation*}
(\MTVertexLabelByPath{\MTMessageIndex}{\MTLayer}^{(b)})_{\MTLayer \in \{0,1,\dots,\MTDepth\}}
\enspace.
\end{equation*}
Finally, $\MTCheck^{\ROFunction}(\MTCommitment,\MTMessageIndex,\MTMessageEntry_{b},\MTAuthPath_{b}) = 1$ implies that $\MTCommitment = \MTVertexLabelByName{0}{1}^{(b)}$.

In particular, we deduce that $\MTVertexLabelByName{0}{1}^{(0)}=\MTVertexLabelByName{0}{1}^{(1)}$.

We consider two cases.
\begin{itemize}

  \item \emph{Case 1: there exists $\MTLayer \in [\MTDepth]$ such that $\MTVertexLabelByPath{\MTMessageIndex}{\MTLayer}^{(0)} \neq \MTVertexLabelByPath{\MTMessageIndex}{\MTLayer}^{(1)}$ or $\MTVertexLabelByCoPath{\MTMessageIndex}{\MTLayer}^{(0)} \neq \MTVertexLabelByCoPath{\MTMessageIndex}{\MTLayer}^{(1)}$.}

Let $\MTLayer$ be the minimal index satisfying this condition.
\begin{itemize}[nolistsep]
  \item If $\MTPathIndex{\MTMessageIndex}{\MTLayer}$ is an odd vertex then for $b \in \Bits$ set $(\MTLeftVertexLabel^{(b)},\MTRightVertexLabel^{(b)}) \DefineEqual (\MTVertexLabelByPath{\MTMessageIndex}{\MTLayer}^{(b)},\MTVertexLabelByCoPath{\MTMessageIndex}{\MTLayer}^{(b)})$.
  \item If $\MTPathIndex{\MTMessageIndex}{\MTLayer}$ is an even vertex then for $b \in \Bits$ set $(\MTLeftVertexLabel^{(b)},\MTRightVertexLabel^{(b)}) \DefineEqual (\MTVertexLabelByCoPath{\MTMessageIndex}{\MTLayer}^{(b)},\MTVertexLabelByPath{\MTMessageIndex}{\MTLayer}^{(b)})$.
\end{itemize}
For $b \in \Bits$ the pair $(\MTLeftVertexLabel^{(b)},\MTRightVertexLabel^{(b)})$ is a query in $\ROHonestTrace_{b}$ (see \Cref{step:mtcheck:internal-query} in the description of $\MTCheck$).

These two queries are a collision for $\ROFunction$ due to the following.
\begin{itemize}[nosep]

  \item $(\MTLeftVertexLabel^{(0)},\MTRightVertexLabel^{(0)}) \neq (\MTLeftVertexLabel^{(1)},\MTRightVertexLabel^{(1)})$. \\ This follows since $\MTVertexLabelByPath{\MTMessageIndex}{\MTLayer}^{(0)} \neq \MTVertexLabelByPath{\MTMessageIndex}{\MTLayer}^{(1)}$ or $\MTVertexLabelByCoPath{\MTMessageIndex}{\MTLayer}^{(0)} \neq \MTVertexLabelByCoPath{\MTMessageIndex}{\MTLayer}^{(1)}$.

  \item $\ROFunction(\MTLeftVertexLabel^{(0)},\MTRightVertexLabel^{(0)})=\ROFunction(\MTLeftVertexLabel^{(1)},\MTRightVertexLabel^{(1)})$. \\ This follows since $\MTVertexLabelByPath{\MTMessageIndex}{\MTLayer-1}^{(0)} = \ROFunction(\MTLeftVertexLabel^{(0)},\MTRightVertexLabel^{(0)})$, $\MTVertexLabelByPath{\MTMessageIndex}{\MTLayer-1}^{(1)} = \ROFunction(\MTLeftVertexLabel^{(1)},\MTRightVertexLabel^{(1)})$, and $\MTVertexLabelByPath{\MTMessageIndex}{\MTLayer-1}^{(0)} = \MTVertexLabelByPath{\MTMessageIndex}{\MTLayer-1}^{(1)}$. The last equality holds because $\MTLayer$ is the minimal index satisfying the condition (and we know that $\MTVertexLabelByName{0}{1}^{(0)}=\MTVertexLabelByName{0}{1}^{(1)}$).

\end{itemize}

  \item \emph{Case 2: $\MTVertexLabelByName{\MTDepth}{\MTMessageIndex}^{(0)}=\MTVertexLabelByName{\MTDepth}{\MTMessageIndex}^{(1)}$, and $\MTMessageEntry_{0} \neq \MTMessageEntry_{1}$ or $\MTSaltString_{\MTMessageIndex}^{(0)} \neq \MTSaltString_{\MTMessageIndex}^{(1)}$.}

For $b \in \Bits$ the pair $(\MTMessageEntry_{b},\MTSaltString_{\MTMessageIndex}^{(b)})$ is a query in $\ROHonestTrace_{b}$ (see \Cref{step:mtcheck:leaf-query} in the description of $\MTCheck$).

These two queries are a collision for $\ROFunction$ due to the following.
\begin{itemize}[nosep]
  \item $(\MTMessageEntry_{0},\MTSaltString_{\MTMessageIndex}^{(0)}) \neq (\MTMessageEntry_{1},\MTSaltString_{\MTMessageIndex}^{(1)})$. This follows since $\MTMessageEntry_{0} \neq \MTMessageEntry_{1}$ or $\MTSaltString_{\MTMessageIndex}^{(0)} \neq \MTSaltString_{\MTMessageIndex}^{(1)}$.
  \item $\ROFunction(\MTMessageEntry_{0},\MTSaltString_{\MTMessageIndex}^{(0)})=\ROFunction(\MTMessageEntry_{1},\MTSaltString_{\MTMessageIndex}^{(1)})$. This follows since $\MTVertexLabelByName{\MTDepth}{\MTMessageIndex}^{(0)} = \ROFunction(\MTMessageEntry_{0},\MTSaltString_{\MTMessageIndex}^{(0)})$, $\MTVertexLabelByName{\MTDepth}{\MTMessageIndex}^{(1)} = \ROFunction(\MTMessageEntry_{1},\MTSaltString_{\MTMessageIndex}^{(1)})$, and $\MTVertexLabelByName{\MTDepth}{\MTMessageIndex}^{(0)}=\MTVertexLabelByName{\MTDepth}{\MTMessageIndex}^{(1)}$.
\end{itemize}

\end{itemize}
In either case we find a collision with one query in $\ROHonestTrace_{0}$ and the other query in $\ROHonestTrace_{1}$.

We are left to argue that the antecedent in \Cref{equation:simple-mt-colliding-paths} implies that Case 1 or Case 2 holds.

If Case 1 holds then we are done. So suppose that Case 1 does not hold, which means that $(\MTVertexLabelByPath{\MTMessageIndex}{\MTLayer}^{(0)})_{\MTLayer \in [\MTDepth]}=(\MTVertexLabelByPath{\MTMessageIndex}{\MTLayer}^{(1)})_{\MTLayer \in [\MTDepth]}$ and $(\MTVertexLabelByCoPath{\MTMessageIndex}{\MTLayer}^{(0)})_{\MTLayer \in [\MTDepth]}=(\MTVertexLabelByCoPath{\MTMessageIndex}{\MTLayer}^{(1)})_{\MTLayer \in [\MTDepth]}$ (and in particular that $\MTVertexLabelByName{\MTDepth}{\MTMessageIndex}^{(0)}=\MTVertexLabelByName{\MTDepth}{\MTMessageIndex}^{(1)}$). If $\MTMessageEntry_{0} \neq \MTMessageEntry_{1}$ then Case 2 holds. Suppose instead that $\MTAuthPath_{0} \neq \MTAuthPath_{1}$. The only way for the two authentication paths to differ in this case is for the Merkle salts to differ ($\MTSaltString_{\MTMessageIndex}^{(0)} \neq \MTSaltString_{\MTMessageIndex}^{(1)}$), in which case again Case 2 holds.
\end{proof}

\parhead{General case: opening multiple leaves}
We now consider the collision lemma for the general case, where opening proofs may authenticate multiple leaves. Consider a Merkle commitment $\MTCommitment$, as well as:
\begin{inparaenum}[\emptylabel]
  \item two index sets $\MTMessageIndexSet_{0},\MTMessageIndexSet_{1}$;
  \item two vectors $\MTMessageSubVector_{0},\MTMessageSubVector_{1}$;
  \item two opening proofs $\MTProof_{0},\MTProof_{1}$.
\end{inparaenum}
Suppose that both openings are valid, that is, $\MTCheck^{\ROFunction}(\MTCommitment,\MTMessageIndexSet_{0},\MTMessageSubVector_{0},\MTProof_{0})=1$ and $\MTCheck^{\ROFunction}(\MTCommitment,\MTMessageIndexSet_{1},\MTMessageSubVector_{1},\MTProof_{1})=1$.

Intuitively, and similarly to the case of a single leaf, if there is a position $\MTMessageIndex \in \MTMessageIndexSet_{0} \cap \MTMessageIndexSet_{1}$ such that $\MTMessageSubVector_{0}[\MTMessageIndex] \neq \MTMessageSubVector_{1}[\MTMessageIndex]$ then we can find a collision in (the union of the relevant) query-answer traces. Informally, the prior argument works applied to the authentication paths for the index $\MTMessageIndex$ in the opening proofs $\MTProof_{0}$ and $\MTProof_{1}$ (and the subset of query-answer pairs used to validate these paths).

Less obvious is the fact that we cannot expect (in general) to find a collision if $\MTProof_{0} \neq \MTProof_{1}$. This is because the two opening proofs are for potentially different index sets $\MTMessageIndexSet_{0}$ and $\MTMessageIndexSet_{1}$, so the opening proofs can be different. Nevertheless, we can prove that if the opening proofs are for the same index set ($\MTMessageIndexSet_{0}=\MTMessageIndexSet_{1}$) and they are different ($\MTProof_{0} \neq \MTProof_{1}$) then we can find a collision.

We formally state and prove this below.

\begin{lemma}
\label{lemma:mt-colliding-paths}
Let $\MTSymbol \DefineEqual \MTConstructor{\ROOutputSize}{\MTAlphabet}{\MTMessageLength}{\MTSaltSize}$. Fix arbitrary choices of oracle $\ROFunction \in \RODistribution{\ROOutputSize}$ and:
\begin{itemize}[nolistsep]
  \item a Merkle commitment $\MTCommitment \in \Bits^{\ROOutputSize}$;
  \item a set $\MTMessageIndexSet_{0} \subseteq [\MTMessageLength]$, vector $\MTMessageSubVector_{0} \in \MTAlphabet^{\MTMessageIndexSet_{0}}$, and opening proof $\MTProof_{0}$;
  \item a set $\MTMessageIndexSet_{1} \subseteq [\MTMessageLength]$, vector $\MTMessageSubVector_{1} \in \MTAlphabet^{\MTMessageIndexSet_{1}}$, and opening proof $\MTProof_{1}$.
\end{itemize}
For $b \in \Bits$, let $\ROHonestTrace_{b}$ be the query-answer trace of $\MTCheck^{\ROFunction}(\MTCommitment,\MTMessageIndexSet_{b},\MTMessageSubVector_{b},\MTProof_{b})$. Then the following implication holds:
\begin{equation}
\label{equation:mt-colliding-paths}
\left[
\begin{array}{l}
\begin{pmatrix}
\exists\;\MTMessageIndex \in \MTMessageIndexSet_{0} \cap \MTMessageIndexSet_{1} : \MTMessageSubVector_{0}[\MTMessageIndex] \neq \MTMessageSubVector_{1}[\MTMessageIndex] \\
\text{or} \\
\MTMessageIndexSet_{0} = \MTMessageIndexSet_{1} \; \land \; \MTProof_{0} \neq \MTProof_{1}
\end{pmatrix} \\
\land \;\MTCheck^{\ROFunction}(\MTCommitment,\MTMessageIndexSet_{0},\MTMessageSubVector_{0},\MTProof_{0})=1 \\
\land \;\MTCheck^{\ROFunction}(\MTCommitment,\MTMessageIndexSet_{1},\MTMessageSubVector_{1},\MTProof_{1})=1
\end{array}
\right]
\Rightarrow\;
\left[
\begin{array}{c}
\exists\,
(\ROQuery_{0},\ROAnswer_{0}) \in \ROHonestTrace_{0}
,
(\ROQuery_{1},\ROAnswer_{1}) \in \ROHonestTrace_{1}
\\
\text{with $\ROQuery_{0} \neq \ROQuery_{1}$}
\text{ such that } \ROAnswer_{0}=\ROAnswer_{1}
\end{array}
\right]
\enspace.
\end{equation}
\end{lemma}

\begin{proof}
Fix $b \in \Bits$. The opening proof $\MTProof_{b}$ for the set $\MTMessageIndexSet_{b}$ has the following structure:
\begin{equation*}
\MTProof_{b}
=
\Big(
  \MTSaltString_{\MTMessageIndex}^{(b)}
  ,
  (\MTVertexLabelByCoPath{\MTMessageIndex}{\MTLayer}^{(b)})_{\MTLayer \in \{1,\dots,\MTDepth\}}
\Big)_{\MTMessageIndex \in \MTMessageIndexSet_{b}}
\enspace.
\end{equation*}
Above, the value $\MTSaltString_{\MTMessageIndex}^{(b)}$ is the Merkle salt used to authenticate leaf $\MTMessageIndex$, and the values $(\MTVertexLabelByCoPath{\MTMessageIndex}{\MTLayer}^{(b)})_{\MTLayer \in \{1,\dots,\MTDepth\}}$ are the commitments labeling vertices in $\MTCoPath(\MTMessageIndex)$ used to authenticate leaf $\MTMessageIndex$.

Moreover, the computation $\MTCheck^{\ROFunction}(\MTCommitment,\MTMessageIndexSet_{b},\MTMessageSubVector_{b},\MTProof_{b})$ obtains as answers from $\ROFunction$ the commitments corresponding to vertices in $\MTPath(\MTMessageIndex)$. We denote these commitments as
\begin{equation*}
(\MTVertexLabelByPath{\MTMessageIndex}{\MTLayer}^{(b)})_{\MTLayer \in \{0,1,\dots,\MTDepth\}}
\enspace.
\end{equation*}

Finally, $\MTCheck^{\ROFunction}(\MTCommitment,\MTMessageIndexSet_{b},\MTMessageSubVector_{b},\MTProof_{b}) = 1$ implies that $\MTCommitment = \MTVertexLabelByName{0}{1}^{(b)}$.

In particular, we deduce that $\MTVertexLabelByName{0}{1}^{(0)}=\MTVertexLabelByName{0}{1}^{(1)}$.

We consider two cases.
\begin{enumerate}

  \item \emph{Case 1: there exist $\MTMessageIndex \in \MTMessageIndexSet_{0} \cap \MTMessageIndexSet_{1}$ and $\MTLayer \in [\MTDepth]$ such that $\MTVertexLabelByPath{\MTMessageIndex}{\MTLayer}^{(0)} \neq \MTVertexLabelByPath{\MTMessageIndex}{\MTLayer}^{(1)}$ or $\MTVertexLabelByCoPath{\MTMessageIndex}{\MTLayer}^{(0)} \neq \MTVertexLabelByCoPath{\MTMessageIndex}{\MTLayer}^{(1)}$.}

Let $\MTLayer$ be the minimal index satisfying this condition (for an arbitrary choice of $\MTMessageIndex$).
\begin{itemize}[nolistsep]
  \item If $\MTPathIndex{\MTMessageIndex}{\MTLayer}$ is an odd vertex then for $b \in \Bits$ set $(\MTLeftVertexLabel^{(b)},\MTRightVertexLabel^{(b)}) \DefineEqual (\MTVertexLabelByPath{\MTMessageIndex}{\MTLayer}^{(b)},\MTVertexLabelByCoPath{\MTMessageIndex}{\MTLayer}^{(b)})$.
  \item If $\MTPathIndex{\MTMessageIndex}{\MTLayer}$ is an even vertex then for $b \in \Bits$ set $(\MTLeftVertexLabel^{(b)},\MTRightVertexLabel^{(b)}) \DefineEqual (\MTVertexLabelByCoPath{\MTMessageIndex}{\MTLayer}^{(b)},\MTVertexLabelByPath{\MTMessageIndex}{\MTLayer}^{(b)})$.
\end{itemize}
For $b \in \Bits$ the pair $(\MTLeftVertexLabel^{(b)},\MTRightVertexLabel^{(b)})$ is the query in $\ROHonestTrace_{b}$ (see \Cref{step:mtcheck:internal-query} in the description of $\MTCheck$).

These two queries are a collision for $\ROFunction$ due to the following.
\begin{itemize}[nosep]

  \item $(\MTLeftVertexLabel^{(0)},\MTRightVertexLabel^{(0)}) \neq (\MTLeftVertexLabel^{(1)},\MTRightVertexLabel^{(1)})$. \\ This follows since $\MTVertexLabelByPath{\MTMessageIndex}{\MTLayer}^{(0)} \neq \MTVertexLabelByPath{\MTMessageIndex}{\MTLayer}^{(1)}$ or $\MTVertexLabelByCoPath{\MTMessageIndex}{\MTLayer}^{(0)} \neq \MTVertexLabelByCoPath{\MTMessageIndex}{\MTLayer}^{(1)}$.

  \item $\ROFunction(\MTLeftVertexLabel^{(0)},\MTRightVertexLabel^{(0)})=\ROFunction(\MTLeftVertexLabel^{(1)},\MTRightVertexLabel^{(1)})$. \\ This follows since $\MTVertexLabelByPath{\MTMessageIndex}{\MTLayer-1}^{(0)} = \ROFunction(\MTLeftVertexLabel^{(0)},\MTRightVertexLabel^{(0)})$, $\MTVertexLabelByPath{\MTMessageIndex}{\MTLayer-1}^{(1)} = \ROFunction(\MTLeftVertexLabel^{(1)},\MTRightVertexLabel^{(1)})$, and $\MTVertexLabelByPath{\MTMessageIndex}{\MTLayer-1}^{(0)} = \MTVertexLabelByPath{\MTMessageIndex}{\MTLayer-1}^{(1)}$. The last equality holds because $\MTLayer$ is a minimal index satisfying the condition for this choice of $\MTMessageIndex$ (and we know that $\MTVertexLabelByName{0}{1}^{(0)}=\MTVertexLabelByName{0}{1}^{(1)}$).

\end{itemize}

  \item \emph{Case 2: there exists $\MTMessageIndex \in \MTMessageIndexSet_{0} \cap \MTMessageIndexSet_{1}$ such that $\MTVertexLabelByName{\MTDepth}{\MTMessageIndex}^{(0)}=\MTVertexLabelByName{\MTDepth}{\MTMessageIndex}^{(1)}$ and such that $\MTMessageSubVector_{0}[\MTMessageIndex] \neq \MTMessageSubVector_{1}[\MTMessageIndex]$ or $\MTSaltString_{\MTMessageIndex}^{(0)} \neq \MTSaltString_{\MTMessageIndex}^{(1)}$.}

For $b \in \Bits$ the pair $(\MTMessageSubVector_{b}[\MTMessageIndex],\MTSaltString_{\MTMessageIndex}^{(b)})$ is a query in $\ROHonestTrace_{b}$ (see \Cref{step:mtcheck:leaf-query} in the description of $\MTCheck$).

These two queries are a collision for $\ROFunction$ due to the following.
\begin{itemize}[nosep]
  \item $(\MTMessageSubVector_{0}[\MTMessageIndex],\MTSaltString_{\MTMessageIndex}^{(0)}) \neq (\MTMessageSubVector_{1}[\MTMessageIndex],\MTSaltString_{\MTMessageIndex}^{(1)})$. \\ This follows since $\MTMessageSubVector_{0}[\MTMessageIndex] \neq \MTMessageSubVector_{1}[\MTMessageIndex]$ or $\MTSaltString_{\MTMessageIndex}^{(0)} \neq \MTSaltString_{\MTMessageIndex}^{(1)}$.
  \item $\ROFunction(\MTMessageSubVector_{0}[\MTMessageIndex],\MTSaltString_{\MTMessageIndex}^{(0)})=\ROFunction(\MTMessageSubVector_{1}[\MTMessageIndex],\MTSaltString_{\MTMessageIndex}^{(1)})$. \\ This follows since $\MTVertexLabelByName{\MTDepth}{\MTMessageIndex}^{(0)} = \ROFunction(\MTMessageSubVector_{0}[\MTMessageIndex],\MTSaltString_{\MTMessageIndex}^{(0)})$, $\MTVertexLabelByName{\MTDepth}{\MTMessageIndex}^{(1)} = \ROFunction(\MTMessageSubVector_{1}[\MTMessageIndex],\MTSaltString_{\MTMessageIndex}^{(1)})$, and $\MTVertexLabelByName{\MTDepth}{\MTMessageIndex}^{(0)}=\MTVertexLabelByName{\MTDepth}{\MTMessageIndex}^{(1)}$.
\end{itemize}

\end{enumerate}
In either case we find a collision with one query in $\ROHonestTrace_{0}$ and the other query in $\ROHonestTrace_{1}$.

We are left to argue that the antecedent in \Cref{equation:mt-colliding-paths} implies that Case 1 or Case 2 holds.

Suppose that there exists $\MTMessageIndex \in \MTMessageIndexSet_{0} \cap \MTMessageIndexSet_{1}$ such that $\MTMessageSubVector_{0}[\MTMessageIndex] \neq \MTMessageSubVector_{1}[\MTMessageIndex]$. One of the following must hold.
\begin{itemize}[nolistsep]
  \item If there exists $\MTLayer \in [\MTDepth]$ such that $\MTVertexLabelByPath{\MTMessageIndex}{\MTLayer}^{(0)} \neq \MTVertexLabelByPath{\MTMessageIndex}{\MTLayer}^{(1)}$, then Case 1 holds.
  \item If instead $(\MTVertexLabelByPath{\MTMessageIndex}{\MTLayer}^{(0)})_{\MTLayer \in [\MTDepth]}=(\MTVertexLabelByPath{\MTMessageIndex}{\MTLayer}^{(1)})_{\MTLayer \in [\MTDepth]}$, then in particular $\MTVertexLabelByName{\MTDepth}{\MTMessageIndex}^{(0)}=\MTVertexLabelByName{\MTDepth}{\MTMessageIndex}^{(1)}$ and so Case 2 holds.
\end{itemize}

Alternatively, suppose that $\MTMessageIndexSet_{0} = \MTMessageIndexSet_{1}$ and $\MTProof_{0} \neq \MTProof_{1}$. If Case 1 holds then we are done. So suppose that Case 1 does not hold; in particular we know that $(\MTVertexLabelByCoPath{\MTMessageIndex}{\MTLayer}^{(0)})_{\MTLayer \in [\MTDepth]}=(\MTVertexLabelByCoPath{\MTMessageIndex}{\MTLayer}^{(1)})_{\MTLayer \in [\MTDepth]}$. If so, the only way for $\MTProof_{0} \neq \MTProof_{1}$ to hold is that there exists $\MTMessageIndex \in \MTMessageIndexSet_{0} \cap \MTMessageIndexSet_{1}$ such that $\MTSaltString_{\MTMessageIndex}^{(0)} \neq \MTSaltString_{\MTMessageIndex}^{(1)}$. Moreover for this $\MTMessageIndex$ we know that $\MTVertexLabelByName{\MTDepth}{\MTMessageIndex}^{(0)}=\MTVertexLabelByName{\MTDepth}{\MTMessageIndex}^{(1)}$ because Case 1 does not hold. Hence Case 2 holds.
\end{proof}





%%%%%%%%%%%%%%%%%%%%%%%%%%%%%%%%%%%%%%%%%%%%%%%%%%%%%%%%%%%%%%%%%%%%%%%%%%%%%%%
%%%%%%%%%%%%%%%%%%%%%%%%%%%%%%%%%%%%%%%%%%%%%%%%%%%%%%%%%%%%%%%%%%%%%%%%%%%%%%%
%%%%%%%%%%%%%%%%%%%%%%%%%%%%%%%%%%%%%%%%%%%%%%%%%%%%%%%%%%%%%%%%%%%%%%%%%%%%%%%
\section{Binding}
\label{section:merkle-commitment-binding}

The basic security property of the Merkle commitment scheme is that it is \emph{binding}: a Merkle commitment cannot be opened at the same locations to different values, regardless of the Merkle opening proofs. In more detail, the probability, over the choice of random oracle, that an adversary succeeds in this task is small relative to its query complexity. The salt size $\MTSaltSize$ does not play any role in this property, and in particular it could be zero. This property is analogous to the one that we studied for the basic commitment scheme $\CMSymbol$ in \Cref{chapter:basic-commitment} (see \Cref{lemma:cm-binding} in \Cref{section:basic-commitment-binding}).

Later sections actually rely on a stronger property known as extractability, which we discuss in \Cref{section:merkle-commitment-extractability}. Nevertheless, it is useful to understand binding first in order to familiarize oneself with the basic structure of the Merkle commitment scheme.

\begin{lemma}[$\MTSymbol$ is binding]
\label{lemma:mt-binding}
Let $\MTSymbol \DefineEqual \MTConstructor{\ROOutputSize}{\MTAlphabet}{\MTMessageLength}{\MTSaltSize}$ and define $\MTDepth \DefineEqual \log_{2} \MTMessageLength$. For every query bound $\ROQueryBound \in \Naturals$ and $\ROQueryBound$-query algorithm $\MTAdversary$,
\begin{equation*}
\Pr\left[
\begin{array}{l}
\exists\;\MTMessageIndex \in \MTMessageIndexSet_{0} \cap \MTMessageIndexSet_{1} \text{ s.t. } \MTMessageSubVector_{0}[\MTMessageIndex] \neq \MTMessageSubVector_{1}[\MTMessageIndex] \\
\land \; \MTCheck^{\ROFunction}(\MTCommitment,\MTMessageIndexSet_{0},\MTMessageSubVector_{0},\MTProof_{0})=1 \\
\land \; \MTCheck^{\ROFunction}(\MTCommitment,\MTMessageIndexSet_{1},\MTMessageSubVector_{1},\MTProof_{1})=1
\end{array}
\GivenExperiment
\StateExperiment{
\ROFunction \gets \RODistribution{\ROOutputSize} \\
(\MTCommitment,
\MTMessageIndexSet_{0},\MTMessageSubVector_{0},\MTProof_{0},
\MTMessageIndexSet_{1},\MTMessageSubVector_{1},\MTProof_{1})
\gets \MTAdversary^{\ROFunction}
}
\right]
\leq
\MTBindingError(\ROOutputSize,\MTMessageLength,\ROQueryBound)
\enspace.
\end{equation*}
Above, $\MTBindingError(\ROOutputSize,\MTMessageLength,\ROQueryBound) \leq \MTBindingExpression{\ROOutputSize}{\ROQueryBound}$ if $\ROQueryBound \geq 2(\MTDepth+1)^{2}$.
\end{lemma}

Informally, by the collision lemma in \Cref{section:merkle-commitment-collision-lemma}, breaking the binding property implies a collision in the query-answer traces of the two executions of $\MTCheck$. The proof below upper bounds the probability that the adversary produces outputs that result in such a collision.

Note that a straightforward analysis would consider the probability of a collision in the query-answer trace resulting from concatenating the adversary's trace and the query-answer traces of the two executions of $\MTCheck$; this would yield a bound of $\binom{\ROQueryBound+2(\MTDepth+1)}{2} \cdot \frac{1}{2^{\ROOutputSize}}$. The proof below improves on this upper bound via a more careful analysis of which collisions are relevant.

\begin{proof}
Let $\ROTrace$ be the query-answer trace of $\MTAdversary^{\ROFunction}$. We define two events:
\begin{itemize}[noitemsep]
  \item $\Event$ is the event that $\MTAdversary$'s output satisfies the conditions in the probability statement; and
  \item $\CollisionEvent$ is the event that $\ROTrace$ contains a collision.
\end{itemize}
Our goal is to upper bound the probability that $\Event$ holds.

We break the analysis into two cases as follows:
\begin{equation*}
\Pr[\Event]
= \Pr[\Event \land \CollisionEvent]
+ \Pr[\Event \land \Negate{\CollisionEvent}]
\enspace.
\end{equation*}
We upper bound each term individually, yielding the claimed upper bound.
\begin{enumerate}

  \item The probability of $\Event \land \CollisionEvent$ is upper bounded by the probability of $\CollisionEvent$, which by \Cref{lemma:rom-cr} is at most $\ROCRExpression{\ROOutputSize}{\ROQueryBound}$.

  \item The event $\Event \land \Negate{\CollisionEvent}$ implies that
\begin{itemize}[nolistsep]
  \item $\ROTrace$ contains no collisions,
  \item $\exists\;\MTMessageIndex \in \MTMessageIndexSet_{0} \cap \MTMessageIndexSet_{1} \text{ s.t. } \MTMessageSubVector_{0}[\MTMessageIndex] \neq \MTMessageSubVector_{1}[\MTMessageIndex]$,
  \item $\MTCheck^{\ROFunction}(\MTCommitment,\MTMessageIndexSet_{0},\MTMessageSubVector_{0},\MTProof_{0})=1$, and
  \item $\MTCheck^{\ROFunction}(\MTCommitment,\MTMessageIndexSet_{1},\MTMessageSubVector_{1},\MTProof_{1})=1$.
\end{itemize}
Let $\MTMessageIndex \in \MTMessageIndexSet_{0} \cap \MTMessageIndexSet_{1}$ be the minimal index such that $\MTMessageSubVector_{0}[\MTMessageIndex] \neq \MTMessageSubVector_{1}[\MTMessageIndex]$. Let $\MTAuthPath_{0}$ and $\MTAuthPath_{1}$ be the authentication paths for location $\MTMessageIndex$ in $\MTProof_{0}$ and $\MTProof_{1}$ respectively. Define the query-answer traces:
\begin{itemize}[noitemsep]
  \item $\ROHonestTrace_{0}$ is the query-answer trace of $\MTCheck^{\ROFunction}(\MTCommitment,\MTMessageIndex,\MTMessageSubVector_{0}[\MTMessageIndex],\MTAuthPath_{0})$; and
  \item $\ROHonestTrace_{1}$ is the query-answer trace of $\MTCheck^{\ROFunction}(\MTCommitment,\MTMessageIndex,\MTMessageSubVector_{1}[\MTMessageIndex],\MTAuthPath_{1})$.
\end{itemize}
Note that $\MTCheck^{\ROFunction}(\MTCommitment,\MTMessageIndexSet_{0},\MTMessageSubVector_{0},\MTProof_{0})=1$ implies that $\MTCheck^{\ROFunction}(\MTCommitment,\MTMessageIndex,\MTMessageSubVector_{0}[\MTMessageIndex],\MTAuthPath_{0})=1$, and similarly $\MTCheck^{\ROFunction}(\MTCommitment,\MTMessageIndexSet_{1},\MTMessageSubVector_{1},\MTProof_{1})=1$ implies that $\MTCheck^{\ROFunction}(\MTCommitment,\MTMessageIndex,\MTMessageSubVector_{1}[\MTMessageIndex],\MTAuthPath_{1})=1$.

By \Cref{lemma:mt-colliding-paths} we get a collision in these traces: there are distinct queries $\ROQuery_{0},\ROQuery_{1}$ and an answer $\ROAnswer$ such that $(\ROQuery_{0},\ROAnswer) \in \ROHonestTrace_{0}$ and $(\ROQuery_{1},\ROAnswer) \in \ROHonestTrace_{1}$. Thus, we show that
\begin{equation*}
\Pr[\Event \land \Negate{\CollisionEvent}]
\leq
\Pr\left[
\begin{array}{c}
\text{$\ROTrace$ contains no collisions and there exist} \\
\text{$(\ROQuery_{0},\ROAnswer) \in \ROHonestTrace_{0}$ and $(\ROQuery_{1},\ROAnswer) \in \ROHonestTrace_{1}$ with $\ROQuery_{0} \neq \ROQuery_{1}$}
\end{array}
\right]
\leq
\frac{(\MTDepth+1)^{2}}{2^{\ROOutputSize}}
\enspace.
\end{equation*}
We are left to argue the second inequality above. Since $\ROTrace$ contains no collisions, it cannot be that $(\ROQuery_{0},\ROAnswer) \in \ROTrace$ and $(\ROQuery_{1},\ROAnswer) \in \ROTrace$ simultaneously. Fix any pair of distinct queries $(\ROQuery_{0},\ROQuery_{1})$ such that $\ROQuery_{0}$ appears as a query in $\ROHonestTrace_{0}$ and $\ROQuery_{1}$ appears as a query in $\ROHonestTrace_{1}$, but not both queries are in $\ROTrace$. Then, the probability that $\ROFunction(\ROQuery_{0}) = \ROFunction(\ROQuery_{1})$ is $\frac{1}{2^{\ROOutputSize}}$. Taking a union bound over all such pairs $(\ROQuery_{0},\ROQuery_{1})$, the probability that there exist $(\ROQuery_{0},\ROAnswer) \in \ROHonestTrace_{0}$ and $(\ROQuery_{1},\ROAnswer) \in \ROHonestTrace_{1}$ with $\ROQuery_{0} \neq \ROQuery_{1}$ is at most $\frac{\Cardinality{\ROHonestTrace_{0}} \cdot \Cardinality{\ROHonestTrace_{1}}}{2^{\ROOutputSize}}$. The claimed bound then follows since $\Cardinality{\ROHonestTrace_{0}},\Cardinality{\ROHonestTrace_{1}} \leq \MTDepth+1$ ($\MTCheck$ makes one query to compute the commitment $\MTVertexLabelByName{\MTDepth}{\MTMessageIndex}$ and then $\MTDepth$ queries to compute the commitments $(\MTVertexLabelByPath{\MTMessageIndex}{\MTLayer})_{\MTLayer \in \{0,1,\dots,\MTDepth-1\}}$).

\end{enumerate}
We conclude that
\begin{align*}
\Pr[\Event]
& = \Pr[\Event \land \CollisionEvent]
+ \Pr[\Event \land \Negate{\CollisionEvent}]
\\ & \leq
\ROCRExpression{\ROOutputSize}{\ROQueryBound} + \frac{(\MTDepth+1)^{2}}{2^{\ROOutputSize}}
\enspace.
\end{align*}
Finally, if $\ROQueryBound \geq 2(\MTDepth+1)^{2}$ then the above expression is at most $\ROCommonExpression{\ROOutputSize}{\ROQueryBound}$.
\end{proof}

We additionally state a lemma that captures binding in the case where the commitment is honestly generated by $\MTCommit$ on a given message, and the adversary wins if the adversary produces a valid opening that is inconsistent with the message. The lemma directly follows from \Cref{lemma:mt-binding}, and we use it in \Cref{chapter:preprocessing-snargs}.

\begin{lemma}
\label{lemma:mt-other-binding}
Let $\MTSymbol \DefineEqual \MTConstructor{\ROOutputSize}{\MTAlphabet}{\MTMessageLength}{\MTSaltSize}$ and define $\MTDepth \DefineEqual \log_{2} \MTMessageLength$. For every query bound $\ROQueryBound \in \Naturals$ and $\ROQueryBound$-query algorithm $\MTAdversary$,
\begin{equation*}
\Pr\left[
\begin{array}{l}
\MTCheck^{\ROFunction}(\MTCommitment,\MTMessageIndexSet,\MTMessageSubVector,\MTProof)=1 \\
\land\;\MTMessageVector[\MTMessageIndexSet] \neq \MTMessageSubVector
\end{array}
\GivenExperiment
\StateExperiment{
\ROFunction \gets \RODistribution{\ROOutputSize} \\
(\MTMessageVector,\ROAdvState) \gets \MTAdversary^{\ROFunction} \\
(\MTCommitment,\MTTrapdoor) \gets \MTCommit^{\ROFunction}(\MTMessageVector) \\
(\MTMessageIndexSet,\MTMessageSubVector,\MTProof) \gets \MTAdversary^{\ROFunction}(\ROAdvState,\MTCommitment,\MTTrapdoor)
}
\right]
\leq
\MTOtherBindingError(\ROOutputSize,\MTMessageLength,\ROQueryBound)
\enspace.
\end{equation*}
Above, $\MTOtherBindingError(\ROOutputSize,\MTMessageLength,\ROQueryBound) \leq \MTHonestBindingExpression{\ROOutputSize}{\ROQueryBound}{\MTMessageLength}$ if $\ROQueryBound \geq 2(\MTDepth+1)^{2}$.
\end{lemma}

\begin{proof}
We reduce to the binding property of $\MTSymbol$ in \Cref{lemma:mt-binding}. Let $\MTAdversary_{2}$ be the algorithm that, using $\MTAdversary$ as a subroutine, works as follows.
\begin{itemize}[noitemsep]
\item[] $\MTAdversary_{2}^{\ROFunction}$:
\begin{enumerate}[noitemsep]
  \item $(\MTMessageVector,\ROAdvState) \gets \MTAdversary^{\ROFunction}$.
  \item $(\MTCommitment,\MTTrapdoor) \gets \MTCommit^{\ROFunction}(\MTMessageVector)$.
  \item $(\MTMessageIndexSet,\MTMessageSubVector,\MTProof_{0}) \gets \MTAdversary^{\ROFunction}(\ROAdvState,\MTCommitment,\MTTrapdoor)$.
  \item $\MTProof_{1} \gets \MTOpen^{\ROFunction}(\MTTrapdoor, \MTMessageIndexSet)$.
  \item Output $(\MTCommitment,
  \MTMessageIndexSet,\MTMessageSubVector,\MTProof_{0},
  \MTMessageIndexSet,\MTMessageVector,\MTProof_{1})$.
\end{enumerate}
\end{itemize}
The algorithm $\MTAdversary_{2}$ makes at most $\ROQueryBound_{2} \DefineEqual \ROQueryBound + 2\MTMessageLength$ queries. Moreover, $\MTCheck^{\ROFunction}(\MTCommitment,\MTMessageIndexSet,\MTMessageVector,\MTProof_{1})=1$ holds always for its output. Therefore,
\begin{align*}
&\Pr\left[
\begin{array}{l}
  \MTCheck^{\ROFunction}(\MTCommitment,\MTMessageIndexSet,\MTMessageSubVector,\MTProof)=1 \\
  \land\;\MTMessageVector[\MTMessageIndexSet] \neq \MTMessageSubVector
\end{array}
\GivenExperiment
\StateExperiment{
  \ROFunction \gets \RODistribution{\ROOutputSize} \\
  (\MTMessageVector,\ROAdvState) \gets \MTAdversary^{\ROFunction} \\
  (\MTCommitment,\MTTrapdoor) \gets \MTCommit^{\ROFunction}(\MTMessageVector) \\
  (\MTMessageIndexSet,\MTMessageSubVector,\MTProof) \gets \MTAdversary^{\ROFunction}(\ROAdvState,\MTCommitment,\MTTrapdoor)
}
\right]
\\ & =
\Pr\left[
\begin{array}{l}
  \exists\;\MTMessageIndex \in \MTMessageIndexSet  \text{ s.t. } \MTMessageSubVector[\MTMessageIndex] \neq \MTMessageVector[\MTMessageIndex] \\
  \land \; \MTCheck^{\ROFunction}(\MTCommitment,\MTMessageIndexSet,\MTMessageSubVector,\MTProof_{0})=1 \\
  \land \; \MTCheck^{\ROFunction}(\MTCommitment,\MTMessageIndexSet,\MTMessageVector,\MTProof_{1})=1
\end{array}
\GivenExperiment
\StateExperiment{
  \ROFunction \gets \RODistribution{\ROOutputSize} \\
  (\MTCommitment,
  \MTMessageIndexSet,\MTMessageSubVector,\MTProof_{0},
  \MTMessageIndexSet,\MTMessageVector,\MTProof_{1})
  \gets \MTAdversary_{2}^{\ROFunction}
}
\right]
\\ & \leq
\MTBindingExpression{\ROOutputSize}{\ROQueryBound_{2}}
\EquationComment{by \Cref{lemma:mt-binding}}
\\ & =
\MTHonestBindingExpression{\ROOutputSize}{\ROQueryBound}{\MTMessageLength}
\enspace.
\end{align*}
\end{proof}

%%%%%%%%%%%%%%%%%%%%%%%%%%%%%%%%%%%%%%%%%%%%%%%%%%%%%%%%%%%%%%%%%%%%%%%%%%%%%%%
%%%%%%%%%%%%%%%%%%%%%%%%%%%%%%%%%%%%%%%%%%%%%%%%%%%%%%%%%%%%%%%%%%%%%%%%%%%%%%%
%%%%%%%%%%%%%%%%%%%%%%%%%%%%%%%%%%%%%%%%%%%%%%%%%%%%%%%%%%%%%%%%%%%%%%%%%%%%%%%
\section{Extractability}
\label{section:merkle-commitment-extractability}

We prove that the Merkle commitment scheme is \emph{extractable}, a security property that informally states that if an adversary outputs a Merkle commitment and then subsequently opens the Merkle commitment at a subset of locations then the adversary \DoQuote{knew} the opening at commitment time. In \Cref{section:merkle-commitment-single-extractability} we study this property for one commitment, in \Cref{section:merkle-commitment-multi-extractability} we study it for multiple commitments, and in \Cref{section:merkle-commitment-multi-configuration-multi-extractability} we study it for multiple commitments across multiple configurations.

%%%%%%%%%%%%%%%%%%%%%%%%%%%%%%%%%%%%%%%%%%%%%%%%%%%%%%%%%%%%%%%%%%%%%%%%%%%%%%%
%%%%%%%%%%%%%%%%%%%%%%%%%%%%%%%%%%%%%%%%%%%%%%%%%%%%%%%%%%%%%%%%%%%%%%%%%%%%%%%
\subsection{Extraction for one commitment}
\label{section:merkle-commitment-single-extractability}

The extractability property is formalized via an efficient algorithm $\MTExtractor$ known as the \emph{extractor}. This is analogous to the extractability property that we proved for the basic commitment scheme $\CMSymbol$ in \Cref{chapter:basic-commitment} (see \Cref{section:basic-commitment-extractability}), but for a major difference. Namely, \emph{we cannot expect to extract an entire message for a given Merkle commitment}.

The adversary may output a Merkle commitment obtained from a partial tree, and then subsequently open parts of this partial tree. Hence $\MTExtractor$ (an efficient algorithm) cannot output an entire message whose Merkle commitment is the one output by the adversary, because that would entail, e.g., inverting the random oracle.

In light of the above, the extractability property only requires $\MTExtractor$ to output a message that agrees with the adversary only at locations that the adversary knows how to open. In particular, if the adversary never provides an opening for certain locations then no requirements are imposed on those locations of the extracted message.

In more detail, we consider an experiment of the following form.
\begin{itemize}

  \item First, in a commitment phase, the adversary, given oracle access to the random oracle, performs a computation and outputs a Merkle commitment $\MTCommitment$ and a private state $\ROAdvState$. Let $\ROTrace$ be the query-answer trace of the adversary in this phase.

  \item Then, in an opening phase, the adversary, given oracle access to the random oracle and as input the private state $\ROAdvState$, continues its computation and outputs an index set $\MTMessageIndexSet$, opening values $\MTMessageSubVector$, and opening proof $\MTProof$.

\end{itemize}
Extractability states that if the adversary's output satisfies $\MTCheck^{\ROFunction}(\MTCommitment,\MTMessageIndexSet,\MTMessageSubVector,\MTProof)=1$ then the extractor $\MTExtractor$, given the Merkle commitment $\MTCommitment$ and the query-answer trace $\ROTrace$, outputs a message $\MTMessageVector$ and opening trapdoor $\MTTrapdoor$ such that $\MTMessageVector$ agrees with $\MTMessageSubVector$ on $\MTMessageIndexSet$ and $\MTTrapdoor$ leads to the opening proof $\MTProof$ when specialized to $\MTMessageIndexSet$. (Up to a small error.) In other words, the adversary \DoQuote{knew} the local opening in the commitment phase because the extractor found it by examining only the query-answer trace of the adversary relevant for producing the commitment (without seeing the subsequent query-answer trace in the opening phase).

The lemma below formally states the property and provides the extraction error.

\begin{lemma}[$\MTSymbol$ is extractable]
\label{lemma:mt-extractability}
Let $\MTSymbol \DefineEqual \MTConstructor{\ROOutputSize}{\MTAlphabet}{\MTMessageLength}{\MTSaltSize}$ and define $\MTDepth \DefineEqual \log_{2} \MTMessageLength$. There exists a deterministic algorithm $\MTExtractor$ such that, for every query bound $\ROQueryBound \in \Naturals$ and $\ROQueryBound$-query algorithm $\MTAdversary$,
\begin{equation*}
\Pr\left[
\begin{array}{l}
\MTCheck^{\ROFunction}(\MTCommitment,\MTMessageIndexSet,\MTMessageSubVector,\MTProof)=1 \\
\land\;(\MTMessageVector[\MTMessageIndexSet] \neq \MTMessageSubVector \lor \MTProof' \neq \MTProof)
\end{array}
\GivenExperiment
\StateExperiment{
\ROFunction \gets \RODistribution{\ROOutputSize} \\
\ROOutputAndTrace{\ROFunction}{\MTAdversary}{\ROTrace}{(\MTCommitment,\ROAdvState)} \\
(\MTMessageVector,\MTTrapdoor) \gets \MTExtractor(\MTCommitment,\ROTrace) \\
(\MTMessageIndexSet,\MTMessageSubVector,\MTProof) \gets \MTAdversary^{\ROFunction}(\ROAdvState) \\
\MTProof' \gets \MTOpen^{\ROFunction}(\MTTrapdoor, \MTMessageIndexSet)
}
\right]
\leq
\MTExtractionError(\ROOutputSize,\MTMessageLength,\ROQueryBound)
\enspace.
\end{equation*}
Above:
\begin{itemize}[nolistsep]
  \item $\MTExtractionError(\ROOutputSize,\MTMessageLength,\ROQueryBound) \leq \MTExtractabilityExpression{\ROOutputSize}{\ROQueryBound}{\MTMessageLength}{\MTDepth}$ (at most $\MTExtractabilityShortExpression{\ROOutputSize}{\ROQueryBound}$ if $\MTExtractabilityCondition{\ROQueryBound}{\MTMessageLength}{\MTDepth}$);
  \item $\MTExtractor$ runs in time $\MTExtractorTimeFunction{\ROOutputSize}{\MTAlphabet}{\MTMessageLength}{\MTSaltSize}{\ROQueryBound} = O\big( \ROQueryBound \cdot \MTMessageLength \cdot (\log\Cardinality{\MTAlphabet} + \MTSaltSize + \ROOutputSize)\big)$.
\end{itemize}
\end{lemma}

We describe the intuition behind the extractor $\MTExtractor$, and then provide a formal description of it. The extractor $\MTExtractor$ is tasked with finding a message that \DoQuote{explains} any future opening to a given Merkle commitment $\MTCommitment$, given the query-answer trace of the adversary that produced the Merkle commitment. Intuitively, all the information about possible openings of the adversary is available in the query-answer trace of the adversary's computation till it outputs the Merkle commitment $\MTCommitment$. So the extractor merely needs to organize the query-answer trace into a partial tree, and read off the leaves of this tree, filling in with an arbitrary value any missing leaves. Because a random oracle is inversion resistant and collision resistant, no adversary is able to provide any other openings other than those contained in the leaves output this way, up to a small probability of error.

In order to define the extractor, it will be convenient to partition the queries into three types: leaf queries (for computing the leaves of the Merkle tree), inner vertex queries (for computing the internal vertices of the tree), and other (queries that do not match these previous categories). After that we describe the extractor $\MTExtractor$.

\begin{definition}
\label{definition:query-types}
Let $\ROTrace = \big( (\ROQuery_{i},\ROAnswer_{i}) \big)_{i \in [\ROQueryBound]}$ be a query-answer trace for an oracle $\ROFunction \colon \Bits^{*} \to \Bits^{\ROOutputSize}$. We partition the query-answer pairs in $\ROTrace$ according to three \defemph{types of queries}.
\begin{enumerate}[nolistsep]
  \item \emph{Leaf vertex queries.}
  A set $\MTLeafQueries$ that contains queries-answer pairs $(\ROQuery,\ROAnswer)$ where $\ROQuery=(\MTMessageEntry,\MTSaltString)$ for some message entry $\MTMessageEntry \in \MTAlphabet$ and salt $\MTSaltString \in \Bits^{\MTSaltSize}$.
  \item \emph{Inner vertex queries.}
  A set $\MTInternalQueries$ that contains queries-answer pairs $(\ROQuery,\ROAnswer)$ where $\ROQuery \in \Bits^{2\ROOutputSize}$.
  \item \emph{Other queries.}
  A set $\MTOtherQueries$ that contains all other query-answer pairs.
\end{enumerate}
\end{definition}

\begin{construction}
\label{construction:mt-extractor}
Given as input a Merkle commitment $\MTCommitment \in \Bits^{\ROOutputSize}$ and query-answer trace $\ROTrace = \big( (\ROQuery_{i},\ROAnswer_{i}) \big)_{i \in [\ROQueryBound]}$, $\MTExtractor$ works as follows.
\begin{enumerate}[noitemsep]

\item If $\MTCommitment$ is not the answer to any query in $\ROTrace$, then output $(\MTMessageVector,\MTTrapdoor) \DefineEqual (\bot,\bot)$.

\item \label{step:mt-extractor-partition-queries}
Let $\MTLeafQueries,\MTInternalQueries,\MTOtherQueries$ be the partitioning of $\ROTrace$ according to \Cref{definition:query-types}.

\item \label{step:mt-extractor-build-tree}
Label the binary tree $\MTGraph_{\MTMessageLength}$ (see \Cref{definition:mt-graph}) as follows:
\begin{itemize}
  \item The root of $\MTGraph_{\MTMessageLength}$ is labeled with $\MTCommitment$.
  \item While there is a query-answer pair $(\ROQuery, \ROAnswer) \in \MTInternalQueries$ where $\ROAnswer$ is a label of a vertex in $\MTGraph_{\MTMessageLength}$, set the label of the left child to be $\StringLeftHalf{\ROQuery}$ and the label of the right child to be $\StringRightHalf{\ROQuery}$ (where $\StringLeftHalf{\ROQuery}$ and $\StringRightHalf{\ROQuery}$ are the left half and right half of $\ROQuery$). Remove $(\ROQuery,\ROAnswer)$ from $\MTInternalQueries$.
  \item While there is a query-answer pair $(\ROQuery,\ROAnswer) \in \MTLeafQueries$ where $\ROAnswer$ is the label of the $\MTMessageIndex$-th leaf of $\MTGraph_{\MTMessageLength}$ for some $\MTMessageIndex \in [\MTMessageLength]$, set $\MTMessageEntry_{\MTMessageIndex} \DefineEqual \MTMessageEntry$ and $\MTSaltString_{\MTMessageIndex} \DefineEqual \MTSaltString$, where $\ROQuery = (\MTMessageEntry,\MTSaltString)$. Remove $(\ROQuery,\ROAnswer)$ from $\MTLeafQueries$.
\end{itemize}

\item \label{step:mt-extractor-fill-rest}
Set arbitrarily the values of any label in $\MTGraph_{\MTMessageLength}$ or pair $(\MTMessageEntry_{\MTMessageIndex},\MTSaltString_{\MTMessageIndex})$ not defined in the previous step.

\item Set the commitments $\MTVertexLabels \DefineEqual ((\MTVertexLabelByName{\MTLayer}{\MTMessageIndex})_{\MTMessageIndex \in [2^{\MTLayer}]})_{\MTLayer=0}^{\MTDepth}$, where $\MTVertexLabelByName{\MTLayer}{\MTMessageIndex}$ is the label of the $\MTMessageIndex$-th vertex in layer $\MTLayer$ of the tree $\MTGraph_{\MTMessageLength}$.

\item Set the message vector $\MTMessageVector \DefineEqual (\MTMessageEntry_{\MTMessageIndex})_{\MTMessageIndex \in [\MTMessageLength]}$.

\item Set the opening trapdoor $\MTTrapdoor \DefineEqual (\MTSaltStrings,\MTVertexLabels)$, where $\MTSaltStrings \DefineEqual (\MTSaltString_{\MTMessageIndex})_{\MTMessageIndex \in [\MTMessageLength]}$.

\item Output $(\MTMessageVector,\MTTrapdoor)$.

\end{enumerate}
\end{construction}


\begin{proof}[Proof of \Cref{lemma:mt-extractability}]
We discuss the time complexity of $\MTExtractor$ in \Cref{remark:mt-extractor-time}. Below we analyze its success probability.

Consider the following query-answer traces:
\begin{itemize}[noitemsep]
  \item $\ROTrace$ is the query-answer trace of the execution $\ROOutputAndTrace{\ROFunction}{\MTAdversary}{\ROTrace}{(\MTCommitment,\ROAdvState)}$;
  \item $\ROOpenTrace$ is the query-answer trace of the execution
  $\ROInputOutputAndTrace{\ROFunction}{\MTAdversary}{\ROAdvState}{\ROOpenTrace}{(\MTMessageIndexSet,\MTMessageSubVector,\MTProof)}$;
  \item $\ROHonestTrace$ is the query-answer trace of the execution $\MTCheck^{\ROFunction}(\MTCommitment,\MTMessageIndexSet,\MTMessageSubVector,\MTProof)$.
\end{itemize}

Fix any $\ROQueryBound_{1},\ROQueryBound_{2} \in \Naturals$ with $\ROQueryBound_{1} + \ROQueryBound_{2} \leq \ROQueryBound$. We prove the claimed upper bound conditioned on the event that $\ROQueryBound_{1} = \Cardinality{\ROTrace}$ and $\ROQueryBound_{2} = \Cardinality{\ROOpenTrace}$. The lemma follows from this because the upper bound holds for every $\ROQueryBound_{1},\ROQueryBound_{2} \in \Naturals$ with $\ROQueryBound_{1} + \ROQueryBound_{2} \leq \ROQueryBound$, and every computation of the $\ROQueryBound$-query adversary $\MTAdversary$ implies a setting of $\ROQueryBound_{1},\ROQueryBound_{2} \in \Naturals$ with $\ROQueryBound_{1} + \ROQueryBound_{2} \leq \ROQueryBound$.

We define several events:
\begin{itemize}[noitemsep]
  \item $\Event$ is the event that the conditions in the probability statement hold;
  \item $\CollisionEvent$ is the event that $\ROTrace$ contains a collision;
  \item $\NotEqualEvent$ is the event that $\MTGraph_{\MTMessageLength} \neq \hat{\MTGraph_{\MTMessageLength}}$ where
  \begin{itemize}[nolistsep]
    \item $\MTGraph_{\MTMessageLength}$ is the tree generated in the execution of $\MTExtractor(\MTCommitment,\ROTrace)$ and
    \item $\hat{\MTGraph_{\MTMessageLength}}$ is the tree generated in the execution of $\MTExtractor(\MTCommitment,\ROTrace \concat \ROOpenTrace)$.
  \end{itemize}
  \item $\SubsetEvent$ is the event that $\ROHonestTrace \not \subseteq \ROTrace$ and $\MTCheck^{\ROFunction}(\MTCommitment,\MTMessageIndexSet,\MTMessageSubVector,\MTProof)=1$.
\end{itemize}
Our goal is to upper bound the probability that $\Event$ holds. We break the analysis into four cases:
\begin{align*}
\Pr[\Event]
\leq
\Pr[\CollisionEvent]
+ \Pr\left[\NotEqualEvent \ConditionedOn \Negate{\CollisionEvent}\right]
+ \Pr\left[\SubsetEvent \ConditionedOn \Negate{\CollisionEvent} \land \Negate{\NotEqualEvent} \right]
+ \Pr\left[\Event \ConditionedOn \Negate{\CollisionEvent} \land \Negate{\NotEqualEvent} \land \Negate{\SubsetEvent}\right]
\enspace.
\end{align*}
We upper bound each term individually, yielding the claimed upper bound.
\begin{enumerate}

  \item An upper bound on the probability of the event $\CollisionEvent$ follows directly from \Cref{lemma:rom-cr}:
  \begin{equation*}
\Pr[\CollisionEvent]
\leq
\ROCRExpression{\ROOutputSize}{\ROQueryBound_{1}}
\enspace.
\end{equation*}

 \item We upper bound the probability of the event $\left.\NotEqualEvent \ConditionedOn \Negate{\CollisionEvent}\right.$.

If $\MTGraph_{\MTMessageLength} \neq \hat{\MTGraph_{\MTMessageLength}}$ then there exists a query in $\ROOpenTrace$ that is not in $\ROTrace$ with an answer that equals a non-dummy label in $\MTGraph_{\MTMessageLength}$. (A non-dummy label is one added in \Cref{step:mt-extractor-build-tree} of $\MTExtractor$, rather than \Cref{step:mt-extractor-fill-rest}.) We bound the probability that such a query exists.

Since the event $\CollisionEvent$ does not hold ($\ROTrace$ contains no collisions), each vertex in the tree $\MTGraph_{\MTMessageLength}$ is labeled at most once. There are at most $\min\{2\ROQueryBound_{1}+1,2\MTMessageLength\}$ non-dummy labels in $\MTGraph_{\MTMessageLength}$: the root vertex is labeled separately, and each query in $\ROTrace$ results in at most $2$ newly-labeled vertices; moreover, there are at most $2\MTMessageLength$ vertices in the tree $\MTGraph_{\MTMessageLength}$. Hence the probability that any fixed query in $\ROOpenTrace$ (that is not in $\ROTrace$) has an answer that equals one of the labels in $\MTGraph_{\MTMessageLength}$ is at most $\frac{\min\{2\ROQueryBound_{1}+1,2\MTMessageLength\}}{2^{\ROOutputSize}}$.

Taking a union bound over all queries in $\ROOpenTrace$,
\begin{align*}
\Pr\left[\NotEqualEvent \ConditionedOn \Negate{\CollisionEvent}\right]
\leq
\Cardinality{\ROOpenTrace} \cdot \frac{\min\{2\ROQueryBound_{1}+1,2\MTMessageLength\}}{2^{\ROOutputSize}}
=
\ROQueryBound_{2} \cdot \frac{\min\{2\ROQueryBound_{1}+1,2\MTMessageLength\}}{2^{\ROOutputSize}}
\enspace.
\end{align*}

 \item We upper bound the probability of the event $\left.\SubsetEvent \ConditionedOn \Negate{\CollisionEvent} \land \Negate{\NotEqualEvent} \right.$.

For every $\MTMessageIndex \in \MTMessageIndexSet$, let $\ROHonestTrace_{\MTMessageIndex} \subseteq \ROHonestTrace$ be the query-answer trace of $\MTCheck^{\ROFunction}(\MTCommitment,\MTMessageIndex,\MTMessageSubVector[\MTMessageIndex],\MTAuthPath_{\MTMessageIndex})$, where $\MTAuthPath_{\MTMessageIndex}$ is the authentication path for location $\MTMessageIndex$ in $\MTProof$.

Suppose that $\SubsetEvent$ holds. Then $\ROHonestTrace \not \subseteq \ROTrace$, so there exists $\MTMessageIndex \in \MTMessageIndexSet$ such that $\ROHonestTrace_{\MTMessageIndex} \not \subseteq \ROTrace$. Moreover, $\MTCheck^{\ROFunction}(\MTCommitment,\MTMessageIndex,\MTMessageSubVector[\MTMessageIndex],\MTAuthPath_{\MTMessageIndex})=1$, so the query-answer trace $\ROHonestTrace_{\MTMessageIndex}$ can be viewed as a path that \DoQuote{ends} in the Merkle commitment $\MTCommitment$, and thus \DoQuote{connects} somewhere into $\MTGraph_{\MTMessageLength}$ (the tree generated in the execution of $\MTExtractor(\MTCommitment,\ROTrace)$). In sum, if $\SubsetEvent$ holds then there exists a query in $\ROHonestTrace_{\MTMessageIndex}$ that is not in $\ROTrace$ with an answer that equals a non-dummy label in $\MTGraph_{\MTMessageLength}$. (A non-dummy label is one added in \Cref{step:mt-extractor-build-tree} of $\MTExtractor$, rather than \Cref{step:mt-extractor-fill-rest}.)

Since the event $\CollisionEvent$ does not hold ($\ROTrace$ contains no collisions), each vertex in the tree $\MTGraph_{\MTMessageLength}$ is labeled at most once. There are at most $\min\{2\ROQueryBound_{1}+1,2\MTMessageLength\}$ non-dummy labels in $\MTGraph_{\MTMessageLength}$ (we argued this in the prior point). Moreover, $\MTGraph_{\MTMessageLength} = \hat{\MTGraph_{\MTMessageLength}}$ because $\NotEqualEvent$ does not hold. Hence there are no queries in $\ROOpenTrace \setminus \ROTrace$ that have an answer that equals a non-dummy label in $\MTGraph_{\MTMessageLength}$ (otherwise, there would be a leaf in $\hat{\MTGraph_{\MTMessageLength}}$ labeled differently from the corresponding leaf in $\MTGraph_{\MTMessageLength}$). Hence the probability that any fixed query in $\ROHonestTrace_{\MTMessageIndex}$ (that is not in $\ROTrace \cup \ROOpenTrace$) has an answer that equals a non-dummy label in $\MTGraph_{\MTMessageLength}$ is at most $\frac{\min\{2\ROQueryBound_{1}+1,2\MTMessageLength\}}{2^{\ROOutputSize}}$.

Considering any choice of $\MTMessageIndexSet \subseteq [\MTMessageLength]$ and $\MTMessageIndex \in \MTMessageIndexSet$ and taking a union bound over all queries in $\ROHonestTrace_{\MTMessageIndex}$,
\begin{align*}
\Pr\left[\SubsetEvent \ConditionedOn \Negate{\CollisionEvent} \land \Negate{\NotEqualEvent} \right]
\leq
\max_{\MTMessageIndexSet \subseteq [\MTMessageLength]}
\max_{\MTMessageIndex \in \MTMessageIndexSet}
\Cardinality{\ROHonestTrace_{\MTMessageIndex}} \cdot \frac{\min\{2\ROQueryBound_{1}+1,2\MTMessageLength\}}{2^{\ROOutputSize}}
\leq
(\MTDepth+1) \cdot \frac{\min\{2\ROQueryBound_{1}+1,2\MTMessageLength\}}{2^{\ROOutputSize}}
\enspace.
\end{align*}

  \item We show that the event $\Event$ cannot happen if $\CollisionEvent$ does not hold, $\NotEqualEvent$ does not hold, and $\SubsetEvent$ does not hold (in fact, it suffices to assume $\Negate{\CollisionEvent} \land \Negate{\SubsetEvent}$ but we include $\Negate{\NotEqualEvent}$ for symmetry). We provide this as a claim (as we reuse it in later sections).

\begin{claim}
\label{lemma:mt-extractability-same-trees}
It holds that
\begin{equation*}
\Pr\left[
\begin{array}{l}
  \MTCheck^{\ROFunction}(\MTCommitment,\MTMessageIndexSet,\MTMessageSubVector,\MTProof)=1 \\
  \land\;(\MTMessageVector[\MTMessageIndexSet] \neq \MTMessageSubVector \lor \MTProof' \neq \MTProof) \\
  \ConditionedOnText \\
  \Negate{\CollisionEvent} \land \Negate{\NotEqualEvent} \land \Negate{\SubsetEvent}
\end{array}
\GivenExperiment
\StateExperiment{
  \ROFunction \gets \RODistribution{\ROOutputSize} \\
  \ROOutputAndTrace{\ROFunction}{\MTAdversary}{\ROTrace}{(\MTCommitment,\ROAdvState)} \\
  (\MTMessageVector,\MTTrapdoor) \gets \MTExtractor(\MTCommitment,\ROTrace) \\
  (\MTMessageIndexSet,\MTMessageSubVector,\MTProof) \gets \MTAdversary^{\ROFunction}(\ROAdvState) \\
  \MTProof' \gets \MTOpen^{\ROFunction}(\MTTrapdoor, \MTMessageIndexSet)
}
\right]
=0
\enspace.
\end{equation*}
\end{claim}

\begin{proof}
The conditioning means that all of the events $\CollisionEvent,\NotEqualEvent,\SubsetEvent$ do not hold. Suppose that $\MTCheck^{\ROFunction}(\MTCommitment,\MTMessageIndexSet,\MTMessageSubVector,\MTProof)=1$ (as otherwise the claim holds trivially). Then, $\ROHonestTrace \subseteq \ROTrace$ (since $\SubsetEvent$ does not hold). Below we argue that (under the conditioning) this implies that $\MTMessageVector[\MTMessageIndexSet] = \MTMessageSubVector$ and $\MTProof' = \MTProof$. Note that the opening proofs $\MTProof$ and $\MTProof'$ consist of authentication paths $(\MTAuthPath_{\MTMessageIndex})_{\MTMessageIndex \in \MTMessageIndexSet}$ and $(\MTAuthPath'_{\MTMessageIndex})_{\MTMessageIndex \in \MTMessageIndexSet}$. Moreover, from \Cref{equation:mt-check-and-for-each-location} we know that
\begin{align*}
&\MTCheck^{\ROFunction}(\MTCommitment,\MTMessageIndexSet,\MTMessageSubVector,\MTProof)
=
\idxland{\MTMessageIndex \in \MTMessageIndexSet}\, \MTCheck^{\ROFunction}(\MTCommitment,\MTMessageIndex,\MTMessageSubVector[\MTMessageIndex],\MTAuthPath_{\MTMessageIndex})
\enspace, \text{ and} \\
&\MTCheck^{\ROFunction}(\MTCommitment,\MTMessageIndexSet,\MTMessageVector[\MTMessageIndexSet],\MTProof')
=
\idxland{\MTMessageIndex \in \MTMessageIndexSet}\, \MTCheck^{\ROFunction}(\MTCommitment,\MTMessageIndex,\MTMessageVector[\MTMessageIndex],\MTAuthPath'_{\MTMessageIndex})
\enspace.
\end{align*}
\begin{itemize}

  \item We show that $\MTMessageVector[\MTMessageIndexSet] = \MTMessageSubVector$.

Fix $\MTMessageIndex \in \MTMessageIndexSet$ and consider the following query-answer traces:
\begin{itemize}[nolistsep]
  \item $\ROHonestTrace_{\MTMessageIndex}$ is the query-answer trace of the execution $\MTCheck^{\ROFunction}(\MTCommitment,\MTMessageIndex,\MTMessageSubVector[\MTMessageIndex],\MTAuthPath_{\MTMessageIndex})$;
  \item $\ROHonestTrace'_{\MTMessageIndex}$ is the query-answer trace of the execution $\MTCheck^{\ROFunction}(\MTCommitment,\MTMessageIndex,\MTMessageVector[\MTMessageIndex],\MTAuthPath'_{\MTMessageIndex})$.
\end{itemize}

We know that $\ROHonestTrace_{\MTMessageIndex} \subseteq \ROHonestTrace$ (since $\MTCheck$ validates each authentication path in the opening proof), so $\ROHonestTrace_{\MTMessageIndex} \subseteq \ROTrace$. Moreover, we know that  $\MTCheck^{\ROFunction}(\MTCommitment,\MTMessageIndex,\MTMessageSubVector[\MTMessageIndex],\MTAuthPath_{\MTMessageIndex})=1$ (since $\MTCheck^{\ROFunction}(\MTCommitment,\MTMessageIndexSet,\MTMessageSubVector,\MTProof)=1$).

Next, the extractor $\MTExtractor$, given input $(\MTCommitment,\ROTrace)$, uses the query-answer pairs in $\ROTrace$ to assign values to the entries of $\MTMessageVector$ that have a valid authentication paths in $\ROTrace$ to the Merkle commitment $\MTCommitment$ (see \Cref{step:mt-extractor-build-tree} of \Cref{construction:mt-extractor}); the remaining entries of $\MTMessageVector$ are assigned arbitrary values (see \Cref{step:mt-extractor-fill-rest} of \Cref{construction:mt-extractor}). Since $\ROHonestTrace_{\MTMessageIndex} \subseteq \ROTrace$, we know that the $\MTMessageIndex$-th entry of $\MTMessageVector$ is assigned a non-arbitrary value (as the trace $\ROTrace$ contains a path from the $\MTMessageIndex$-th leaf to the root), and hence $\ROHonestTrace'_{\MTMessageIndex} \subseteq \ROTrace$. Moreover, the opening trapdoor $\MTTrapdoor$ output by $\MTExtractor$ enables $\MTOpen$ to output valid authentication paths for entries that are assigned non-arbitrary values, including the authentication path $\MTAuthPath'_{\MTMessageIndex}$, so we also deduce that $\MTCheck^{\ROFunction}(\MTCommitment,\MTMessageIndex,\MTMessageVector[\MTMessageIndex],\MTAuthPath'_{\MTMessageIndex})=1$.

We have established that:
\begin{itemize}[nolistsep]
  \item $\ROHonestTrace_{\MTMessageIndex}$ is contained in $\ROTrace$, and $\MTCheck^{\ROFunction}(\MTCommitment,\MTMessageIndex,\MTMessageSubVector[\MTMessageIndex],\MTAuthPath_{\MTMessageIndex})=1$; and
  \item $\ROHonestTrace'_{\MTMessageIndex}$ is contained in $\ROTrace$, and $\MTCheck^{\ROFunction}(\MTCommitment,\MTMessageIndex,\MTMessageVector[\MTMessageIndex],\MTAuthPath'_{\MTMessageIndex})= 1$.
\end{itemize}

Now suppose by way of contradiction that $\MTMessageVector[\MTMessageIndexSet] \neq \MTMessageSubVector$. By \Cref{lemma:mt-colliding-paths}, $\ROTrace$ contains a collision, which contradicts the conditioning.

  \item We show that $\MTProof' = \MTProof$.

Suppose by way of contradiction that $\MTProof' \neq \MTProof$. This means that there exists an index $\MTMessageIndex \in \MTMessageIndexSet$ such that $\MTAuthPath'_{\MTMessageIndex} \neq \MTAuthPath_{\MTMessageIndex}$. Above we argued that $\MTMessageVector[\MTMessageIndexSet] = \MTMessageSubVector$, which means that $\MTMessageVector[\MTMessageIndex] = \MTMessageSubVector[\MTMessageIndex]$. Above we also argued that $\MTCheck^{\ROFunction}(\MTCommitment,\MTMessageIndex,\MTMessageVector[\MTMessageIndex],\MTAuthPath'_{\MTMessageIndex})=1$ and that $\MTCheck^{\ROFunction}(\MTCommitment,\MTMessageIndex,\MTMessageSubVector[\MTMessageIndex],\MTAuthPath_{\MTMessageIndex})=1$, and that the corresponding traces $\ROHonestTrace'_{\MTMessageIndex}$ and $\ROHonestTrace_{\MTMessageIndex}$ are in $\ROTrace$. By \Cref{lemma:simple-mt-colliding-paths}, $\ROTrace$ contains a collision, which contradicts the conditioning.

\end{itemize}
\end{proof}
\end{enumerate}
We conclude that
\begin{align*}
&\Pr[\Event]
\\ & \leq
\Pr[\CollisionEvent]
+ \Pr\left[\NotEqualEvent \ConditionedOn \Negate{\CollisionEvent}\right]
+ \Pr\left[\SubsetEvent \ConditionedOn \Negate{\CollisionEvent} \land \Negate{\NotEqualEvent}\right]
+ \Pr\left[\Event \ConditionedOn \Negate{\CollisionEvent} \land \Negate{\NotEqualEvent} \land \Negate{\SubsetEvent}\right]
\\ & \leq
\ROCRExpression{\ROOutputSize}{\ROQueryBound_{1}}
+ \frac{\ROQueryBound_{2} \cdot \min\{2\ROQueryBound_{1}+1,2\MTMessageLength\}}{2^{\ROOutputSize}}
+ \frac{(\MTDepth+1) \cdot \min\{2\ROQueryBound_{1}+1,2\MTMessageLength\}}{2^{\ROOutputSize}}
+ 0
\\ & \leq
\ROCRExpression{\ROOutputSize}{\ROQueryBound_{1}}
+ \frac{(\ROQueryBound-\ROQueryBound_{1} + \MTDepth+1) \cdot \min\{2\ROQueryBound_{1}+1,2\MTMessageLength\}}{2^{\ROOutputSize}}
\EquationComment{since $\ROQueryBound_{1} + \ROQueryBound_{2} \leq \ROQueryBound$}
\\ & \leq
\ROCRExpression{\ROOutputSize}{\ROQueryBound_{1}}
+ \frac{(\ROQueryBound-\ROQueryBound_{1} + \MTDepth+1) \cdot 2\MTMessageLength}{2^{\ROOutputSize}}
\enspace.
\EquationComment{since $\min\{a,b\}\leq a,b$}
\end{align*}
The expression is a convex function in $\ROQueryBound_{1}$ so its maximum on the interval $\ROQueryBound_{1} \in \{0,1,\dots,\ROQueryBound\}$ is achieved at either $\ROQueryBound_{1}=0$ or $\ROQueryBound_{1} = \ROQueryBound$. Hence the expression is upper bounded from the maximum of these two options:
\begin{equation*}
\Pr[\Event]
\leq
\max\left\{
\frac{(\ROQueryBound + \MTDepth+1) \cdot 2\MTMessageLength}{2^{\ROOutputSize}}
,
\MTExtractabilityExpression{\ROOutputSize}{\ROQueryBound}{\MTMessageLength}{\MTDepth}
\right\}
\enspace.
\end{equation*}
We conclude that if $\ROQueryBound \geq 4\MTMessageLength+1$ (a very mild condition) then the first term is no more than the second term, in which case
\begin{equation*}
\Pr[\Event]
\leq
\MTExtractabilityExpression{\ROOutputSize}{\ROQueryBound}{\MTMessageLength}{\MTDepth}
\enspace.
\end{equation*}
\end{proof}

\begin{remark}[time complexity of $\MTExtractor$]
\label{remark:mt-extractor-time}
We discuss how a straightforward realization of $\MTExtractor$ leads to the time complexity claimed in \Cref{lemma:mt-extractability}, and how in common parameter regimes this time complexity can be improved.

The dominant contribution comes from \Cref{step:mt-extractor-build-tree}, which we discuss below. Indeed, all other steps (partitioning queries, giving arbitrary values to unlabeled vertices in the tree, and assembling the extractor outputs) can be performed in time at most $O\big((\ROQueryBound + \MTMessageLength) \cdot (\log\Cardinality{\MTAlphabet} + \MTSaltSize + \ROOutputSize)\big)$.

\Cref{step:mt-extractor-build-tree} requires structuring certain information from the query-trace into a partial binary tree. This involves at most $\min\{2\MTMessageLength,\ROQueryBound\}=O(\MTMessageLength)$ iterations (each iteration corresponds to a vertex in the tree and an entry in the trace), proceeding layer by layer from the root towards the leaves. Each iteration involves searching the query-answer trace for an entry in which the answer equals the label of the vertex. Each of these $O(\MTMessageLength)$ searches takes at most time $O\big(\ROQueryBound \cdot (\log\Cardinality{\MTAlphabet} + \MTSaltSize + \ROOutputSize)\big)$, which leads to the time complexity $O\big( \ROQueryBound \cdot \MTMessageLength \cdot (\log\Cardinality{\MTAlphabet} + \MTSaltSize + \ROOutputSize)\big)$ claimed in \Cref{lemma:mt-extractability}.

We can do better by using suitable data structures. Realizing \Cref{step:mt-extractor-build-tree} essentially requires a \emph{static dictionary} because we need to perform $\min\{2\MTMessageLength,\ROQueryBound\}$ lookups into a (static) list of at most $\ROQueryBound$ entries, with each entry consisting of at most $\MTTraceEntrySize \DefineEqual O(\log\Cardinality{\MTAlphabet} + \MTSaltSize + \ROOutputSize)$ bits. One way to achieve this is to sort the query-answer trace by answers, and then perform lookups via binary search. Sorting the trace takes time $O(\ROQueryBound \cdot \log \ROQueryBound \cdot \MTTraceEntrySize)$ and performing each of the $\min\{2\MTMessageLength,\ROQueryBound\}$ lookups takes time $O(\log \ROQueryBound \cdot \MTTraceEntrySize)$; this leads to a total time complexity of $O((\ROQueryBound+\min\{2\MTMessageLength,\ROQueryBound\}) \cdot \log \ROQueryBound \cdot \MTTraceEntrySize) = O(\ROQueryBound \cdot \log \ROQueryBound \cdot \MTTraceEntrySize)$.

The above expression improves on the time in \Cref{lemma:mt-extractability} provided that $\log \ROQueryBound = O(\MTMessageLength)$.
\end{remark}

\begin{remark}[extractability implies binding]
\label{remark:extractability-implies-binding}
The extractability property in \Cref{lemma:mt-extractability} is a strong security guarantee: any bounded-query adversary that opens a commitment must have \DoQuote{known} the opening at commitment time. In fact, \emph{extractability is stronger than binding}: we show that if $\MTSymbol$ satisfies extractability with error $\MTExtractionError$ then $\MTSymbol$ satisfies binding with error $\MTBindingError \leq 2 \cdot \MTExtractionError$ (a closely related error).\footnote{This implication loses a factor of $2$, and we present it merely for didactic purposes. Indeed, for $\MTSymbol$ we prove almost the same errors in \Cref{lemma:mt-binding} (binding error) and \Cref{lemma:mt-extractability} (extraction error).}

Let $\MTAdversary$ be a $\ROQueryBound$-query algorithm that opens a commitment in two ways with probability $\MTBindingError$:
\begin{equation*}
\MTBindingError
\DefineEqual
\Pr\left[
\begin{array}{l}
\exists\;\MTMessageIndex \in \MTMessageIndexSet_{0} \cap \MTMessageIndexSet_{1} \text{ s.t. } \MTMessageSubVector_{0}[\MTMessageIndex] \neq \MTMessageSubVector_{1}[\MTMessageIndex] \\
\land \; \MTCheck^{\ROFunction}(\MTCommitment,\MTMessageIndexSet_{0},\MTMessageSubVector_{0},\MTProof_{0})=1 \\
\land \; \MTCheck^{\ROFunction}(\MTCommitment,\MTMessageIndexSet_{1},\MTMessageSubVector_{1},\MTProof_{1})=1
\end{array}
\GivenExperiment
\StateExperiment{
\ROFunction \gets \RODistribution{\ROOutputSize} \\
(\MTCommitment,
\MTMessageIndexSet_{0},\MTMessageSubVector_{0},\MTProof_{0},
\MTMessageIndexSet_{1},\MTMessageSubVector_{1},\MTProof_{1})
\gets \MTAdversary^{\ROFunction}
}
\right]
\enspace.
\end{equation*}

Let $\MTAdversary_{e}$ be the $\ROQueryBound$-query algorithm for the extraction game that is defined as follows.
\begin{itemize}[noitemsep]
  \item $\MTAdversary_{e}^{\ROFunction}$: Run $\MTAdversary^{\ROFunction}$ to get $(\MTCommitment,
\MTMessageIndexSet_{0},\MTMessageSubVector_{0},\MTProof_{0},
\MTMessageIndexSet_{1},\MTMessageSubVector_{1},\MTProof_{1})$, set $\ROAdvState \DefineEqual (\MTMessageIndexSet_{0},\MTMessageSubVector_{0},\MTProof_{0},
\MTMessageIndexSet_{1},\MTMessageSubVector_{1},\MTProof_{1})$, and output $(\MTCommitment,\ROAdvState)$.
  \item $\MTAdversary_{e}^{\ROFunction}(\ROAdvState)$: Sample a random bit $b$ (by using private randomness or via a fresh additional query to the random oracle) and output the tuple $(\MTMessageIndexSet_{b},\MTMessageSubVector_{b},\MTProof_{b})$ stored in $\ROAdvState$.
\end{itemize}
Fix any candidate extractor $\MTExtractor$. Fix any output $(\MTCommitment,\MTMessageIndexSet_{0},\MTMessageSubVector_{0},\MTProof_{0},\MTMessageIndexSet_{1},\MTMessageSubVector_{1},\MTProof_{1})$  of $\MTAdversary$ that breaks the binding property (i.e., satisfies the conditions in the probability statement above). If we run $\MTExtractor$ on $\MTCommitment$ (the commitment output by $\MTAdversary_{e}$ in the first step) and $\ROTrace$ (the query-answer trace by $\MTAdversary_{e}$ in the first step) then we obtain an output $(\MTMessageVector,\MTTrapdoor)$ that, with probability at least $1/2$, satisfies $\MTMessageVector[\MTMessageIndexSet] \neq \MTMessageSubVector_{b} \lor \MTProof \neq \MTProof_{b}$, where $\MTProof \DefineEqual \MTOpen^{\ROFunction}(\MTTrapdoor, \MTMessageIndexSet_{b})$. This is because the inputs to the extractor $\MTExtractor$ are independent of the bit $b$ (which is sampled by $\MTAdversary_{e}$ after the inputs to $\MTExtractor$ are determined). Moreover, by the definition of $\MTAdversary_{e}$, we know that $\MTCheck^{\ROFunction}(\MTCommitment,\MTMessageIndexSet_{b},\MTMessageSubVector_{b},\MTProof_{b})=1$ regardless of the choice of bit $b$. We deduce that
\begin{equation*}
\frac{1}{2} \cdot \MTBindingError \leq
\Pr\left[
\begin{array}{l}
  \MTCheck^{\ROFunction}(\MTCommitment,\MTMessageIndexSet_{b},\MTMessageSubVector_{b},\MTProof_{b})=1 \\
  \land\;(\MTMessageVector[\MTMessageIndexSet] \neq \MTMessageSubVector_{b} \lor \MTProof \neq \MTProof_{b})
\end{array}
\GivenExperiment
\StateExperiment{
  \ROFunction \gets \RODistribution{\ROOutputSize} \\
  \ROOutputAndTrace{\ROFunction}{\MTAdversary_{e}}{\ROTrace}{(\MTCommitment,\ROAdvState)} \\
  (\MTMessageVector,\MTTrapdoor) \gets \MTExtractor(\MTCommitment,\ROTrace) \\
  (\MTMessageIndexSet_{b},\MTMessageSubVector_{b},\MTProof_{b}) \gets \MTAdversary_{e}^{\ROFunction}(\ROAdvState) \\
  \MTProof \gets \MTOpen^{\ROFunction}(\MTTrapdoor, \MTMessageIndexSet_{b})
}
\right]
\enspace.
\end{equation*}
We conclude that any extraction error $\MTExtractionError$ that can be proved must satisfy the relation $\frac{1}{2} \cdot \MTBindingError \leq \MTExtractionError$ which gives us the desired bound on the binding error $\MTBindingError$.
\end{remark}

%%%%%%%%%%%%%%%%%%%%%%%%%%%%%%%%%%%%%%%%%%%%%%%%%%%%%%%%%%%%%%%%%%%%%%%%%%%%%%%
%%%%%%%%%%%%%%%%%%%%%%%%%%%%%%%%%%%%%%%%%%%%%%%%%%%%%%%%%%%%%%%%%%%%%%%%%%%%%%%
\subsection{Extraction for multiple commitments}
\label{section:merkle-commitment-multi-extractability}

We discuss how extractability extends to adversaries that output multiple Merkle commitments. Intuitively, this should follow directly from \Cref{lemma:mt-extractability} via a union bound: if the adversary outputs $\MTNumCommitments$ commitments and subsequently an opening for one of the commitments, then the probability that the opening is valid but the extractor did not succeed for the corresponding commitment is at most $\MTNumCommitments$ times the error in \Cref{lemma:mt-extractability}. This error can, however, be improved. Each execution of the Merkle commitment extractor is applied to a (possibly) different Merkle commitment, but with (possibly different prefixes of) the \emph{same} query-answer trace $\ROTrace$. This enables a tighter analysis.

The lemma below extends \Cref{lemma:mt-extractability} to the case where an adversary outputs $\MTNumCommitments$ Merkle commitments and the extractor is tasked to extract for the commitment that the adversary chooses to open (successfully). Similarly to \Cref{lemma:cm-multi-extractability} for the basic commitment scheme $\CMSymbol$, we let the extractor $\MTMultiExtractor$ be a stateful algorithm.

\begin{lemma}[$\MTSymbol$ is multi-extractable]
\label{lemma:mt-multi-extractability}
Let $\MTSymbol \DefineEqual \MTConstructor{\ROOutputSize}{\MTAlphabet}{\MTMessageLength}{\MTSaltSize}$ and define $\MTDepth \DefineEqual \log_{2} \MTMessageLength$. There exists a deterministic \textbf{stateful} algorithm $\MTMultiExtractor$ such that for every query bound $\ROQueryBound \in \Naturals$, $\ROQueryBound$-query algorithm $\MTAdversary$, and number of commitments $\MTNumCommitments \in \Naturals$,
\begin{align*}
&\Pr\left[
\begin{array}{l}
\left(
\begin{array}{l}
\MTCheck^{\ROFunction}(\MTCommitment_{\MTCommitmentIndex},\MTMessageIndexSet,\MTMessageSubVector,\MTProof)=1 \\
\land\,(\MTMessageVector_{\MTCommitmentIndex}[\MTMessageIndexSet] \neq \MTMessageSubVector \lor \MTProof' \neq \MTProof)
\end{array}
\right)
\\
\lor \\
\left(
\begin{array}{l}
\exists\,\MTCommitmentIndex_{1},\MTCommitmentIndex_{2} \in [\MTNumCommitments] : \\
\MTCommitment_{\MTCommitmentIndex_{1}} = \MTCommitment_{\MTCommitmentIndex_{2}} \\
\land (\MTMessageVector_{\MTCommitmentIndex_{1}},\MTTrapdoor_{\MTCommitmentIndex_{1}}) \neq (\MTMessageVector_{\MTCommitmentIndex_{2}},\MTTrapdoor_{\MTCommitmentIndex_{2}})
\end{array}
\right)
\end{array}
\GivenExperiment
\StateExperiment{
\ROFunction \gets \RODistribution{\ROOutputSize} \\
\ROAdvState_{0} \DefineEqual \bot \\
\text{For $\MTCommitmentIndex=1,\dots,\MTNumCommitments$:} \\
\IndentedBullet\ROInputOutputAndTrace{\ROFunction}{\MTAdversary}{\ROAdvState_{\MTCommitmentIndex-1}}{\ROTrace_{\MTCommitmentIndex}}{(\MTCommitment_{\MTCommitmentIndex},\ROAdvState_{\MTCommitmentIndex})} \\
\IndentedBullet(\MTMessageVector_{\MTCommitmentIndex},\MTTrapdoor_{\MTCommitmentIndex}) \gets \MTMultiExtractor(\MTCommitment_{\MTCommitmentIndex},\ROTrace_{\MTCommitmentIndex}) \\
(\MTCommitmentIndex \in [\MTNumCommitments],\MTMessageIndexSet,\MTMessageSubVector,\MTProof) \gets \MTAdversary^{\ROFunction}(\ROAdvState_{\MTNumCommitments}) \\
\MTProof' \gets \MTOpen^{\ROFunction}(\MTTrapdoor_{\MTCommitmentIndex}, \MTMessageIndexSet)
}
\right]
\\&\leq
\MTMultiExtractionError(\ROOutputSize,\MTMessageLength,\ROQueryBound,\MTNumCommitments)
\enspace.
\end{align*}
Above:
\begin{itemize}[nolistsep]
  \item $\MTMultiExtractionError(\ROOutputSize,\MTMessageLength,\ROQueryBound,\MTNumCommitments) \leq \MTMultiExtractabilityExpression{\ROOutputSize}{\ROQueryBound}{\MTMessageLength}{\MTDepth}{\MTNumCommitments}$ (at most $\MTMultiExtractabilityShortExpression{\ROOutputSize}{\ROQueryBound}{\MTNumCommitments}{\MTMessageLength}$ if $\MTMultiExtractabilityCondition{\ROQueryBound}{\MTDepth}{\MTMessageLength}$);
  \item the total running time of $\MTMultiExtractor$ across the $\MTNumCommitments$ invocations is
\begin{equation*}
\MTMultiExtractorTimeFunction{\ROOutputSize}{\MTAlphabet}{\MTMessageLength}{\MTSaltSize}{\ROQueryBound}{\MTNumCommitments} = O\big(\MTNumCommitments \cdot \ROQueryBound \cdot \MTMessageLength \cdot (\log\Cardinality{\MTAlphabet} + \MTSaltSize + \ROOutputSize)\big)
\enspace.
\end{equation*}
\end{itemize}
\end{lemma}

\begin{proof}
The stateful extractor $\MTMultiExtractor$ is similar to the (stateless) extractor $\MTExtractor$ in \Cref{construction:mt-extractor}. The extractor $\MTMultiExtractor$ internally maintains the concatenation of all query-answer traces received so far (after the $i$-th invocation it stores the trace $\ROTrace \DefineEqual \ROTrace_{1} \concat \cdots \concat \ROTrace_{i}$), and uses this trace to construct the message and trapdoor corresponding to the given commitment.
\begin{itemize}
\item[] $\MTMultiExtractor(\MTCommitment_{\MTCommitmentIndex},\ROTrace_{\MTCommitmentIndex})$:
\begin{enumerate}[nolistsep]
  \item Append $\ROTrace_{\MTCommitmentIndex}$ to the query-answer trace $\ROTrace$ stored in the internal state.
  \item Compute and output $(\MTMessageVector_{\MTCommitmentIndex},\MTTrapdoor_{\MTCommitmentIndex}) \DefineEqual \MTExtractor(\MTCommitment_{\MTCommitmentIndex},\ROTrace)$.
\end{enumerate}
\end{itemize}
By \Cref{lemma:mt-extractability} each invocation of $\MTExtractor$ runs in time at most $O\big( \ROQueryBound \cdot \MTMessageLength \cdot (\log\Cardinality{\MTAlphabet} + \MTSaltSize + \ROOutputSize)\big)$, because the stored query-answer trace contains at most $\ROQueryBound$ entries. The $\MTNumCommitments$ invocations yield the claimed total running time. The time complexity can be improved via a suitable use of dynamic data structures in the internal state (and foregoing $\MTExtractor$ as a subroutine); see \Cref{remark:mt-multi-extractor-efficient-implementation}.

Now we analyze the extraction error. Consider the following query-answer traces:
\begin{itemize}[noitemsep]
  \item $\ROTrace \DefineEqual \ROTrace_{1} \concat \cdots \concat \ROTrace_{\MTNumCommitments}$;
  \item $\ROOpenTrace$ is the query-answer trace of the execution $(\MTCommitmentIndex, \MTMessageIndexSet,\MTMessageSubVector,\MTProof) \gets \MTAdversary^{\ROFunction}(\ROAdvState_{\MTNumCommitments})$;
  \item $\ROHonestTrace$ is the query-answer trace of the execution $\MTCheck^{\ROFunction}(\MTCommitment_{\MTCommitmentIndex},\MTMessageIndexSet,\MTMessageSubVector,\MTProof)$.
\end{itemize}

Fix any $\ROQueryBound_{1},\ROQueryBound_{2} \in \Naturals$ with $\ROQueryBound_{1} + \ROQueryBound_{2} \leq \ROQueryBound$. We prove the claimed upper bound conditioned on the event that $\ROQueryBound_{1} = \Cardinality{\ROTrace}$ and $\ROQueryBound_{2} = \Cardinality{\ROOpenTrace}$. The lemma follows from this because the upper bound holds for every $\ROQueryBound_{1},\ROQueryBound_{2} \in \Naturals$ with $\ROQueryBound_{1} + \ROQueryBound_{2} \leq \ROQueryBound$, and every computation of the $\ROQueryBound$-query adversary $\MTAdversary$ implies a setting of $\ROQueryBound_{1},\ROQueryBound_{2} \in \Naturals$ with $\ROQueryBound_{1} + \ROQueryBound_{2} \leq \ROQueryBound$.

We define several events.
\begin{itemize}[noitemsep]
  \item $\Event$ is the event that the conditions in the probability statement hold.
  \item $\CollisionEvent$ is the event that $\ROTrace$ contains a collision.
  \item $\NotEqualEventA$ is the event that $\MTGraph_{\MTMessageLength}^{\MTCommitmentIndex} \neq \hat{\MTGraph}_{\MTMessageLength}^{\MTCommitmentIndex}$ where
  \begin{itemize}[nosep]
    \item $\MTGraph_{\MTMessageLength}^{\MTCommitmentIndex}$ is the tree generated during the execution of $\MTExtractor(\MTCommitment_{\MTCommitmentIndex},\ROTrace_{1} \concat \cdots \concat \ROTrace_{\MTCommitmentIndex})$, and
    \item $\hat{\MTGraph}_{\MTMessageLength}^{\MTCommitmentIndex}$ is the tree generated during the execution of $\MTExtractor(\MTCommitment_{\MTCommitmentIndex},\ROTrace \concat \ROOpenTrace)$.
  \end{itemize}
  \item $\NotEqualEventB$ is the event that there exist $\MTCommitmentIndex_{1}, \MTCommitmentIndex_{2} \in [\MTNumCommitments]$ with $\MTCommitment_{\MTCommitmentIndex_{1}} = \MTCommitment_{\MTCommitmentIndex_{2}}$ and $\MTGraph_{\MTMessageLength}^{\MTCommitmentIndex_{1}} \neq \MTGraph_{\MTMessageLength}^{\MTCommitmentIndex_{2}}$ where
  \begin{itemize}[nosep]
    \item $\MTGraph_{\MTMessageLength}^{\MTCommitmentIndex_{1}}$ is the tree generated during the execution of $\MTExtractor(\MTCommitment_{\MTCommitmentIndex_{1}},\ROTrace_{1} \concat \cdots \concat \ROTrace_{\MTCommitmentIndex_{1}})$, and
    \item $\MTGraph_{\MTMessageLength}^{\MTCommitmentIndex_{2}}$ is the tree generated during the execution of $\MTExtractor(\MTCommitment_{\MTCommitmentIndex_{2}},\ROTrace_{1} \concat \cdots \concat \ROTrace_{\MTCommitmentIndex_{2}})$.
  \end{itemize}
  \item $\NotEqualEvent$ is the event $\NotEqualEventA \vee \NotEqualEventB$.
  \item $\SubsetEvent$ is the event that $\ROHonestTrace \not \subseteq \ROTrace$ and $\MTCheck^{\ROFunction}(\MTCommitment_{\MTCommitmentIndex},\MTMessageIndexSet,\MTMessageSubVector,\MTProof)=1$.
\end{itemize}
The trees $(\MTGraph_{\MTMessageLength}^{\MTCommitmentIndex})_{\MTCommitmentIndex \in [\MTNumCommitments]}$ defined in $\NotEqualEvent$ are the trees generated during the execution of   $\MTMultiExtractor$: for every $\MTCommitmentIndex \in [\MTNumCommitments]$, $\MTMultiExtractor(\MTCommitment_{\MTCommitmentIndex},\ROTrace_{\MTCommitmentIndex})$ corresponds to $\MTExtractor(\MTCommitment_{\MTCommitmentIndex},\ROTrace_{1} \concat \cdots \concat \ROTrace_{\MTCommitmentIndex})$.

Our goal is to upper bound the probability that $\Event$ holds. We break the analysis into several cases:
\begin{align*}
\Pr[\Event] &
\leq
\Pr[\CollisionEvent] + \Pr\left[\NotEqualEventA \ConditionedOn \Negate{\CollisionEvent}\right]
+ \Pr\left[\NotEqualEventB \ConditionedOn \Negate{\CollisionEvent}\right]
\\ &\quad
+ \Pr\left[\SubsetEvent \ConditionedOn \Negate{\CollisionEvent} \land \Negate{\NotEqualEvent} \right]
+ \Pr\left[\Event \ConditionedOn \Negate{\CollisionEvent} \land \Negate{\NotEqualEvent} \land \Negate{\SubsetEvent} \right]
\enspace.
\end{align*}
We upper bound each term individually, yielding the claimed upper bound.
\begin{enumerate}

  \item An upper bound on the probability of the event $\CollisionEvent$ follows directly from \Cref{lemma:rom-cr}:
\begin{equation*}
\Pr[\CollisionEvent]
\leq
\ROCRExpression{\ROOutputSize}{\ROQueryBound_{1}}
\enspace.
\end{equation*}

  \item We upper bound the probability of the event $\left.\NotEqualEventA \ConditionedOn \Negate{\CollisionEvent}\right.$.

Since $\NotEqualEventA$ holds ($\MTGraph_{\MTMessageLength}^{\MTCommitmentIndex} \neq \hat{\MTGraph}_{\MTMessageLength}^{\MTCommitmentIndex}$), $\ROTrace_{\MTCommitmentIndex+1} \concat \cdots \concat \ROTrace_{\MTNumCommitments} \concat \ROOpenTrace$ contains a query that is not in $\ROTrace_{1} \concat \cdots \concat \ROTrace_{\MTCommitmentIndex}$ with an answer that equals a non-dummy label in $\MTGraph_{\MTMessageLength}^{\MTCommitmentIndex}$.

Since the event $\CollisionEvent$ does not hold ($\ROTrace$ contains no collisions), the aforementioned query cannot be in $\ROTrace_{\MTCommitmentIndex+1} \concat \cdots \concat \ROTrace_{\MTNumCommitments}$ because that would form a collision in $\ROTrace$. Hence it suffices to upper bound the probability that there exists a query in $\ROOpenTrace$ that is not in $\ROTrace_{1} \concat \cdots \concat \ROTrace_{\MTCommitmentIndex}$ with an answer that equals a non-dummy label in $\MTGraph_{\MTMessageLength}^{\MTCommitmentIndex}$.

Since the event $\CollisionEvent$ does not hold ($\ROTrace$ contains no collisions), each vertex in any of the trees $(\MTGraph_{\MTMessageLength}^{\MTCommitmentIndex})_{\MTCommitmentIndex \in [\MTNumCommitments]}$ is labeled at most once. The number of non-dummy labels in $\MTGraph_{\MTMessageLength}^{\MTCommitmentIndex}$ is upper bounded by $\min\{3\ROQueryBound_{1},2\MTMessageLength\}$. The upper bound $3\ROQueryBound_{1}$ is since trees are constructed from a trace of length at most $\ROQueryBound_{1}$ and each query adds at most three labels in one tree (most queries add two labels but the first one in a tree adds three since the root is added as a label as well). The upper bound $2\MTMessageLength$ is since a tree has at most $2\MTMessageLength$ labels. For every query in $\ROOpenTrace$, the probability that the query answer equals one of the these non-dummy labels is at most $\frac{\min\{3\ROQueryBound_{1},2\MTMessageLength\}}{2^{\ROOutputSize}}$.

Taking a union bound over all queries in $\ROOpenTrace$,
\begin{align*}
\Pr\left[\NotEqualEventA \ConditionedOn \Negate{\CollisionEvent}\right]
\leq
\Cardinality{\ROOpenTrace} \cdot \frac{\min\{3\ROQueryBound_{1},2\MTMessageLength\}}{2^{\ROOutputSize}}
=
\ROQueryBound_{2} \cdot \frac{\min\{3\ROQueryBound_{1},2\MTMessageLength\}}{2^{\ROOutputSize}}
\enspace.
\end{align*}

  \item We upper bound the probability of the event $\left.\NotEqualEventB \ConditionedOn \Negate{\CollisionEvent}\right.$.

Since the event $\CollisionEvent$ does not hold ($\ROTrace$ contains no collisions), each vertex in any of the trees $(\MTGraph_{\MTMessageLength}^{\MTCommitmentIndex})_{\MTCommitmentIndex \in [\MTNumCommitments]}$ is labeled at most once.

For every $\MTCommitmentIndex_{1},\MTCommitmentIndex_{2} \in [\MTNumCommitments]$ with $\MTCommitmentIndex_{1} < \MTCommitmentIndex_{2}$ and $\MTCommitment_{\MTCommitmentIndex_{1}} = \MTCommitment_{\MTCommitmentIndex_{2}}$, if $\MTGraph_{\MTMessageLength}^{\MTCommitmentIndex_{1}} \neq \MTGraph_{\MTMessageLength}^{\MTCommitmentIndex_{2}}$ then there must be a query in $\ROTrace_{1} \concat \cdots \concat \ROTrace_{\MTCommitmentIndex_{2}}$ that is not in $\ROTrace_{1} \concat \cdots \concat \ROTrace_{\MTCommitmentIndex_{1}}$ with an answer that equals a non-dummy label in $\MTGraph_{\MTMessageLength}^{\MTCommitmentIndex_{1}}$. Below we upper bound the probability that such a pair of indexes $\MTCommitmentIndex_{1},\MTCommitmentIndex_{2} \in [\MTNumCommitments]$ exist.

Suppose that, for $j \in [\ROQueryBound_{1}]$, the adversary's $j$-th query happens in iteration $\MTCommitmentIndex_{2}$ of the loop (i.e., at the end of this iteration the adversary outputs $\MTCommitment_{\MTCommitmentIndex_{2}}$). The number of non-dummy labels in the first $\MTCommitmentIndex_{2}-1$ trees $\MTGraph_{\MTMessageLength}^{1},\dots,\MTGraph_{\MTMessageLength}^{\MTCommitmentIndex_{2}-1}$ is bounded in two ways.
\begin{itemize}[nolistsep]
\item Each query-answer pair leads to $\leq 2$ new labels in a single tree plus at most $\MTCommitmentIndex_{2}-1$ root labels. This is at most $2(j-1) + \MTCommitmentIndex_{2}-1 \leq 2(j-1) + \MTNumCommitments-1$ non-dummy labels.
\item Each tree has at most $2\MTMessageLength$ labels and there are $\MTCommitmentIndex_{2}-1 \leq \MTNumCommitments-1$ trees. This is at most $(\MTNumCommitments-1) \cdot 2\MTMessageLength$ non-dummy labels.
\end{itemize}
Therefore the number of non-dummy labels is at most
\begin{equation*}
\min\{2(j-1) + \MTNumCommitments-1,(\MTNumCommitments-1) \cdot 2\MTMessageLength\}
\enspace.
\end{equation*}
Hence the probability that the answer of the $j$-th query matches any of the non-dummy labels is at most $\frac{\min\{2(j-1) + \MTNumCommitments-1,(\MTNumCommitments-1) \cdot 2\MTMessageLength\}}{2^{\ROOutputSize}}$.

Summing over all $\ROQueryBound_{1}$ queries in $\ROTrace$,
\begin{align*}
\Pr\left[\NotEqualEventB \ConditionedOn \Negate{\CollisionEvent}\right]
\leq
\sum_{j=1}^{\ROQueryBound_{1}} \frac{\min\{2(j-1) + \MTNumCommitments-1,(\MTNumCommitments-1) \cdot 2\MTMessageLength\}}{2^{\ROOutputSize}}
\leq \ROQueryBound_{1} \cdot \frac{\min\{\ROQueryBound_{1} + \MTNumCommitments-2,(\MTNumCommitments-1) \cdot 2\MTMessageLength\} }{2^{\ROOutputSize}}
\enspace.
\end{align*}

  \item We upper bound the probability of the event $\left.\SubsetEvent \ConditionedOn \Negate{\CollisionEvent} \land \Negate{\NotEqualEvent}  \right.$.

For every $\MTMessageIndexAlt \in \MTMessageIndexSet$, we define the subset $\ROHonestTrace_{\MTMessageIndexAlt} \subseteq \ROHonestTrace$ be the query-answer trace of the execution $\MTCheck^{\ROFunction}(\MTCommitment_{\MTCommitmentIndex},\MTMessageIndexAlt,\MTMessageSubVector[\MTMessageIndexAlt],\MTAuthPath_{\MTMessageIndexAlt})$, where $\MTAuthPath_{\MTMessageIndexAlt}$ is the authentication path for location $\MTMessageIndexAlt$ in $\MTProof$.

Suppose that $\SubsetEvent$ holds. Then $\ROHonestTrace \not \subseteq \ROTrace$, so there exists $\MTMessageIndexAlt \in \MTMessageIndexSet$ such that $\ROHonestTrace_{\MTMessageIndexAlt} \not \subseteq \ROTrace$. Moreover, $\MTCheck^{\ROFunction}(\MTCommitment,\MTMessageIndex,\MTMessageSubVector[\MTMessageIndex],\MTAuthPath_{\MTMessageIndex})=1$, so the query-answer trace $\ROHonestTrace_{\MTMessageIndexAlt}$ can be viewed as a path that \DoQuote{ends} in the Merkle commitment $\MTCommitment_{\MTCommitmentIndex}$, and thus \DoQuote{connects} somewhere into $\MTGraph_{\MTMessageLength}^{\MTCommitmentIndex}$. In sum, if $\SubsetEvent$ holds then there exists a query in $\ROHonestTrace_{\MTMessageIndexAlt}$ that is not in $\ROTrace$ with an answer that equals a non-dummy label in $\MTGraph_{\MTMessageLength}^{\MTCommitmentIndex}$.

Since the event $\CollisionEvent$ does not hold ($\ROTrace$ contains no collisions), each vertex in any of the trees $(\MTGraph_{\MTMessageLength}^{\MTCommitmentIndex})_{\MTCommitmentIndex \in [\MTNumCommitments]}$ is labeled at most once. The number of non-dummy labels in $\MTGraph_{\MTMessageLength}^{\MTCommitmentIndex}$ is upper bounded by $\min\{3\ROQueryBound_{1},2\MTMessageLength\}$. The upper bound $3\ROQueryBound_{1}$ is since the trees are constructed from a trace of length at most $\ROQueryBound_{1}$ and each query can add at most three labels in one tree (most queries add two labels but the first one in a tree adds three since the root is added as a label as well). The upper bound $2\MTMessageLength$ is since each tree has at most $2\MTMessageLength$ labels. For every query in $\ROHonestTrace_{\MTMessageIndexAlt}$ (that is not in $\ROTrace$), the probability that the query answer equals one of the these non-dummy labels is at most $\frac{\min\{3\ROQueryBound_{1},2\MTMessageLength\}}{2^{\ROOutputSize}}$.

Since $\Negate{\NotEqualEvent}$ (and thus $\Negate{\NotEqualEventA}$), there are no queries in $\ROTrace'$ with an answer that equals a non-dummy label in $\MTGraph_{\MTMessageLength}^{\MTCommitmentIndex}$. Recall that $\ROHonestTrace_{\MTMessageIndexAlt}$ has length $\MTDepth+1$. Taking a union bound over all queries in $\ROHonestTrace_{\MTMessageIndexAlt}$, the probability that there exists a query in $\ROHonestTrace_{\MTMessageIndexAlt}$ (that is not in $\ROTrace \cup \ROOpenTrace$) with an answer that equals a non-dummy label in $\MTGraph_{\MTMessageLength}^{\MTCommitmentIndex}$ is at most $(\MTDepth+1) \cdot \frac{\min\{3\ROQueryBound_{1},2\MTMessageLength\}}{2^{\ROOutputSize}}$.

Hence
\begin{align*}
\Pr\left[\SubsetEvent \ConditionedOn \Negate{\CollisionEvent} \land \Negate{\NotEqualEvent} \right]
\leq
(\MTDepth+1) \cdot \frac{\min\{3\ROQueryBound_{1},2\MTMessageLength\}}{2^{\ROOutputSize}}
\enspace.
\end{align*}

  \item We show that the event $\left.\Event \ConditionedOn \Negate{\CollisionEvent} \land \Negate{\NotEqualEvent} \land \Negate{\SubsetEvent}\right.$ cannot happen.

First we argue that for every $\MTCommitmentIndex_{1},\MTCommitmentIndex_{2} \in [\MTNumCommitments]$ with $\MTCommitment_{\MTCommitmentIndex_{1}} = \MTCommitment_{\MTCommitmentIndex_{2}}$ it holds that $(\MTMessageVector_{\MTCommitmentIndex_{1}},\MTTrapdoor_{\MTCommitmentIndex_{1}}) = (\MTMessageVector_{\MTCommitmentIndex_{2}},\MTTrapdoor_{\MTCommitmentIndex_{2}})$. This follows directly from the fact that $\Negate{\NotEqualEventB}$ implies that $\MTGraph_{\MTMessageLength}^{\MTCommitmentIndex_{1}} = \MTGraph_{\MTMessageLength}^{\MTCommitmentIndex_{2}}$, so the extracted message vectors and trapdoor information are the same, namely, $(\MTMessageVector_{\MTCommitmentIndex_{1}},\MTTrapdoor_{\MTCommitmentIndex_{1}}) = (\MTMessageVector_{\MTCommitmentIndex_{2}},\MTTrapdoor_{\MTCommitmentIndex_{2}})$.

Next we argue that if $\MTCheck^{\ROFunction}(\MTCommitment_{\MTCommitmentIndex},\MTMessageIndexSet,\MTMessageSubVector,\MTProof)=1$ then $\MTMessageVector_{\MTCommitmentIndex}[\MTMessageIndexSet] = \MTMessageSubVector$ and $\MTProof' = \MTProof$. This follows from applying \Cref{lemma:mt-extractability-same-trees} for all commitments $(\MTCommitment_{\MTCommitmentIndex})_{\MTCommitmentIndex \in [\MTNumCommitments]}$, as we now explain. Fix any $\MTCommitmentIndex \in [\MTNumCommitments]$, and consider the trace up to the point where the adversary outputs $\MTCommitment_{\MTCommitmentIndex}$ (i.e., the trace $\ROTrace_{1} \concat \cdots \concat \ROTrace_{\MTCommitmentIndex}$), and the trace of all queries performed after (i.e., the trace $\ROTrace_{\MTCommitmentIndex+1} \concat \cdots \concat \ROTrace_{\MTNumCommitments} \concat \ROOpenTrace)$. Since we assume $\Negate{\CollisionEvent} \land \Negate{\NotEqualEvent} \land \Negate{\SubsetEvent}$, both conditions in the probability statement of \Cref{lemma:mt-extractability-same-trees} hold as well. This implies that if $\MTCheck^{\ROFunction}(\MTCommitment_{\MTCommitmentIndex},\MTMessageIndexSet,\MTMessageSubVector,\MTProof)=1$ then $\MTMessageVector_{\MTCommitmentIndex}[\MTMessageIndexSet] = \MTMessageSubVector$ and $\MTProof' = \MTProof$ as required.
\end{enumerate}
We conclude that
\begin{align*}
&\Pr[\Event]
\\ & \leq
\Pr[\CollisionEvent]
+ \Pr\left[\NotEqualEventA \ConditionedOn \Negate{\CollisionEvent}\right]
+ \Pr\left[\NotEqualEventB \ConditionedOn \Negate{\CollisionEvent}\right]
+ \Pr\left[\SubsetEvent \ConditionedOn \Negate{\CollisionEvent} \land \Negate{\NotEqualEvent} \right]
\\ & +
\Pr\left[\Event \ConditionedOn \Negate{\CollisionEvent} \land \Negate{\NotEqualEvent} \land \Negate{\SubsetEvent}\right]
\\ & \leq
\ROCRExpression{\ROOutputSize}{\ROQueryBound_{1}}
+ \frac{\ROQueryBound_{2} \cdot \min\{3\ROQueryBound_{1},2\MTMessageLength\}}{2^{\ROOutputSize}}
+ \ROQueryBound_{1} \cdot \frac{\min\{\ROQueryBound_{1} + \MTNumCommitments-2,(\MTNumCommitments-1) \cdot 2\MTMessageLength\} }{2^{\ROOutputSize}}
+ (\MTDepth+1) \cdot \frac{\min\{3\ROQueryBound_{1},2\MTMessageLength\}}{2^{\ROOutputSize}}
+ 0
\\ & \leq
\ROCRExpression{\ROOutputSize}{\ROQueryBound_{1}}
+ \frac{(\ROQueryBound - \ROQueryBound_{1} + \MTDepth+1) \cdot \min\{3\ROQueryBound_{1},2\MTMessageLength\}}{2^{\ROOutputSize}}
 + \ROQueryBound_{1} \cdot \frac{\min\{\ROQueryBound_{1} + \MTNumCommitments-2,(\MTNumCommitments-1) \cdot 2\MTMessageLength\} }{2^{\ROOutputSize}}
\EquationComment{since $\ROQueryBound_{1} + \ROQueryBound_{2} \leq \ROQueryBound$}
\\ & \leq
\ROCRExpression{\ROOutputSize}{\ROQueryBound_{1}}
+ \frac{(\ROQueryBound - \ROQueryBound_{1} + \MTDepth+1) \cdot 2\MTMessageLength}{2^{\ROOutputSize}}
+ \frac{\ROQueryBound_{1} \cdot (\ROQueryBound_{1} + \MTNumCommitments-2)}{2^{\ROOutputSize}}
\EquationComment{since $\min\{a,b\}\leq a,b$}
\\ & =
\frac{3}{2} \cdot \frac{(\ROQueryBound_{1}-1) \cdot \ROQueryBound_{1}}{2^{\ROOutputSize}}
+ \frac{(\ROQueryBound - \ROQueryBound_{1} + \MTDepth+1) \cdot 2\MTMessageLength}{2^{\ROOutputSize}}
+ \frac{\ROQueryBound_{1} \cdot (\MTNumCommitments-1)}{2^{\ROOutputSize}}
\enspace.
\end{align*}
The expression is a convex function in $\ROQueryBound_{1}$ so its maximum on the interval $\ROQueryBound_{1} \in \{0,1,\dots,\ROQueryBound\}$ is achieved at either $\ROQueryBound_{1}=0$ or $\ROQueryBound_{1} = \ROQueryBound$. Hence the expression is upper bounded from the maximum of these two options (where the second term is no smaller than the first):
\begin{align*}
\Pr[\Event]
& \leq
 \max\left\{
\frac{(\ROQueryBound + \MTDepth+1) \cdot 2\MTMessageLength}{2^{\ROOutputSize}}
,
\MTMultiExtractabilityExpression{\ROOutputSize}{\ROQueryBound}{\MTMessageLength}{\MTDepth}{\MTNumCommitments}
\right\}
\\ & \leq
\MTMultiExtractabilityExpression{\ROOutputSize}{\ROQueryBound}{\MTMessageLength}{\MTDepth}{\MTNumCommitments}
\enspace,
\end{align*}
where the second inequality holds under the mild condition that $3\ROQueryBound \geq 4\MTMessageLength - 2\MTNumCommitments + 3$.
\end{proof}

\begin{remark}[time complexity of $\MTMultiExtractor$]
\label{remark:mt-multi-extractor-efficient-implementation}
A straightforward realization of procedure $\MTMultiExtractor$ leads to the time complexity claimed in \Cref{lemma:mt-multi-extractability}, because $\MTExtractor$ is run at most $\MTNumCommitments$ times.

If $\log \ROQueryBound = O(\MTNumCommitments \cdot \MTMessageLength)$ (a natural regime) then we can do better by using suitable data structures. The idea is similar to the improved realization of $\MTExtractor$ described in \Cref{remark:mt-extractor-time}. The main difference is that $\MTMultiExtractor$ receives, with each invocation, an additional query-answer trace. We use a dynamic dictionary (instead of a static dictionary), for which inserts and searches take time $\log \ROQueryBound \cdot \MTTraceEntrySize$ (when there are $\ROQueryBound$ entries each consisting of $\MTTraceEntrySize$ bits); these are the same costs as for a static dictionary. The total number of inserts into the dictionary is $\ROQueryBound$, and the total number of searches is $O(\MTNumCommitments \cdot \MTMessageLength)$. This leads to a time complexity that is $O((\ROQueryBound + \MTNumCommitments \cdot \MTMessageLength) \cdot \log \ROQueryBound \cdot \MTTraceEntrySize)$ where $\MTTraceEntrySize \DefineEqual \log\Cardinality{\MTAlphabet} + \MTSaltSize + \ROOutputSize$. This expression improves on the one in \Cref{lemma:mt-multi-extractability} provided that $\log \ROQueryBound = O(\MTNumCommitments \cdot \MTMessageLength)$.
\end{remark}

\begin{remark}[multiple openings]
\label{remark:mt-multi-extractability-with-multiple-openings}
\Cref{lemma:mt-multi-extractability} considers the case where the adversary outputs one opening, which suffices for the purposes of this book. More generally, one can consider the case where the adversary outputs $\MTNumOpenings \in \Naturals$ openings. The proof of \Cref{lemma:mt-multi-extractability} can be extended to this case, establishing an upper bound $\MTMultiExtractionError(\ROOutputSize,\MTMessageLength,\ROQueryBound,\MTNumCommitments,\MTNumOpenings)$ for the following probability:
\begin{equation*}
\Pr\left[
\begin{array}{l}
\left(
\begin{array}{l}
\exists\,\MTOpeningIndex \in [\MTNumOpenings] : \\
\MTCheck^{\ROFunction}(\MTCommitment_{\MTCommitmentIndex_{\MTOpeningIndex}},\MTMessageIndexSet_{\MTOpeningIndex},\MTMessageSubVector_{\MTOpeningIndex},\MTProof_{\MTOpeningIndex})=1 \\
\land\,(\MTMessageVector_{\MTCommitmentIndex_{\MTOpeningIndex}}[\MTMessageIndexSet_{\MTOpeningIndex}] \neq \MTMessageSubVector_{\MTOpeningIndex} \lor \MTProof'_{\MTOpeningIndex} \neq \MTProof_{\MTOpeningIndex})
\end{array}
\right)
\\
\lor \\
\left(
\begin{array}{l}
\exists\,\MTCommitmentIndex_{1},\MTCommitmentIndex_{2} \in [\MTNumCommitments] : \\
\MTCommitment_{\MTCommitmentIndex_{1}}
=
\MTCommitment_{\MTCommitmentIndex_{2}} \\
\land
(\MTMessageVector_{\MTCommitmentIndex_{1}},\MTTrapdoor_{\MTCommitmentIndex_{1}})
\neq
(\MTMessageVector_{\MTCommitmentIndex_{2}},\MTTrapdoor_{\MTCommitmentIndex_{2}})
\end{array}
\right)
\end{array}
\GivenExperiment
\StateExperiment{
\ROFunction \gets \RODistribution{\ROOutputSize} \\
\ROAdvState_{0} \DefineEqual \bot \\
\text{For $\MTCommitmentIndex=1,\dots,\MTNumCommitments$:} \\
\IndentedBullet\ROInputOutputAndTrace{\ROFunction}{\MTAdversary}{\ROAdvState_{\MTCommitmentIndex-1}}{\ROTrace_{\MTCommitmentIndex}}{(\MTCommitment_{\MTCommitmentIndex},\ROAdvState_{\MTCommitmentIndex})} \\
\IndentedBullet(\MTMessageVector_{\MTCommitmentIndex},\MTTrapdoor_{\MTCommitmentIndex}) \gets \MTMultiExtractor(\MTCommitment_{\MTCommitmentIndex},\ROTrace_{\MTCommitmentIndex}) \\
\big((\MTCommitmentIndex_{\MTOpeningIndex} \in [\MTNumCommitments], \MTMessageIndexSet_{\MTOpeningIndex},\MTMessageSubVector_{\MTOpeningIndex},\MTProof_{\MTOpeningIndex})\big)_{\MTOpeningIndex \in [\MTNumOpenings]} \gets \MTAdversary^{\ROFunction}(\ROAdvState_{\MTNumCommitments}) \\
\text{For $\MTOpeningIndex \in [\MTNumOpenings]$: } \MTProof'_{\MTOpeningIndex} \gets \MTOpen^{\ROFunction}(\MTTrapdoor_{\MTCommitmentIndex_{\MTOpeningIndex}}, \MTMessageIndexSet_{\MTOpeningIndex})
}
\right]
\enspace.
\end{equation*}
Specifically, one can show the upper bound
\begin{equation*}
\MTMultiExtractionError(\ROOutputSize,\MTMessageLength,\ROQueryBound,\MTNumCommitments,\MTNumOpenings)
\leq
\MTMultiExtractabilityMultiOpeningExpression{\ROOutputSize}{\ROQueryBound}{\MTMessageLength}{\MTDepth}{\MTNumCommitments}{\MTNumOpenings}
\enspace.
\end{equation*}
The analysis is a straightforward extension of the proof of \Cref{lemma:mt-multi-extractability}.

For every $\MTOpeningIndex \in [\MTNumOpenings]$, let $\ROHonestTrace_{\MTOpeningIndex}$ be the query-answer trace of the execution $\MTCheck^{\ROFunction}(\MTCommitment_{\MTCommitmentIndex_{\MTOpeningIndex}},\MTMessageIndexSet_{\MTOpeningIndex},\MTMessageSubVector_{\MTOpeningIndex},\MTProof_{\MTOpeningIndex})$. The event $\Event$ is defined to be the new condition in the probability statement above; the events $\CollisionEvent,\NotEqualEventA,\NotEqualEventB,\NotEqualEvent$ are defined as before; and $\SubsetEvent$ is extended to be the event that there exists $\MTOpeningIndex \in [\MTNumOpenings]$ with $\ROHonestTrace_{\MTOpeningIndex} \not \subseteq \ROTrace$ and $\MTCheck^{\ROFunction}(\MTCommitment_{\MTCommitmentIndex_{\MTOpeningIndex}},\MTMessageIndexSet_{\MTOpeningIndex},\MTMessageSubVector_{\MTOpeningIndex},\MTProof_{\MTOpeningIndex})=1$. The probability that $\Event$ holds is broken into the same cases as in the proof of \Cref{lemma:mt-multi-extractability}, and all but one of the terms are upper bounded in the same way.

The only difference is the upper bound for the event $\left.\SubsetEvent \ConditionedOn \Negate{\CollisionEvent} \land \Negate{\NotEqualEvent}  \right.$ (due to the new definition of $\SubsetEvent$). Via an additional union bound on the $\MTNumOpenings$ openings, one can show that
\begin{align*}
\Pr\left[\SubsetEvent \ConditionedOn \Negate{\CollisionEvent} \land \Negate{\NotEqualEvent} \right]
\leq
\MTNumOpenings \cdot (\MTDepth+1) \cdot \frac{\min\{3\ROQueryBound_{1},2\MTMessageLength\}}{2^{\ROOutputSize}}
\enspace.
\end{align*}
This different upper bound leads to the bound for $\MTMultiExtractionError(\ROOutputSize,\MTMessageLength,\ROQueryBound,\MTNumCommitments,\MTNumOpenings)$ claimed above:
\begin{align*}
&\Pr[\Event]
\\ & \leq
\Pr[\CollisionEvent]
+ \Pr\left[\NotEqualEventA \ConditionedOn \Negate{\CollisionEvent}\right]
+ \Pr\left[\NotEqualEventB \ConditionedOn \Negate{\CollisionEvent}\right]
+ \Pr\left[\SubsetEvent \ConditionedOn \Negate{\CollisionEvent} \land \Negate{\NotEqualEvent} \right]
\\ & +
\Pr\left[\Event \ConditionedOn \Negate{\CollisionEvent} \land \Negate{\NotEqualEvent} \land \Negate{\SubsetEvent}\right]
\\ & \leq
\ROCRExpression{\ROOutputSize}{\ROQueryBound_{1}}
+ \frac{\ROQueryBound_{2} \cdot \min\{3\ROQueryBound_{1},2\MTMessageLength\}}{2^{\ROOutputSize}}
+ \ROQueryBound_{1} \cdot \frac{\min\{\ROQueryBound_{1} + \MTNumCommitments - 2,(\MTNumCommitments-1) \cdot 2\MTMessageLength\} }{2^{\ROOutputSize}}
+ \MTNumOpenings \cdot (\MTDepth+1) \cdot \frac{\min\{3\ROQueryBound_{1},2\MTMessageLength\}}{2^{\ROOutputSize}}
+ 0
\\ & \leq
\frac{1}{2} \cdot \frac{(\ROQueryBound_{1}-1)\cdot\ROQueryBound_{1}}{2^{\ROOutputSize}}
+ \frac{(\ROQueryBound - \ROQueryBound_{1} + \MTNumOpenings \cdot (\MTDepth+1)) \cdot \min\{3\ROQueryBound_{1},2\MTMessageLength\}}{2^{\ROOutputSize}}
+ \ROQueryBound_{1} \cdot \frac{\min\{\ROQueryBound_{1} + \MTNumCommitments - 2,(\MTNumCommitments-1) \cdot 2\MTMessageLength\} }{2^{\ROOutputSize}}
\EquationComment{since $\ROQueryBound_{1} + \ROQueryBound_{2} \leq \ROQueryBound$}
\\ & \leq
\frac{1}{2} \cdot \frac{(\ROQueryBound_{1}-1)\cdot\ROQueryBound_{1}}{2^{\ROOutputSize}}
+ \frac{(\ROQueryBound - \ROQueryBound_{1} + \MTNumOpenings \cdot (\MTDepth+1)) \cdot 2\MTMessageLength}{2^{\ROOutputSize}}
+ \frac{\ROQueryBound_{1} \cdot (\ROQueryBound_{1} + \MTNumCommitments-2)}{2^{\ROOutputSize}}
\EquationComment{since $\min\{a,b\}\leq a,b$}
\\ & \leq
\max\left\{
\frac{(\ROQueryBound + \MTNumOpenings \cdot (\MTDepth+1) ) \cdot 2\MTMessageLength}{2^{\ROOutputSize}}
,
\MTMultiExtractabilityMultiOpeningExpression{\ROOutputSize}{\ROQueryBound}{\MTMessageLength}{\MTDepth}{\MTNumCommitments}{\MTNumOpenings}
\right\}
\EquationComment{convexity in $\ROQueryBound_{1}$}
\\ & \leq
\MTMultiExtractabilityMultiOpeningExpression{\ROOutputSize}{\ROQueryBound}{\MTMessageLength}{\MTDepth}{\MTNumCommitments}{\MTNumOpenings}
\EquationComment{provided $3\ROQueryBound \geq 4\MTMessageLength - 2\MTNumCommitments + 3$}
\enspace.
\end{align*}
\end{remark}

%%%%%%%%%%%%%%%%%%%%%%%%%%%%%%%%%%%%%%%%%%%%%%%%%%%%%%%%%%%%%%%%%%%%%%%%%%%%%%%
%%%%%%%%%%%%%%%%%%%%%%%%%%%%%%%%%%%%%%%%%%%%%%%%%%%%%%%%%%%%%%%%%%%%%%%%%%%%%%%
\subsection{Extraction for multiple configurations and commitments}
\label{section:merkle-commitment-multi-configuration-multi-extractability}

We discuss how extractability further extends to adversaries that output multiple Merkle commitments across multiple configurations. We consider a setting with multiple configurations $(\MTSymbol_{\MTConfigIndex} \DefineEqual \MTConstructor{\ROOutputSize}{\MTAlphabet_{\MTConfigIndex}}{\MTMessageLength_{\MTConfigIndex}}{\MTSaltSize})_{\MTConfigIndex \in [\MTNumConfigurations]}$ of the Merkle commitment scheme each using a corresponding random oracle $(\ROFunction_{\MTConfigIndex})_{\MTConfigIndex \in [\MTNumConfigurations]} \in \RODistribution{\ROOutputSize}^{\MTNumConfigurations}$. The adversary outputs at most $\MTNumCommitments$ commitments for each of the $\MTNumConfigurations$ configurations; and subsequently the adversary outputs a single opening per configuration. The special case $\MTNumConfigurations=1$ corresponds to the setting in \Cref{lemma:mt-multi-extractability}.


\begin{lemma}[$\MTSymbol$ is multi-extractable across configurations]
\label{lemma:mt-multi-configuration-multi-extractability}
Let $(\MTSymbol_{\MTConfigIndex} \DefineEqual \MTConstructor{\ROOutputSize}{\MTAlphabet_{\MTConfigIndex}}{\MTMessageLength_{\MTConfigIndex}}{\MTSaltSize})_{\MTConfigIndex \in [\MTNumConfigurations]}$ and define $(\MTDepth_{\MTConfigIndex} \DefineEqual \log_{2} \MTMessageLength_{\MTConfigIndex})_{\MTConfigIndex \in [\MTNumConfigurations]}$. There exists a deterministic \textbf{stateful} algorithm $\MTMultiMultiExtractor{\MTNumConfigurations}$ such that for every query bound $\ROQueryBound \in \Naturals$, $\ROQueryBound$-query algorithm $\MTAdversary$, and number of commitments per configuration $\MTNumCommitments \in \Naturals$,
\begin{align*}
&\Pr\left[
\begin{array}{l}
\Cardinality{\MTConfigurationIdxSet{1}},\dots,\Cardinality{\MTConfigurationIdxSet{\MTNumConfigurations}} \leq \MTNumCommitments \\
\text{and the following disjunction holds:} \\
\left(
\begin{array}{l}
\exists\,\MTOpeningIndex \in [\MTNumConfigurations] : \\
\MTIdxCheck{\MTOpeningIndex}^{\ROFunction_{\MTOpeningIndex}}(\MTCommitment_{\MTCommitmentIndex_{\MTOpeningIndex}},\MTMessageIndexSet_{\MTOpeningIndex},\MTMessageSubVector_{\MTOpeningIndex},\MTProof_{\MTOpeningIndex})=1 \\
\land\,(\MTMessageVector_{\MTCommitmentIndex_{\MTOpeningIndex}}[\MTMessageIndexSet_{\MTOpeningIndex}] \neq \MTMessageSubVector_{\MTOpeningIndex} \lor \MTProof'_{\MTOpeningIndex} \neq \MTProof_{\MTOpeningIndex})
\end{array}
\right)
\\
\lor \\
\left(
\begin{array}{l}
\exists\,\MTOpeningIndex \in [\MTNumConfigurations]\;
\exists\,\MTCommitmentIndex_{1},\MTCommitmentIndex_{2} \in \MTConfigurationIdxSet{\MTOpeningIndex} : \\
\MTCommitment_{\MTCommitmentIndex_{1}}
=
\MTCommitment_{\MTCommitmentIndex_{2}} \\
\land
(\MTMessageVector_{\MTCommitmentIndex_{1}},\MTTrapdoor_{\MTCommitmentIndex_{1}})
\neq
(\MTMessageVector_{\MTCommitmentIndex_{2}},\MTTrapdoor_{\MTCommitmentIndex_{2}})
\end{array}
\right)
\end{array}
\GivenExperiment
\StateExperiment{
\OracleFunction=(\ROFunction_{\MTConfigIndex})_{\MTConfigIndex \in [\MTNumConfigurations]} \gets \RODistribution{\ROOutputSize}^{\MTNumConfigurations} \\
\ROAdvState_{0} \DefineEqual \bot \\
\MTConfigurationIdxSet{1},\dots,\MTConfigurationIdxSet{\MTNumConfigurations} \DefineEqual \emptyset \\
\text{For $\MTCommitmentIndex = 1,2,\dots$:} \\
\IndentedBullet\ROInputOutputAndTrace{\OracleFunction}{\MTAdversary}{\ROAdvState_{\MTCommitmentIndex-1}}{\ROTrace_{\MTCommitmentIndex}}{(\MTConfigIndex_{\MTCommitmentIndex} \in [\MTNumConfigurations],\MTCommitment_{\MTCommitmentIndex},\ROAdvState_{\MTCommitmentIndex})} \\
\IndentedBullet(\MTMessageVector_{\MTCommitmentIndex},\MTTrapdoor_{\MTCommitmentIndex}) \gets \MTMultiMultiExtractor{\MTNumConfigurations}(\MTConfigIndex_{\MTCommitmentIndex},\MTCommitment_{\MTCommitmentIndex},\ROTrace_{\MTCommitmentIndex}) \\
\IndentedBullet\text{add $\MTCommitmentIndex$ to $\MTConfigurationIdxSet{\MTConfigIndex_{\MTCommitmentIndex}}$} \\
\big((\MTCommitmentIndex_{\MTOpeningIndex} \in \MTConfigurationIdxSet{\MTOpeningIndex}, \MTMessageIndexSet_{\MTOpeningIndex},\MTMessageSubVector_{\MTOpeningIndex},\MTProof_{\MTOpeningIndex})\big)_{\MTOpeningIndex \in [\MTNumConfigurations]} \gets \MTAdversary^{\OracleFunction}(\ROAdvState_{\MTNumCommitments}) \\
\text{For $\MTOpeningIndex \in [\MTNumConfigurations]$: } \MTProof'_{\MTOpeningIndex} \gets \MTIdxOpen{\MTOpeningIndex}^{\ROFunction_{\MTOpeningIndex}}(\MTTrapdoor_{\MTCommitmentIndex_{\MTOpeningIndex}}, \MTMessageIndexSet_{\MTOpeningIndex})
}
\right]
\\&\leq
\MTMultiMultiExtractionError(\ROOutputSize,\MTMessageLengths,\ROQueryBound,\MTNumCommitments)
\enspace.
\end{align*}
Above:
\begin{itemize}[nolistsep]
  \item $\MTMultiMultiExtractionError(\ROOutputSize,\MTMessageLengths,\ROQueryBound,\MTNumCommitments) \leq \MTMultiExtractabilityExpression{\ROOutputSize}{\ROQueryBound}{\MTMessageSumLength}{\MTMaxDepth}{\MTNumCommitments}$ (at most $\MTMultiMultiExtractabilityShortExpression{\ROOutputSize}{\ROQueryBound}{\MTNumCommitments}{\MTMessageLength}$ if $\MTMultiMultiExtractabilityCondition{\ROQueryBound}{\MTMaxDepth}{\MTMessageSumLength}$);
  \item the total running time of $\MTMultiMultiExtractor{\MTNumConfigurations}$ across the $\MTNumCommitments$ invocations is
\begin{equation*}
\MTMultiMultiExtractorTimeFunction{\ROOutputSize}{\MTAlphabets}{\MTMessageLengths}{\MTSaltSize}{\ROQueryBound}{\MTNumCommitments}
= \MTMultiMultiExtractorTimeExpression
\enspace.
\end{equation*}
\end{itemize}
\end{lemma}

\begin{proof}
The stateful extractor $\MTMultiMultiExtractor{\MTNumConfigurations}$ is a straightforward generalization of the stateful extractor in the proof \Cref{lemma:mt-multi-extractability}. The extractor $\MTMultiMultiExtractor{\MTNumConfigurations}$ internally maintains the concatenation of all query-answer traces received so far (after the $\MTCommitmentIndex$-th invocation it stores the trace $\ROTrace \DefineEqual \ROTrace_{1} \concat \cdots \concat \ROTrace_{\MTCommitmentIndex}$), and uses the relevant Merkle commitment extractor on the appropriate subtrace to construct the message and trapdoor corresponding to the given commitment.
\begin{itemize}
\item[] $\MTMultiMultiExtractor{\MTNumConfigurations}(\MTConfigIndex_{\MTCommitmentIndex},\MTCommitment_{\MTCommitmentIndex},\ROTrace_{\MTCommitmentIndex})$:
\begin{enumerate}[nolistsep]
  \item Append $\ROTrace_{\MTCommitmentIndex}$ to the query-answer trace $\ROTrace$ stored in the internal state.
  \item Let $\ROTrace_{\MTConfigIndex_{\MTCommitmentIndex}}$ be the sub-trace corresponding to the oracle $\ROFunction_{\MTConfigIndex_{\MTCommitmentIndex}}$.
  \item Compute and output $(\MTMessageVector_{\MTCommitmentIndex},\MTTrapdoor_{\MTCommitmentIndex}) \DefineEqual \MTIdxExtractor{\MTConfigIndex_{\MTCommitmentIndex}}(\MTCommitment_{\MTCommitmentIndex},\ROTrace_{\MTConfigIndex_{\MTCommitmentIndex}})$.
\end{enumerate}
\end{itemize}
By \Cref{lemma:mt-extractability}, for every $\MTConfigIndex \in [\MTNumConfigurations]$, each invocation of $\MTIdxExtractor{\MTConfigIndex}$ runs in time at most $O\big( \ROQueryBound_{\MTConfigIndex} \cdot \MTMessageLength_{\MTConfigIndex} \cdot (\log\Cardinality{\MTAlphabet_{\MTConfigIndex}} + \MTSaltSize + \ROOutputSize)\big)$, where $\ROQueryBound_{\MTConfigIndex}$ is the number of queries by $\MTAdversary$ to the oracle $\ROFunction_{\MTConfigIndex}$; note that $\sum_{\MTConfigIndex \in [\MTNumConfigurations]} \ROQueryBound_{\MTConfigIndex} \leq \ROQueryBound$.\footnote{Recall that the time complexity of each $\MTIdxExtractor{\MTConfigIndex}$ can be improved via a suitable use of dynamic data structures in the internal state (and foregoing $\MTIdxExtractor{\MTConfigIndex}$ as a subroutine); see \Cref{remark:mt-multi-extractor-efficient-implementation}.} The $\MTNumCommitments$ invocations, each for some $\MTConfigIndex \in [\MTNumConfigurations]$, yield the claimed total running time:
\begin{align*}
\MTNumCommitments
\cdot
\max_{\MTConfigIndex \in [\MTNumConfigurations]}
O\big(\ROQueryBound_{\MTConfigIndex} \cdot \MTMessageLength_{\MTConfigIndex} \cdot (\log\Cardinality{\MTAlphabet_{\MTConfigIndex}} + \MTSaltSize + \ROOutputSize)\big)
= \MTMultiMultiExtractorTimeExpression
\enspace.
\end{align*}

Next, we discuss the extraction error. The analysis relies on \Cref{lemma:mt-multi-extractability} and the fact that each Merkle commitment has its own random oracle.

For each $\MTConfigIndex \in [\MTNumConfigurations]$ let $\ROQueryBound_{\MTConfigIndex}$ is the number of queries by $\MTAdversary$ to the oracle $\ROFunction_{\MTConfigIndex}$, and let $\MTNumCommitments_{\MTConfigIndex}$ be a bound on the number of commitments for $\MTSymbol_{\MTConfigIndex}$ output by $\MTAdversary$. Further below we argue that the probability in the lemma statement is at most $\sum_{\MTConfigIndex \in [\MTNumConfigurations]} \MTMultiExtractionError(\ROOutputSize,\MTMessageLength_{\MTConfigIndex},\ROQueryBound_{\MTConfigIndex},\MTNumCommitments_{\MTConfigIndex})$, where $\MTMultiExtractionError$ is the Merkle commitment multi-extraction error from \Cref{lemma:mt-multi-extractability}. This suffices to show the claimed bound because:
\begin{align*}
&\sum_{\MTConfigIndex \in [\MTNumConfigurations]} \MTMultiExtractionError(\ROOutputSize,\MTMessageLength_{\MTConfigIndex},\ROQueryBound_{\MTConfigIndex},\MTNumCommitments_{\MTConfigIndex})
\\& \leq \sum_{\MTConfigIndex \in [\MTNumConfigurations]}  \MTMultiExtractabilityExpression{\ROOutputSize}{\ROQueryBound_{\MTConfigIndex}}{\MTMessageLength_{\MTConfigIndex}}{\MTDepth_{\MTConfigIndex}}{\MTNumCommitments_{\MTConfigIndex}}
\EquationComment{by \Cref{lemma:mt-multi-extractability}}
\\& \leq \sum_{\MTConfigIndex \in [\MTNumConfigurations]}  \MTMultiExtractabilityExpression{\ROOutputSize}{\ROQueryBound_{\MTConfigIndex}}{\MTMessageLength_{\MTConfigIndex}}{\MTDepth_{\MTConfigIndex}}{\MTNumCommitments}
\EquationComment{since $\max_{\MTConfigIndex \in [\MTNumConfigurations]} \MTNumCommitments_{\MTConfigIndex} \leq \MTNumCommitments$}
\\& \leq
\frac{3}{2} \cdot \frac{(\ROQueryBound-1)\cdot\ROQueryBound}{2^{\ROOutputSize}}
+ \sum_{\MTConfigIndex \in [\MTNumConfigurations]} \frac{2 \cdot (\MTDepth_{\MTConfigIndex}+1) \cdot \MTMessageLength_{\MTConfigIndex}}{2^{\ROOutputSize}}
+ \frac{\MTNumCommitments \cdot \ROQueryBound}{2^{\ROOutputSize}}
\EquationComment{since $\sum_{\MTConfigIndex \in [\MTNumConfigurations]} \ROQueryBound_{\MTConfigIndex} \leq \ROQueryBound$}
\\& \leq
\MTMultiExtractabilityExpression{\ROOutputSize}{\ROQueryBound}{\MTMessageSumLength}{\MTMaxDepth}{\MTNumCommitments}
\enspace.
\end{align*}
We are left to argue the upper bound in terms of $\MTMultiExtractionError$.

For each $\MTOpeningIndex \in [\MTNumConfigurations]$ let $\Event_{\MTOpeningIndex}$ denote the event that, in the experiment of the lemma statement, $\Cardinality{\MTConfigurationIdxSet{\MTOpeningIndex}} \leq \MTNumCommitments_{\MTOpeningIndex}$ and either of the other two conditions holds for $\MTOpeningIndex$. The event in the lemma statement implies that $\lor_{\MTOpeningIndex \in [\MTNumConfigurations]} \Event_{\MTOpeningIndex}$, so its probability is at most $\Pr[\lor_{\MTOpeningIndex \in [\MTNumConfigurations]} \Event_{\MTOpeningIndex}] \leq \sum_{\MTOpeningIndex \in [\MTNumConfigurations]} \Pr[\Event_{\MTOpeningIndex}]$. We now argue that $\Pr[\Event_{\MTOpeningIndex}] \leq \MTMultiExtractionError(\ROOutputSize,\MTMessageLength_{\MTOpeningIndex},\ROQueryBound_{\MTOpeningIndex},\MTNumCommitments_{\MTOpeningIndex})$.

Consider the adversary $\MTAdversary_{\MTOpeningIndex}$ that, given oracle $\ROFunction_{\MTOpeningIndex}$, emulates $\MTAdversary$ as follows.
\begin{itemize}[noitemsep]
\item[] $\MTAdversary_{\MTOpeningIndex}^{\ROFunction_{\MTOpeningIndex}}$:
\begin{enumerate}[nolistsep]
  \item Lazily sample $\MTNumConfigurations-1$ oracles from $\RODistribution{\ROOutputSize}$ to extend the single oracle $\ROFunction_{\MTOpeningIndex}$ into a vector of oracles $(\ROFunction_{\MTConfigIndex})_{\MTConfigIndex \in [\MTNumConfigurations]}$ distributed according to $\RODistribution{\ROOutputSize}^{\MTNumConfigurations}$.
  \item Start emulating $\MTAdversary^{(\ROFunction_{\MTConfigIndex})_{\MTConfigIndex \in [\MTNumConfigurations]}}$.
  \item Whenever $\MTAdversary$ outputs a tuple $(\MTConfigIndex_{\MTCommitmentIndex},\MTCommitment_{\MTCommitmentIndex},\ROAdvState_{\MTCommitmentIndex})$ do: if $\MTConfigIndex_{\MTCommitmentIndex} \neq \MTOpeningIndex$ then do nothing; else if $\MTConfigIndex_{\MTCommitmentIndex} = \MTOpeningIndex$ then output $(\MTCommitment_{\MTCommitmentIndex},\ROAdvState_{\MTCommitmentIndex})$.
  \item When $\MTAdversary$ halts and outputs $\big((\MTCommitmentIndex_{\MTOpeningIndex}, \MTMessageIndexSet_{\MTOpeningIndex},\MTMessageSubVector_{\MTOpeningIndex},\MTProof_{\MTOpeningIndex})\big)_{\MTOpeningIndex \in [\MTNumConfigurations]}$, halt and output $(\MTCommitmentIndex_{\MTOpeningIndex}, \MTMessageIndexSet_{\MTOpeningIndex},\MTMessageSubVector_{\MTOpeningIndex},\MTProof_{\MTOpeningIndex})$.
\end{enumerate}
\end{itemize}
The adversary $\MTAdversary_{\MTOpeningIndex}$ makes at most $\ROQueryBound_{\MTOpeningIndex}$ queries to the oracle $\ROFunction_{\MTOpeningIndex}$ and outputs at most $\MTNumCommitments_{\MTConfigIndex}$ commitments. Moreover, if $\Event_{\MTOpeningIndex}$ holds for $\MTAdversary$ then $\MTAdversary_{\MTOpeningIndex}$ satisfies the conditions in the probability statement of \Cref{lemma:mt-multi-extractability}. We conclude that $\Pr[\Event_{\MTOpeningIndex}] \leq \MTMultiExtractionError(\ROOutputSize,\MTMessageLength_{\MTOpeningIndex},\ROQueryBound_{\MTOpeningIndex},\MTNumCommitments_{\MTOpeningIndex})$.
\end{proof}

\begin{remark}[multiple openings per configuration]
\label{remark:mt-multi-configuration-multi-extractability-with-multiple-openings}
\Cref{lemma:mt-multi-configuration-multi-extractability} considers the case where the adversary outputs at most $\MTNumCommitments$ commitments per configuration and then one opening per configuration, which suffices for the purposes of this book. More generally (and analogously to \Cref{remark:mt-multi-extractability-with-multiple-openings}), one can consider the case where, for each $\MTConfigIndex \in [\MTNumConfigurations]$, the adversary outputs at most $\MTNumCommitments_{\MTConfigIndex}$ commitments and $\MTNumOpenings_{\MTConfigIndex}$ openings for configuration $\MTSymbol_{\MTConfigIndex}$. The proof of \Cref{lemma:mt-multi-configuration-multi-extractability} can be extended to this case, establishing an upper bound $\MTMultiMultiExtractionError(\ROOutputSize,\MTMessageLengths,\ROQueryBound,(\MTNumCommitments_{\MTConfigIndex})_{\MTConfigIndex \in [\MTNumConfigurations]},(\MTNumOpenings_{\MTConfigIndex})_{\MTConfigIndex \in [\MTNumConfigurations]})$ for the following probability:
\begin{equation*}
\Pr\left[
\begin{array}{l}
\Cardinality{\MTConfigurationIdxSet{1}} \leq \MTNumCommitments_{1},\dots,\Cardinality{\MTConfigurationIdxSet{\MTNumConfigurations}} \leq \MTNumCommitments_{\MTNumConfigurations} \\
\text{and the following disjunction holds:} \\
\left(
\begin{array}{l}
\exists\,\MTConfigIndex \in [\MTNumConfigurations]\;
\exists\,\MTOpeningIndex \in [\MTNumOpenings_{\MTConfigIndex}] : \\
\MTIdxCheck{\MTConfigIndex}^{\ROFunction_{\MTConfigIndex}}(\MTCommitment_{\MTCommitmentIndex_{\MTConfigIndex,\MTOpeningIndex}},\MTMessageIndexSet_{\MTConfigIndex,\MTOpeningIndex},\MTMessageSubVector_{\MTConfigIndex,\MTOpeningIndex},\MTProof_{\MTConfigIndex,\MTOpeningIndex})=1 \\
\land\,(\MTMessageVector_{\MTCommitmentIndex_{\MTConfigIndex,\MTOpeningIndex}}[\MTMessageIndexSet_{\MTConfigIndex,\MTOpeningIndex}] \neq \MTMessageSubVector_{\MTConfigIndex,\MTOpeningIndex} \lor \MTProof'_{\MTConfigIndex,\MTOpeningIndex} \neq \MTProof_{\MTConfigIndex,\MTOpeningIndex})
\end{array}
\right)
\\
\lor \\
\left(
\begin{array}{l}
\exists\,\MTOpeningIndex \in [\MTNumConfigurations]\;
\exists\,\MTCommitmentIndex_{1},\MTCommitmentIndex_{2} \in \MTConfigurationIdxSet{\MTOpeningIndex} : \\
\MTCommitment_{\MTCommitmentIndex_{1}}
=
\MTCommitment_{\MTCommitmentIndex_{2}} \\
\land
(\MTMessageVector_{\MTCommitmentIndex_{1}},\MTTrapdoor_{\MTCommitmentIndex_{1}})
\neq
(\MTMessageVector_{\MTCommitmentIndex_{2}},\MTTrapdoor_{\MTCommitmentIndex_{2}})
\end{array}
\right)
\end{array}
\GivenExperiment
\StateExperiment{
\OracleFunction=(\ROFunction_{\MTConfigIndex})_{\MTConfigIndex \in [\MTNumConfigurations]} \gets \RODistribution{\ROOutputSize}^{\MTNumConfigurations} \\
\ROAdvState_{0} \DefineEqual \bot \\
\MTConfigurationIdxSet{1},\dots,\MTConfigurationIdxSet{\MTNumConfigurations} \DefineEqual \emptyset \\
\text{For $\MTCommitmentIndex=1,2,\dots$:} \\
\IndentedBullet\ROInputOutputAndTrace{\OracleFunction}{\MTAdversary}{\ROAdvState_{\MTCommitmentIndex-1}}{\ROTrace_{\MTCommitmentIndex}}{(\MTConfigIndex_{\MTCommitmentIndex} \in [\MTNumConfigurations],\MTCommitment_{\MTCommitmentIndex},\ROAdvState_{\MTCommitmentIndex})} \\
\IndentedBullet(\MTMessageVector_{\MTCommitmentIndex},\MTTrapdoor_{\MTCommitmentIndex}) \gets \MTMultiMultiExtractor{\MTNumConfigurations}(\MTConfigIndex_{\MTCommitmentIndex},\MTCommitment_{\MTCommitmentIndex},\ROTrace_{\MTCommitmentIndex}) \\
\IndentedBullet\text{add $\MTCommitmentIndex$ to $\MTConfigurationIdxSet{\MTConfigIndex_{\MTCommitmentIndex}}$} \\
\big((\MTCommitmentIndex_{\MTConfigIndex,\MTOpeningIndex} \in \MTConfigurationIdxSet{\MTOpeningIndex}, \MTMessageIndexSet_{\MTConfigIndex,\MTOpeningIndex},\MTMessageSubVector_{\MTConfigIndex,\MTOpeningIndex},\MTProof_{\MTConfigIndex,\MTOpeningIndex})\big)_{\substack{\MTConfigIndex \in [\MTNumConfigurations] \\ \MTOpeningIndex \in [\MTNumOpenings_{\MTConfigIndex}]}} \gets \MTAdversary^{\OracleFunction}(\ROAdvState_{\MTNumCommitments}) \\
\text{For $\MTConfigIndex \in [\MTNumConfigurations], \MTOpeningIndex \in [\MTNumOpenings_{\MTConfigIndex}]$: } \\
\IndentedBullet \MTProof'_{\MTConfigIndex,\MTOpeningIndex} \gets \MTIdxOpen{\MTConfigIndex}^{\ROFunction_{\MTConfigIndex}}(\MTTrapdoor_{\MTCommitmentIndex_{\MTConfigIndex,\MTOpeningIndex}}, \MTMessageIndexSet_{\MTConfigIndex,\MTOpeningIndex})
}
\right]
\enspace.
\end{equation*}
An argument similar to the proof of \Cref{lemma:mt-multi-configuration-multi-extractability} shows that if the adversary $\MTAdversary$ makes at most $(\ROQueryBound_{\MTConfigIndex})_{\MTConfigIndex \in [\MTNumConfigurations]}$ queries to the oracles $(\ROFunction_{\MTConfigIndex})_{\MTConfigIndex \in [\MTNumConfigurations]}$ respectively then $\MTMultiMultiExtractionError(\ROOutputSize,\MTMessageLengths,\ROQueryBound,(\MTNumCommitments_{\MTConfigIndex})_{\MTConfigIndex \in [\MTNumConfigurations]},(\MTNumOpenings_{\MTConfigIndex})_{\MTConfigIndex \in [\MTNumConfigurations]}) \leq \sum_{\MTConfigIndex \in [\MTNumConfigurations]} \MTMultiExtractionError(\ROOutputSize,\MTMessageLength_{\MTConfigIndex},\ROQueryBound_{\MTConfigIndex},\MTNumCommitments_{\MTConfigIndex},\MTNumOpenings_{\MTConfigIndex})$, where $\MTMultiExtractionError$ is the bound for multiple openings discussed in \Cref{remark:mt-multi-extractability-with-multiple-openings}. The expression can then be simplified using the fact that $\sum_{\MTConfigIndex \in [\MTNumConfigurations]} \ROQueryBound_{\MTConfigIndex} \leq \ROQueryBound$.
\end{remark}

%%%%%%%%%%%%%%%%%%%%%%%%%%%%%%%%%%%%%%%%%%%%%%%%%%%%%%%%%%%%%%%%%%%%%%%%%%%%%%%
%%%%%%%%%%%%%%%%%%%%%%%%%%%%%%%%%%%%%%%%%%%%%%%%%%%%%%%%%%%%%%%%%%%%%%%%%%%%%%%
%%%%%%%%%%%%%%%%%%%%%%%%%%%%%%%%%%%%%%%%%%%%%%%%%%%%%%%%%%%%%%%%%%%%%%%%%%%%%%%
\section{Hiding}
\label{section:merkle-commitment-hiding}

The binding and extractability properties of Merkle commitments protect the receiver from malicious senders. Merkle commitments additionally satisfy certain \emph{hiding} properties, which instead protect the sender. Informally, hiding ensures that information about the committed message that the sender did not open is hidden from the receiver. The mechanism that enables this is the random salts used to commit to each message entry at the leaf layer of the tree (see \Cref{step:mt-sample-salts} in $\MTCommit$); in fact, this mechanism can be viewed as using the basic commitment scheme $\CMSymbol$ (from \Cref{chapter:basic-commitment}) to commit to each message entry.

We discuss hiding properties of Merkle commitments in two steps. In \Cref{section:merkle-commitment-hiding-1} we explain how a Merkle commitment reveals no information about the committed message. Then in \Cref{section:merkle-commitment-hiding-2} we explain how additionally providing an opening proof reveals no information about the message beyond the opened entries of the message.

%%%%%%%%%%%%%%%%%%%%%%%%%%%%%%%%%%%%%%%%%%%%%%%%%%%%%%%%%%%%%%%%%%%%%%%%%%%%%%%
%%%%%%%%%%%%%%%%%%%%%%%%%%%%%%%%%%%%%%%%%%%%%%%%%%%%%%%%%%%%%%%%%%%%%%%%%%%%%%%
\subsection{Merkle commitments hide the message}
\label{section:merkle-commitment-hiding-1}

We prove that a Merkle commitment $\MTCommitment$ reveals essentially no information about the underlying committed message $\MTMessageVector$. Formally this is captured via the simulation paradigm: there exists a polynomial-time probabilistic algorithm $\MTRootSimulator$ that samples a string, without any information about the message $\MTMessageVector$, whose distribution is statistically close to $\MTCommitment$. This is analogous to the hiding property for the basic commitment scheme $\CMSymbol$ in \Cref{chapter:basic-commitment} (see \Cref{lemma:cm-hiding} in \Cref{section:basic-commitment-hiding}). Below we state the hiding property and prove the statistical error bound. Note that the statistical error depends on several parameters:
\begin{inparaenum}[(a)]
  \item the output size $\ROOutputSize \in \Naturals$ of the random oracle;
  \item the salt size $\CMSaltSize \in \Naturals$;
  \item the message size $\MTMessageLength \in \Naturals$;
  \item the number of queries $\ROQueryBound$ to the random oracle made by an algorithm trying to distinguish between the two distributions.
\end{inparaenum}

\begin{lemma}[$\MTCommitment$ is hiding]
\label{lemma:mt-root-hiding}
Let $\MTSymbol \DefineEqual \MTConstructor{\ROOutputSize}{\MTAlphabet}{\MTMessageLength}{\MTSaltSize}$. There exists a polynomial-time probabilistic algorithm $\MTRootSimulator$ such that for every query bound $\ROQueryBound \in \Naturals$ and $\ROQueryBound$-query algorithm $\MTAdversary$, the following two distributions are $\MTRootHidingError(\ROOutputSize,\MTAlphabet,\MTMessageLength,\MTSaltSize,\ROQueryBound)$-close in statistical distance:
\begin{equation*}
\left\{
  \MTAdversary^{\ROFunction}(\ROAdvState,\MTCommitment)
\GivenExperiment
\StateExperiment{
  \ROFunction \gets \RODistribution{\ROOutputSize} \\
  (\MTMessageVector \in \MTAlphabet^{\MTMessageLength},\ROAdvState) \gets \MTAdversary^{\ROFunction} \\
  (\MTCommitment,\MTTrapdoor) \gets \MTCommit^{\ROFunction}(\MTMessageVector)
}
\right\}
\TextAndInMath
\left\{
  \MTAdversary^{\ROFunction}(\ROAdvState,\MTCommitment)
\GivenExperiment
\StateExperiment{
  \ROFunction \gets \RODistribution{\ROOutputSize} \\
  (\MTMessageVector \in \MTAlphabet^{\MTMessageLength},\ROAdvState) \gets \MTAdversary^{\ROFunction} \\
  \MTCommitment \gets \MTRootSimulator^{\ROFunction}
}
\right\}
\enspace.
\end{equation*}
Above,
\begin{equation*}
\MTRootHidingError(\ROOutputSize,\MTAlphabet,\MTMessageLength,\MTSaltSize,\ROQueryBound)
\leq
\MTMessageLength \cdot \CMHidingError(\ROOutputSize,\log\Cardinality{\MTAlphabet}, \MTSaltSize,\ROQueryBound)
+ \MTMessageLength \cdot \CMHidingError(\ROOutputSize,0,2\ROOutputSize,\ROQueryBound)
\end{equation*}
where $\CMHidingError$ is the hiding error of $\CMSymbol = \CMConstructor{\ROOutputSize}{\CMMessageLength}{\CMSaltSize}$ from \Cref{lemma:cm-hiding}. In particular, plugging-in the bounds from the lemma we get that
\begin{equation*}
\MTRootHidingError(\ROOutputSize,\MTAlphabet,\MTMessageLength,\MTSaltSize,\ROQueryBound)
\leq
\begin{cases}
\MTRootHidingErrorExpression
& \text{ if $\ROQueryBound<\infty$} \\
\MTMessageLength \cdot (2^{-\frac{\MTSaltSize-\ROOutputSize}{2}} + 2^{-\frac{\MTSaltSize}{2}}
+ 2^{-\frac{\ROOutputSize}{2}}) & \text{ if $\ROQueryBound=\infty$ and $\sqrt{\log\Cardinality{\MTAlphabet}} \leq \MTMessageLength$}
\end{cases}
\enspace.
\end{equation*}
\end{lemma}



\begin{proof}
The simulator $\MTRootSimulator$ outputs a random string in $\Bits^{\ROOutputSize}$. The proof follows a straightforward hybrid argument, in which we view each vertex in the Merkle tree as a basic commitment.

Consider the leaf layer. Each leaf is a basic commitment as defined in \Cref{chapter:basic-commitment}. Using \Cref{lemma:cm-hiding}, we can replace a leaf by a uniform random string over $\Bits^{\ROOutputSize}$ and lose an additive term of $\CMHidingError(\ROOutputSize,\log\Cardinality{\MTAlphabet},\CMSaltSize,\ROQueryBound)$ in the statistical distance. Hence replacing all leaves with random strings gives us the error of $\MTMessageLength \cdot \CMHidingError(\ROOutputSize,\log\Cardinality{\MTAlphabet},\CMSaltSize,\ROQueryBound)$.

We continue layer by layer up to the root. Each vertex of the tree can be viewed as a basic commitment where the left child is the empty message and the right child is a salt of size $2\ROOutputSize$. Replacing that vertex with a random string gives an error of $\CMHidingError(\ROOutputSize,0,2\ROOutputSize,\ROQueryBound)$. We replace each vertex, one by one, with a random string, until we reach the root. The number of vertices is at most $\MTMessageLength$, so overall this process yields an error of $\MTMessageLength \cdot \CMHidingError(\ROOutputSize,0,2\ROOutputSize,\ROQueryBound)$.

Overall, the total error accumulated is at most $\MTMessageLength \cdot \CMHidingError(\ROOutputSize,\log\Cardinality{\MTAlphabet},\MTSaltSize,\ROQueryBound) + \MTMessageLength \cdot \CMHidingError(\ROOutputSize,0,2\ROOutputSize,\ROQueryBound)$.
\end{proof}



































































































\begin{remark}
In many cases, it is useful to have the simulator output a uniformly random string, which is the case of our simulator. However, the second term in the error expression for $\ROQueryBound=\infty$ can be avoided at the price of changing the simulator to output the root of a Merkle commitment starting from random leaves instead of a uniformly random string. In this case, we only need to replace the leaves with simulated ones and incur only an error of $\MTMessageLength \cdot \CMHidingError(\ROOutputSize,\log\Cardinality{\MTAlphabet},\MTSaltSize,\infty)$.
\end{remark}





%%%%%%%%%%%%%%%%%%%%%%%%%%%%%%%%%%%%%%%%%%%%%%%%%%%%%%%%%%%%%%%%%%%%%%%%%%%%%%%
%%%%%%%%%%%%%%%%%%%%%%%%%%%%%%%%%%%%%%%%%%%%%%%%%%%%%%%%%%%%%%%%%%%%%%%%%%%%%%%
\subsection{Merkle opening proofs hide the rest of the message}
\label{section:merkle-commitment-hiding-2}

We prove that a Merkle opening proof $\MTProof$, which is used to authenticate the entries of a message in a subset $\MTMessageIndexSet \subseteq [\MTMessageLength]$ of locations, hides the entries of the message outside of $\MTMessageIndexSet$ (i.e., with locations in $[\MTMessageLength] \setminus \MTMessageIndexSet$). Formally, this is captured via the simulation paradigm: there exists a polynomial-time probabilistic algorithm $\MTSimulator$ that, given as input the entries of the message in $\MTMessageIndexSet$, samples a string whose distribution is statistically close to $\MTProof$. For convenience to later sections, the property that we consider below additionally includes the Merkle commitment $\MTCommitment$ in order to provide a \emph{joint} hiding property across $\MTCommitment$ and $\MTProof$. In particular, the property below implies the one above (\Cref{lemma:mt-root-hiding}) by taking $\MTMessageIndexSet=\emptyset$ (as in this case the Merkle opening proof is empty and the simulator receives no entries of the message).

\begin{lemma}[$\MTCommitment$ and $\MTProof$ are hiding]
\label{lemma:mt-privacy}
Let $\MTSymbol \DefineEqual \MTConstructor{\ROOutputSize}{\MTAlphabet}{\MTMessageLength}{\MTSaltSize}$. There exists a polynomial-time probabilistic algorithm $\MTSimulator$ such that for every opening size $\MTMessageIndexSetSize \in [\MTMessageLength]$, query bound $\ROQueryBound \in \Naturals$, and $\ROQueryBound$-query algorithm $\MTAdversary$, the following two distributions are $\MTHidingError(\ROOutputSize,\MTAlphabet,\MTMessageLength,\MTSaltSize,\MTMessageIndexSetSize,\ROQueryBound)$-close in statistical distance:
\begin{align*}
&\left\{
  \MTAdversary^{\ROFunction}(\ROAdvState,\MTCommitment,\MTProof)
\GivenExperiment
\StateExperiment{
  \ROFunction \gets \RODistribution{\ROOutputSize} \\
  \left(\MTMessageVector \in \MTAlphabet^{\MTMessageLength},\MTMessageIndexSet \in \binom{[\MTMessageLength]}{\MTMessageIndexSetSize},\ROAdvState\right) \gets \MTAdversary^{\ROFunction} \\
  (\MTCommitment,\MTTrapdoor) \gets \MTCommit^{\ROFunction}(\MTMessageVector) \\
  \MTProof \gets \MTOpen^{\ROFunction}(\MTTrapdoor, \MTMessageIndexSet)
}
\right\}
\TextAndInMath
\\
&\left\{
  \MTAdversary^{\ROFunction}(\ROAdvState,\MTCommitment,\MTProof)
\GivenExperiment
\StateExperiment{
  \ROFunction \gets \RODistribution{\ROOutputSize} \\
  \left(\MTMessageVector \in \MTAlphabet^{\MTMessageLength},\MTMessageIndexSet \in \binom{[\MTMessageLength]}{\MTMessageIndexSetSize},\ROAdvState\right) \gets \MTAdversary^{\ROFunction} \\
  (\MTCommitment,\MTProof) \gets \MTSimulator^{\ROFunction}(\MTMessageIndexSet,\MTMessageVector[\MTMessageIndexSet])
}
\right\}
\enspace.
\end{align*}
Letting $\MTRootHidingError$ be the error defined and upper bounded in \Cref{lemma:mt-root-hiding}, we define the error function $\MTHidingError$ used above as follows:
\begin{itemize}
  \item if $\MTMessageIndexSetSize=0$ then $\MTHidingError(\ROOutputSize,\MTAlphabet,\MTMessageLength,\MTSaltSize,\MTMessageIndexSetSize,\ROQueryBound) \DefineEqual \MTRootHidingError(\ROOutputSize,\MTAlphabet,\MTMessageLength,\MTSaltSize,\ROQueryBound)$; and
  \item if $\MTMessageIndexSetSize>0$ then $\MTHidingError(\ROOutputSize,\MTAlphabet,\MTMessageLength,\MTSaltSize,\MTMessageIndexSetSize,\ROQueryBound) \DefineEqual \max_{\MTMessageIndexSet \in \binom{[\MTMessageLength]}{\MTMessageIndexSetSize}} \sum_{(\MTLayer,\MTMessageIndex) \in \MTCoPathOptVertexSet} \MTRootHidingError(\ROOutputSize,\MTAlphabet,2^{\MTDepth-\MTLayer},\MTSaltSize,\ROQueryBound)$, where $\MTCoPathOptVertexSet$ is the set of vertices in the copaths but not the paths for $\MTMessageIndexSet$, i.e.,
\begin{equation*}
\MTCoPathOptVertexSet
\DefineEqual
\MTCoPath(\MTMessageIndexSet) \setminus \MTPath(\MTMessageIndexSet)
=
\left(\idxcup{\MTMessageIndex \in \MTMessageIndexSet}\MTCoPath(\MTMessageIndex)\right)
\setminus
\left(\idxcup{\MTMessageIndex \in \MTMessageIndexSet}\MTPath(\MTMessageIndex)\right)
\enspace.
\end{equation*}
\end{itemize}
In particular, if $\MTMessageIndexSetSize>0$ then taking $\MTHidingError(\ROOutputSize,\MTAlphabet,\MTMessageLength,\MTSaltSize,\MTMessageIndexSetSize,\ROQueryBound) \DefineEqual \MTHidingErrorExpression$ suffices because every copath contains one vertex per layer (except the root layer):
\begin{align*}
\sum_{(\MTLayer,\MTMessageIndex) \in \MTCoPath(\MTMessageIndexSet) \setminus \MTPath(\MTMessageIndexSet)} \MTRootHidingError(\ROOutputSize,\MTAlphabet,2^{\MTDepth-\MTLayer},\MTSaltSize,\ROQueryBound)
&\leq
\sum_{(\MTLayer,\MTMessageIndex) \in \MTCoPath(\MTMessageIndexSet)} \MTRootHidingError(\ROOutputSize,\MTAlphabet,2^{\MTDepth-\MTLayer},\MTSaltSize,\ROQueryBound)
\\&\leq
\sum_{\MTMessageIndex \in \MTMessageIndexSet} \sum_{\MTLayer \in [\MTDepth]} \MTRootHidingError(\ROOutputSize,\MTAlphabet,2^{\MTDepth-\MTLayer},\MTSaltSize,\ROQueryBound)
\\&=
\MTHidingErrorExpression
\enspace;
\end{align*}
by \Cref{remark:mtzk-plugging-in-root-hiding-bound} if $\ROQueryBound<\infty$ this expression is at most $\MTHidingErrorLooseExpression$.
\end{lemma}



Note that if $\MTMessageIndexSetSize=0$ then the adversary receives only a Merkle commitment $\MTCommitment$; the adversary receives no message entries or opening proof $\MTProof$ (it is empty). This degenerate case corresponds to the setting of \Cref{lemma:mt-root-hiding}, so the statistical distance is upper bounded by $\MTRootHidingError(\ROOutputSize,\MTAlphabet,\MTMessageLength,\MTSaltSize,\ROQueryBound)=\MTRootHidingError(\ROOutputSize,\MTAlphabet,2^{\MTDepth},\MTSaltSize,\ROQueryBound)$. Henceforth we assume that $\MTMessageIndexSetSize>0$, which is the \DoQuote{interesting} case of \Cref{lemma:mt-privacy}.

First we study the case of simulating (the Merkle commitment and) the opening proof for a single leaf and then study the case of multiple leaves.

\parhead{Single leaf}
In the case of a single leaf, the simulator receives a leaf index $\MTMessageIndex \in [\MTMessageLength]$ and a message entry $\MTMessageSubVector \in \MTAlphabet$, and is tasked to sample a Merkle commitment $\MTCommitment$ and authentication path $\MTAuthPath$ whose distribution is close to that obtained by honestly committing to any message whose $\MTMessageIndex$-th entry is $\MTMessageSubVector$ and then opening that location. Informally, the simulator does this by sampling the salt for the leaf at random (just like $\MTCommit$ does), sampling the commitments along $\MTCoPath(\MTMessageIndex)$ via $\MTRootSimulator$ (rather than obtaining them from a computation on a message), and computing from these the implied Merkle commitment $\MTCommitment$ according to the way an authentication path is checked. In more detail, the simulator is constructed as follows.

\begin{construction}
\label{construction:mt-simulator-single-leaf}
The simulator for the case of a single leaf works as follows.
\begin{itemize}

  \item[] $\MTSimulator^{\ROFunction}(\MTMessageIndex,\MTMessageSubVector)$:
  \begin{enumerate}[nolistsep]

    \item Sample a random salt $\MTSaltString_{\MTMessageIndex} \in \Bits^{\MTSaltSize}$.

    \item Compute the commitment $\MTVertexLabelByName{\MTDepth}{\MTMessageIndex} \DefineEqual \ROFunction(\MTMessageSubVector,\MTSaltString_{\MTMessageIndex})$.

    \item For $\MTLayer = \MTDepth-1,\dots,1,0$:
    \begin{enumerate}[nolistsep]

      \item Sample the commitment $\MTVertexLabelByCoPath{\MTMessageIndex}{\MTLayer+1} \gets \MTRootSimulator^{\ROFunction}$.

      \item If $\MTPathIndex{\MTMessageIndex}{\MTLayer+1}$ is an odd vertex then set $\MTLeftVertexLabel \DefineEqual \MTVertexLabelByPath{\MTMessageIndex}{\MTLayer+1}$ and $\MTRightVertexLabel \DefineEqual \MTVertexLabelByCoPath{\MTMessageIndex}{\MTLayer+1}$.

      \item If $\MTPathIndex{\MTMessageIndex}{\MTLayer+1}$ is an even vertex then set $\MTLeftVertexLabel \DefineEqual \MTVertexLabelByCoPath{\MTMessageIndex}{\MTLayer+1}$ and $\MTRightVertexLabel \DefineEqual \MTVertexLabelByPath{\MTMessageIndex}{\MTLayer+1}$.

      \item Compute the commitment $\MTVertexLabelByPath{\MTMessageIndex}{\MTLayer} \DefineEqual \ROFunction(\MTLeftVertexLabel,\MTRightVertexLabel)$.

    \end{enumerate}

    \item Set $\MTCommitment \DefineEqual \MTVertexLabelByName{0}{1}$.

    \item Set $\MTAuthPath \DefineEqual \big(\MTSaltString_{\MTMessageIndex},(\MTVertexLabelByCoPath{\MTMessageIndex}{\MTLayer})_{\MTLayer \in \{1,\dots,\MTDepth\}}\big)$. (For convenience $\MTLayer$ now ranges from $1$ to $\MTDepth$.)

    \item Output $(\MTCommitment,\MTAuthPath)$.

  \end{enumerate}

\end{itemize}
\end{construction}

The output of this simulator always validates, i.e., it is such that $\MTCheck^{\ROFunction}(\MTCommitment,\MTMessageIndex,\MTMessageSubVector,\MTAuthPath)=1$. The number of queries to the random oracle is $\MTDepth+1$. Intuitively, the statistical distance from the real distribution is at most $\MTPathHidingErrorExpression$ because of the following reasons.
\begin{itemize}[noitemsep]
  \item The salt $\MTSaltString_{\MTMessageIndex}$ in $\MTAuthPath$ is distributed as in the real distribution.
  \item For every $\MTLayer \in [\MTDepth]$, the commitment $\MTVertexLabelByCoPath{\MTMessageIndex}{\MTLayer}$ in $\MTAuthPath$ can be viewed as a root of a subtree of size $2^{\MTDepth-\MTLayer}$, and simulating it via $\MTRootSimulator^{\ROFunction}$ contributes a statistical error of $\MTRootHidingError(\ROOutputSize,\MTAlphabet,2^{\MTDepth-\MTLayer},\MTSaltSize,\ROQueryBound)$.
  \item The Merkle commitment $\MTCommitment$ is determined by the other values and the random oracle.
\end{itemize}

\parhead{Multiple leaves}
The case of simulating a Merkle commitment and opening proof for multiple leaves $\MTMessageIndexSet$ may superficially look rather straightforward: for each leaf index $\MTMessageIndex \in \MTMessageIndexSet$, sample an authentication path for $\MTMessageIndex$ by following the simulator for a single leaf in \Cref{construction:mt-simulator-single-leaf}.

While intuitive, the above strategy does \emph{not} work. The problem is that the commitments across multiple authentication paths in $\MTProof$ may not be consistent with one another. First, the same vertex may appear in multiple copaths, and thus its commitment is sampled multiple times inconsistently; this problem is straightforwardly addressed by not sampling the commitment for a vertex if it was already sampled.\footnote{For example, for a message of size $\MTMessageLength=4$, both $\MTCoPath(1)$ and $\MTCoPath(2)$ contain the vertex $(1,2)$.} Second, and this is more subtle, a commitment sampled as part of the copath for a leaf index may be the result of an oracle query to check the copath of another leaf index.\footnote{For example, for a message of size $\MTMessageLength=4$, $\MTCoPath(1)$ contains the vertex $(1,2)$, and if the opening proof contains authentication paths for entries $1$ and $3$, then the commitment of the vertex $(1,2)$ must be the output of the random oracle applied to the commitment of vertices $(2,3)$ (derived from entry $3$ of the message and its salt) and of $(2,4)$ (provided in $\MTCoPath(3)$).} These considerations imply that the simulator should produce the opening proof $\MTProof$ in a more global way: sample all the commitments in $\MTProof$ that do not depend on other commitments, and then derive from these any remaining commitments in $\MTProof$ using the random oracle.

In more detail, an opening proof $\MTProof$ for an index set $\MTMessageIndexSet$ contains commitments for vertices in $\MTCoPath(\MTMessageIndexSet)$ (vertices that belong to the copath corresponding to some location in $\MTMessageIndexSet$). For every $\MTMessageIndex \in \MTMessageIndexSet$, checking the authentication path corresponding to $\MTMessageIndex$ involves computing the commitments for vertices in $\MTPath(\MTMessageIndex)$ from the commitments for vertices in $\MTCoPath(\MTMessageIndex)$ (and the message entry and salt). Hence, the commitments in $\MTProof$ that are derived from other commitments are those in $\MTPath(\MTMessageIndexSet)$. We conclude that the commitments in $\MTProof$ that are to be sampled are those corresponding to vertices in $\MTCoPath(\MTMessageIndexSet) \setminus \MTPath(\MTMessageIndexSet)$. These considerations lead to the following simulator.

\begin{construction}
\label{construction:mt-simulator}
Given query access to a random oracle $\ROFunction \in \RODistribution{\ROOutputSize}$, the simulator $\MTSimulator$ receives as input a subset $\MTMessageIndexSet$ and corresponding entries $(\MTMessageEntry_{\MTMessageIndex})_{\MTMessageIndex \in \MTMessageIndexSet}$ of a message, and works as follows.
\begin{itemize}[noitemsep]

  \item[] $\MTSimulator^{\ROFunction}(\MTMessageIndexSet,(\MTMessageEntry_{\MTMessageIndex})_{\MTMessageIndex \in \MTMessageIndexSet})$:
\begin{enumerate}[noitemsep]
  \item If $\MTMessageIndexSet=\emptyset$, set $\MTCommitment \gets \MTRootSimulator^{\ROFunction}$ and $\MTProof \DefineEqual \bot$, and output $(\MTCommitment,\MTProof)$.
  \item For every $\MTMessageIndex \in \MTMessageIndexSet$, sample a random salt $\MTSaltString_{\MTMessageIndex} \in \Bits^{\MTSaltSize}$.
  \item For every $\MTMessageIndex \in \MTMessageIndexSet$, compute the commitment $\MTVertexLabelByName{\MTDepth}{\MTMessageIndex} \DefineEqual \ROFunction(\MTMessageEntry_{\MTMessageIndex},\MTSaltString_{\MTMessageIndex})$.
  \item For every $(\MTLayer,\MTMessageIndex) \in \MTCoPath(\MTMessageIndexSet) \setminus \MTPath(\MTMessageIndexSet)$, set $\MTVertexLabelByName{\MTLayer}{\MTMessageIndex} \gets \MTRootSimulator^{\ROFunction}$.
  \item For $\MTLayer = \MTDepth,\dots,1$ (in this order) and $\MTMessageIndex \in [2^{\MTLayer}]$: \\ \phantom{aa} if $(\MTLayer,\MTMessageIndex) \in \MTPath(\MTMessageIndexSet)$ then set $\MTVertexLabelByName{\MTLayer}{\MTMessageIndex} \DefineEqual \ROFunction(\MTVertexLabelByName{\MTLayer+1}{2\MTMessageIndex-1},\MTVertexLabelByName{\MTLayer+1}{2\MTMessageIndex}) \in \Bits^{\ROOutputSize}$.
    \item Set $\MTCommitment \DefineEqual \MTVertexLabelByName{0}{1}$. (Note that $(0,1) \in \MTPath(\MTMessageIndexSet)$.)
    \item For every $\MTMessageIndex \in \MTMessageIndexSet$, set $\MTAuthPath_{\MTMessageIndex} \DefineEqual \big(\MTSaltString_{\MTMessageIndex},(\MTVertexLabel_{\MTGraphVertex})_{\MTGraphVertex \in \MTCoPath(\MTMessageIndex)}\big) = \big(\MTSaltString_{\MTMessageIndex},(\MTVertexLabelByCoPath{\MTMessageIndex}{\MTLayer})_{\MTLayer \in \{1,\dots,\MTDepth\}}\big)$.
  \item Output $(\MTCommitment,\MTProof)$ where $\MTProof \DefineEqual (\MTAuthPath_{\MTMessageIndex})_{\MTMessageIndex \in \MTMessageIndexSet}$.
\end{enumerate}
\end{itemize}
\end{construction}

\begin{definition}
\label{definition:independent-vertices-in-tree}
Two vertices $\MTGraphVertex_{1},\MTGraphVertex_{2}$ in the tree $\MTGraph_{\MTMessageLength}$ are \defemph{independent} if the subtree rooted at $\MTGraphVertex_{1}$ is disjoint from the subtree rooted at $\MTGraphVertex_{2}$.
\end{definition}

\begin{claim}
\label{claim:copaths-minus-paths-are-independent}
The vertices in $\MTCoPath(\MTMessageIndexSet) \setminus \MTPath(\MTMessageIndexSet)$ are pairwise independent.
\end{claim}

\begin{proof}
Suppose there exist $\MTGraphVertex_{1}$ and $\MTGraphVertex_{2}$ in $\MTCoPath(\MTMessageIndexSet) \setminus \MTPath(\MTMessageIndexSet)$ that are not independent. Since $\MTGraph_{\MTMessageLength}$ is a tree, one must be a descendant of the other; without loss of generality, $\MTGraphVertex_{2}$ is a descendant of $\MTGraphVertex_{1}$. Since $\MTGraphVertex_{2} \in \MTCoPath(\MTMessageIndexSet)$, there exists $\MTMessageIndex \in \MTMessageIndexSet$ such that $\MTGraphVertex_{2} \in \MTCoPath(\MTMessageIndex)$, and hence the sibling $\MTGraphVertex_{2}'$ of $\MTGraphVertex_{2}$ is in $\MTPath(\MTMessageIndex)$. The vertex $\MTGraphVertex_{2}'$ is also a descendant of $\MTGraphVertex_{1}$. This tells us that $\MTGraphVertex_{1} \in \MTPath(\MTMessageIndex)$, a contradiction.
\end{proof}

\begin{claim}
\label{claim:in-paths-are-determined}
For every $(\MTLayer,\MTMessageIndex) \in \MTPath(\MTMessageIndexSet)$, $(\MTLayer+1,2\MTMessageIndex-1),(\MTLayer+1,2\MTMessageIndex) \in \MTCoPath(\MTMessageIndexSet) \cup \MTPath(\MTMessageIndexSet)$.
\end{claim}

\begin{proof}
If $(\MTLayer,\MTMessageIndex)$ is on some path from a leaf vertex to the root vertex,
one of the two children $(\MTLayer+1,2\MTMessageIndex-1)$ and $(\MTLayer+1,2\MTMessageIndex)$ of $(\MTLayer,\MTMessageIndex)$ is on that path and the other is on the corresponding copath.
\end{proof}

\begin{proof}[Proof of \Cref{lemma:mt-privacy}]
Fix $\MTMessageVector \in \MTAlphabet^{\MTMessageLength}$ and $\MTMessageIndexSet \in \binom{[\MTMessageLength]}{\MTMessageIndexSetSize}$. We argue that if the adversary $\MTAdversary$ outputs the message vector $\MTMessageVector$ and index set $\MTMessageIndexSet$ then the statistical distance between the two distributions (conditioned on this output) is at most $\sum_{(\MTLayer,\MTMessageIndex) \in \MTCoPathOptVertexSet} \MTRootHidingError(\ROOutputSize,\MTAlphabet,2^{\MTDepth-\MTLayer},\MTSaltSize,\ROQueryBound)$ where $\MTCoPathOptVertexSet = \MTCoPath(\MTMessageIndexSet) \setminus \MTPath(\MTMessageIndexSet)$. The lemma then follows by taking the maximum across all possible outputs $\MTMessageVector$ and $\MTMessageIndexSet$.

The simulator $\MTSimulator$ proceeds layer by layer from the leaf vertices towards the root vertex, setting the values of commitments in the opening proof $\MTProof$ either by sampling them via $\MTRootSimulator$ or by deriving them via the random oracle $\ROFunction$ from other commitments.

First we discuss the leaf layer. For every $\MTMessageIndex \in \MTMessageIndexSet$, the simulator samples a random salt $\MTSaltString_{\MTMessageIndex} \in \Bits^{\MTSaltSize}$ and sets the commitment $\MTVertexLabelByName{\MTDepth}{\MTMessageIndex} \DefineEqual \ROFunction(\MTMessageEntry_{\MTMessageIndex},\MTSaltString_{\MTMessageIndex})$. This is distributed exactly as in the real commitment experiment, so no simulation errors are incurred in the leaf layer.

Next we discuss the commitments $((\MTVertexLabel_{\MTGraphVertex})_{\MTGraphVertex \in \MTCoPath(\MTMessageIndex)})_{\MTMessageIndex \in \MTMessageIndexSet}=((\MTVertexLabelByCoPath{\MTMessageIndex}{\MTLayer})_{\MTLayer \in \{1,\dots,\MTDepth\}})_{\MTMessageIndex \in \MTMessageIndexSet}$.
\begin{itemize}
  \item By \Cref{claim:copaths-minus-paths-are-independent}, vertices in $\MTCoPath(\MTMessageIndexSet) \setminus \MTPath(\MTMessageIndexSet)$ are pairwise independent. Hence, for every $(\MTLayer,\MTMessageIndex) \in \MTCoPath(\MTMessageIndexSet) \setminus \MTPath(\MTMessageIndexSet)$, setting $\MTVertexLabelByName{\MTLayer}{\MTMessageIndex} \gets \MTRootSimulator^{\ROFunction}$ incurs a simulation error of $\MTRootHidingError(\ROOutputSize,\MTAlphabet,2^{\MTDepth-\MTLayer},\MTSaltSize,\ROQueryBound)$. Indeed, $(\MTLayer,\MTMessageIndex)$ is the root vertex of a perfect binary tree of depth $\MTDepth-\MTLayer$, and this perfect binary tree is disjoint from the perfect binary trees of any other vertex in $\MTCoPath(\MTMessageIndexSet) \setminus \MTPath(\MTMessageIndexSet)$. Hence, the total simulation error here is $\sum_{(\MTLayer,\MTMessageIndex) \in \MTCoPathOptVertexSet} \MTRootHidingError(\ROOutputSize,\MTAlphabet,2^{\MTDepth-\MTLayer},\MTSaltSize,\ROQueryBound)$.
  \item By \Cref{claim:in-paths-are-determined}, for every $(\MTLayer,\MTMessageIndex) \in \MTPath(\MTMessageIndexSet)$, $(\MTLayer+1,2\MTMessageIndex-1),(\MTLayer+1,2\MTMessageIndex) \in \MTCoPath(\MTMessageIndexSet) \cup \MTPath(\MTMessageIndexSet)$. Hence, the value of the commitment $\MTVertexLabelByName{\MTLayer}{\MTMessageIndex}$ is determined by the values of the commitments $\MTVertexLabelByName{\MTLayer+1}{2\MTMessageIndex-1}$ and $\MTVertexLabelByName{\MTLayer+1}{2\MTMessageIndex}$, which the simulator has either previously sampled or computed. Specifically, for correctness of the simulator's output, the value of the commitment $\MTVertexLabelByName{\MTLayer}{\MTMessageIndex}$ must equal the answer $\ROFunction(\MTVertexLabelByName{\MTLayer+1}{2\MTMessageIndex-1},\MTVertexLabelByName{\MTLayer+1}{2\MTMessageIndex})$. No simulation errors are incurred by setting values for these commitments.
\end{itemize}

Finally, for correctness of the simulator's output, the root commitment $\MTCommitment \DefineEqual \MTVertexLabelByName{0}{1}$ must equal the answer $\ROFunction(\MTVertexLabelByName{1}{1},\MTVertexLabelByName{1}{2})$. The values for both commitments are set by the simulator because $(0,1) \in \MTPath(\MTMessageIndexSet)$. No simulation error is incurred by setting the root commitment.
\end{proof}

\begin{remark}[expressions for $\MTHidingError$]
\label{remark:mtzk-plugging-in-root-hiding-bound}
The expressions for $\MTHidingError$ in \Cref{lemma:mt-privacy} are in terms of $\MTRootHidingError$. Here we describe the expression that we obtain for $\MTHidingError$ in the case $\ROQueryBound<\infty$ when we plug in the expression for $\MTRootHidingError$ in \Cref{lemma:mt-root-hiding} (for $\ROQueryBound<\infty$).

First we discuss the case of a single leaf. We get the expression $\MTHidingError(\ROOutputSize,\MTAlphabet,\MTMessageLength,\MTSaltSize,\MTMessageIndexSetSize=1,\ROQueryBound) = \MTPathHidingErrorExpression \leq (\frac{\ROQueryBound}{2^{\MTSaltSize}} + \frac{\ROQueryBound}{2^{2\ROOutputSize}}) \cdot \sum_{\MTLayer \in [\MTDepth]} 2^{\MTDepth-\MTLayer} = (\frac{\ROQueryBound}{2^{\MTSaltSize}} + \frac{\ROQueryBound}{2^{2\ROOutputSize}}) \cdot (2^{\MTDepth}-1) \leq \frac{\MTMessageLength \cdot \ROQueryBound}{2^{\MTSaltSize}} + \frac{\MTMessageLength \cdot \ROQueryBound}{2^{2\ROOutputSize}}$.

Next we discuss the case of multiple leaves, using the simplified upper bound at the end of \Cref{lemma:mt-privacy}. We get the expression $\MTHidingErrorExpression \leq \MTMessageIndexSetSize \cdot (\frac{\ROQueryBound}{2^{\MTSaltSize}} + \frac{\ROQueryBound}{2^{2\ROOutputSize}}) \cdot \sum_{\MTLayer \in [\MTDepth]} 2^{\MTDepth-\MTLayer} = \MTMessageIndexSetSize \cdot (\frac{\ROQueryBound}{2^{\MTSaltSize}} + \frac{\ROQueryBound}{2^{2\ROOutputSize}}) \cdot (2^{\MTDepth}-1) \leq \MTHidingErrorLooseExpression$.
\end{remark}

\begin{remark}[superadditive and monotonic $\MTHidingError$]
\label{remark:mtzk-superadditive-and-monotonic}
For convenience we assume that the error bound $\MTHidingError(\ROOutputSize,\MTAlphabet,\MTMessageLength,\MTSaltSize,\MTMessageIndexSetSize,\ROQueryBound)$ is \emph{superadditive} in the message length $\MTMessageLength$ and opening size $\MTMessageIndexSetSize$, which means that
\begin{equation*}
\MTHidingError(\ROOutputSize,\MTAlphabet,\MTMessageLength_{1},\MTSaltSize,\MTMessageIndexSetSize_{1},\ROQueryBound)
+
\MTHidingError(\ROOutputSize,\MTAlphabet,\MTMessageLength_{2},\MTSaltSize,\MTMessageIndexSetSize_{2},\ROQueryBound)
\leq
\MTHidingError(\ROOutputSize,\MTAlphabet,\MTMessageLength_{1}+\MTMessageLength_{2},\MTSaltSize,\MTMessageIndexSetSize_{1}+\MTMessageIndexSetSize_{2},\ROQueryBound)
\enspace.
\end{equation*}
In particular, $\MTHidingError$ is \emph{monotonic} (more precisely, nondecreasing) in the message length and opening size. This holds for the upper bound $\MTHidingError(\ROOutputSize,\MTAlphabet,\MTMessageLength,\MTSaltSize,\MTMessageIndexSetSize,\ROQueryBound) = \MTHidingErrorExpression$ in \Cref{lemma:mt-privacy}.

We use this to simplify certain expressions later on (in \Cref {section:micali-zero-knowledge,section:bcs-zero-knowledge}). However, tighter upper bounds on $\MTHidingError$ might not be monotonic (and thus also not superadditive) in the message length $\MTMessageLength$ or opening size $\MTMessageIndexSetSize$. If so one can consider the non-simplified expressions in our analyses.

For example, in the upper bound $\max_{\MTMessageIndexSet \in \binom{[\MTMessageLength]}{\MTMessageIndexSetSize}} \sum_{(\MTLayer,\MTMessageIndex) \in \MTCoPathOptVertexSet} \MTRootHidingError(\ROOutputSize,\MTAlphabet,2^{\MTDepth-\MTLayer},\MTSaltSize,\ROQueryBound)$, the size of the set $\MTCoPathOptVertexSet$ does \emph{not} monotonically increase in the size $\MTMessageIndexSetSize$ of $\MTMessageIndexSet$ (informally, $\MTCoPathOptVertexSet$ first gets bigger and then gets smaller). In the extreme setting where $\MTMessageIndexSet=[\MTMessageLength]$ (i.e., the size of $\MTMessageIndexSet$ is the largest possible value $\MTMessageIndexSetSize=\MTMessageLength)$ the set $\MTCoPathOptVertexSet$ is empty, and the hiding error is zero ($\MTHidingError(\ROOutputSize,\MTAlphabet,\MTMessageLength,\MTSaltSize,\MTMessageIndexSetSize=\MTMessageLength,\ROQueryBound)=0$).
\end{remark}

%%%%%%%%%%%%%%%%%%%%%%%%%%%%%%%%%%%%%%%%%%%%%%%%%%%%%%%%%%%%%%%%%%%%%%%%%%%%%%%
%%%%%%%%%%%%%%%%%%%%%%%%%%%%%%%%%%%%%%%%%%%%%%%%%%%%%%%%%%%%%%%%%%%%%%%%%%%%%%%
%%%%%%%%%%%%%%%%%%%%%%%%%%%%%%%%%%%%%%%%%%%%%%%%%%%%%%%%%%%%%%%%%%%%%%%%%%%%%%%
\section{Equivocation}
\label{section:merkle-commitment-equivocation}

We prove that the Merkle commitment scheme is \emph{equivocable}, a security property that informally states that, given a simulated commitment, one can (inefficiently) sample a simulated opening proof for any desired message fragment, in such a way that an adversary cannot tell the difference from an honestly computed commitment and opening proof. We rely on this property in \Cref{section:wi-for-constructions}, to when studying witness indistinguishability for constructions based on the Merkle commitment scheme.

\begin{lemma}[equivocation]
\label{lemma:mt-equivocation}
Let $\MTSymbol \DefineEqual \MTConstructor{\ROOutputSize}{\MTAlphabet}{\MTMessageLength}{\MTSaltSize}$. There exist (inefficient) probabilistic algorithms $\MTRootInefficientSimulator$ and $\MTOpeningInefficientSimulator$ such that for every opening size $\MTMessageIndexSetSize \in [\MTMessageLength]$, query bound $\ROQueryBound \in \Naturals$, and $\ROQueryBound$-query algorithm $\MTAdversary$, the following two distributions are $\MTEquivocationError(\ROOutputSize,\MTAlphabet,\MTMessageLength,\MTSaltSize,\MTMessageIndexSetSize,\ROQueryBound)$-close in statistical distance:
\begin{equation*}
\Distribution_{\RealSymbol}
\DefineEqual
\left\{
  \MTAdversary^{\ROFunction}(\ROAdvState,\MTProof)
\GivenExperiment
\StateExperiment{
  \ROFunction \gets \RODistribution{\ROOutputSize} \\
  \left(\MTMessageVector \in \MTAlphabet^{\MTMessageLength},\ROAdvState\right) \gets \MTAdversary^{\ROFunction} \\
  (\MTCommitment,\MTTrapdoor) \gets \MTCommit^{\ROFunction}(\MTMessageVector) \\
  \left(\MTMessageIndexSet \in \binom{[\MTMessageLength]}{\MTMessageIndexSetSize},\ROAdvState\right) \gets \MTAdversary^{\ROFunction}(\ROAdvState,\MTCommitment) \\
  \MTProof \gets \MTOpen^{\ROFunction}(\MTTrapdoor,\MTMessageIndexSet)
}
\right\}
\end{equation*}
and
\begin{equation*}
\Distribution_{\SimSymbol}
\DefineEqual
\left\{
  \MTAdversary^{\ROFunction}(\ROAdvState,\MTProof)
\GivenExperiment
\StateExperiment{
  \ROFunction \gets \RODistribution{\ROOutputSize} \\
  \left(\MTMessageVector \in \MTAlphabet^{\MTMessageLength},\ROAdvState\right) \gets \MTAdversary^{\ROFunction} \\
  \MTCommitment \gets \MTRootInefficientSimulator^{\ROFunction} \\
  \left(\MTMessageIndexSet \in \binom{[\MTMessageLength]}{\MTMessageIndexSetSize},\ROAdvState\right) \gets \MTAdversary^{\ROFunction}(\ROAdvState,\MTCommitment) \\
  \MTProof \gets \MTOpeningInefficientSimulator^{\ROFunction}(\MTCommitment,\MTMessageIndexSet,\MTMessageVector[\MTMessageIndexSet])
}
\right\}
\end{equation*}
where $\MTEquivocationError(\ROOutputSize,\MTAlphabet,\MTMessageLength,\MTSaltSize,\MTMessageIndexSetSize,\ROQueryBound) \leq 2 \cdot \MTMessageLength \cdot \ROPseudorandomnessError(\ROOutputSize,0,2\ROOutputSize,\ROQueryBound) + \MTMessageLength \cdot \ROPseudorandomnessError(\ROOutputSize,\log\Cardinality{\MTAlphabet},\MTSaltSize,\ROQueryBound)$. Here $\ROPseudorandomnessError$ is the regularity error of the random oracle from \Cref{lemma:ro-pseudorandomness}.
\end{lemma}

\begin{proof}
The algorithm $\MTOpeningInefficientSimulator$ relies on $\ROInverter$ from \Cref{lemma:ro-pseudorandomness} to sample commitments for all vertices in the Merkle tree, layer by layer. When we invoke $\ROInverter$ on two inputs we mean that the middle input (the target prefix) is the empty string.
\begin{itemize}[nosep]
\item[] $\MTOpeningInefficientSimulator(\MTCommitment,\MTMessageIndexSet,\MTMessageVector[\MTMessageIndexSet])$:
\begin{enumerate}[noitemsep]
\item Set $\MTVertexLabelByName{0}{1} \DefineEqual \MTCommitment$.
\item For each layer $\MTLayer = 0,1,\ldots,\MTDepth-1$ and for each $\MTMessageIndex = 1,\dots,2^{\MTLayer}$:
\begin{enumerate}[noitemsep]
\item Sample $\ROQuery \gets \ROInverter^{\ROFunction}(\MTVertexLabelByName{\MTLayer}{\MTMessageIndex},2\ROOutputSize)$.
\item Parse $\ROQuery = (\ROQuery_{\LeftSymbol},\ROQuery_{\RightSymbol})$ be the left and right half of $\ROQuery$.
\item Set $\MTVertexLabelByName{\MTLayer-1}{2\MTMessageIndex-1} \DefineEqual \ROQuery_{\LeftSymbol}$.
\item Set $\MTVertexLabelByName{\MTLayer-1}{2\MTMessageIndex} \DefineEqual \ROQuery_{\RightSymbol}$.
\end{enumerate}
\item For every $\MTMessageIndex \in \MTMessageIndexSet$, sample $\MTSaltString_{\MTMessageIndex} \gets \ROInverter^{\ROFunction}(\MTVertexLabelByName{\MTDepth}{\MTMessageIndex},\MTMessageVector[\MTMessageIndex],\MTSaltSize)$.
 \item Set the salts $\MTSaltStrings \DefineEqual (\MTSaltString_{\MTMessageIndex})_{\MTMessageIndex \in \MTMessageIndexSet}$.
 \item Set the commitments $\MTVertexLabels \DefineEqual ((\MTVertexLabelByName{\MTLayer}{\MTMessageIndex})_{\MTMessageIndex \in [2^{\MTLayer}]})_{\MTLayer=0}^{\MTDepth}$.
\item Set the opening trapdoor $\MTTrapdoor \DefineEqual (\MTSaltStrings,\MTVertexLabels)$.
\item Output $\MTProof \gets \MTOpen^{\ROFunction}(\MTTrapdoor, \MTMessageIndexSet)$.
\end{enumerate}
\end{itemize}
We define a sequence of hybrid distributions. For every $\MTLayer \in \{0,1,\ldots,\MTDepth\}$, we define a hybrid distribution $\Distribution_{\MTLayer}$ where we sample the $\MTLayer$-th layer at random, compute the Merkle tree from the $\MTLayer$-th layer to the root, and invert using $\ROInverter$ from the $\MTLayer$-th layer to the leaves:
\begin{equation*}
\Distribution_{\MTLayer} \DefineEqual
\left\{
  \MTAdversary^{\ROFunction}(\ROAdvState,\MTProof)
\GivenExperiment
\StateExperiment{
  \ROFunction \gets \RODistribution{\ROOutputSize} \\
  \left(\MTMessageVector \in \MTAlphabet^{\MTMessageLength},\ROAdvState\right) \gets \MTAdversary^{\ROFunction} \\
\text{For} \; \MTCommitmentIndex = 1, \ldots, 2^{\MTLayer},\; \MTVertexLabelByName{\MTLayer}{\MTCommitmentIndex} \gets \MTRootInefficientSimulator^{\ROFunction} \\
  (\MTCommitment,\MTTrapdoor) \gets \MTCommit_{\MTLayer}^{\ROFunction}(\MTVertexLabelByName{\MTLayer}{1},\ldots,\MTVertexLabelByName{\MTLayer}{2^{\MTLayer}}) \\
  \left(\MTMessageIndexSet \in \binom{[\MTMessageLength]}{\MTMessageIndexSetSize},\ROAdvState\right) \gets \MTAdversary^{\ROFunction}(\ROAdvState,\MTCommitment) \\
   \MTProof \gets \MTOpeningInefficientSimulator_{\MTLayer}^{\ROFunction}(\MTVertexLabelByName{\MTLayer}{1},\ldots,\MTVertexLabelByName{\MTLayer}{2^{\MTLayer}},\MTMessageIndexSet,\MTMessageVector[\MTMessageIndexSet],\MTTrapdoor)
}
\right\}
\enspace.
\end{equation*}
Above, $\MTCommit_{\MTLayer}$ receives the $\MTLayer$-th layer commitments and computes as $\MTCommit$ does, and $\MTOpeningInefficientSimulator_{\MTLayer^*}$ is defined below.
\begin{itemize}[noitemsep]
\item[] $\MTOpeningInefficientSimulator_{\MTLayer^*}(\MTCommitment,\MTMessageIndexSet,\MTMessageVector,\MTTrapdoor)$:
\begin{enumerate}[noitemsep]
\item For each layer $\MTLayer = \MTLayer^{*},\ldots,\MTDepth-1$ and for each $\MTMessageIndex = 1,\dots,2^{\MTLayer}$:
\begin{enumerate}[noitemsep]
\item Sample $\ROQuery \gets \ROInverter^{\ROFunction}(\MTVertexLabelByName{\MTLayer}{\MTMessageIndex},2\ROOutputSize)$.
\item Parse $\ROQuery$ as a pair $(\ROQuery_{\LeftSymbol},\ROQuery_{\RightSymbol})$ consisting of its left and right halves.
\item Set $\MTVertexLabelByName{\MTLayer-1}{2\MTMessageIndex-1} \DefineEqual \ROQuery_{\LeftSymbol}$.
\item Set $\MTVertexLabelByName{\MTLayer-1}{2\MTMessageIndex} \DefineEqual \ROQuery_{\RightSymbol}$.
\end{enumerate}
\item For every $\MTMessageIndex \in \MTMessageIndexSet$, sample $\MTSaltString_{\MTMessageIndex} \gets \ROInverter^{\ROFunction}(\MTVertexLabelByName{\MTDepth}{\MTMessageIndex},\MTMessageVector[\MTMessageIndex],\MTSaltSize)$.
\item Set the salts $\MTSaltStrings \DefineEqual (\MTSaltString_{\MTMessageIndex})_{\MTMessageIndex \in \MTMessageIndexSet}$.
\item Set the commitments $\MTVertexLabels \DefineEqual ((\MTVertexLabelByName{\MTLayer}{\MTMessageIndex})_{\MTMessageIndex \in [2^{\MTLayer}]})_{\MTLayer=0}^{\MTDepth}$.
\item Set the opening trapdoor $\MTTrapdoor' \DefineEqual (\MTSaltStrings,\MTVertexLabels \cup \MTTrapdoor)$.
\item Output $\MTProof \gets \MTOpen^{\ROFunction}(\MTTrapdoor', \MTMessageIndexSet)$.
\end{enumerate}
\end{itemize}

\begin{claim}
For every $\MTLayer \in \{0,1,\ldots,\MTDepth-1\}$,
$\StatisticalDistance{\Distribution_{\MTLayer}}{\Distribution_{\MTLayer+1}} \leq 2^{\MTLayer} \cdot \ROPseudorandomnessError(\ROOutputSize,0,2\ROOutputSize,\ROQueryBound)$.
\end{claim}

\begin{proof}
We rely on further hybrid distributions $\Distribution_{\MTLayer,0},\ldots,\Distribution_{\MTLayer,2^{\MTLayer}}$, where $\Distribution_{\MTLayer} = \Distribution_{\MTLayer,0}$ and $\Distribution_{\MTLayer,2^{\MTLayer}} = \Distribution_{\MTLayer+1}$. Recall that in $\Distribution_{\MTLayer}$ the $\MTLayer$-th layer is sampled at random, and layer $\MTLayer+1$ is sampled via $\ROInverter$. We define the hybrid distribution $\Distribution_{\MTLayer,\MTMessageIndex}$ so that we sample all but the first $\MTMessageIndex$ vertices of the $\MTLayer$-th layer at random; specifically, the first $\MTMessageIndex$ vertices are sampled as in $\Distribution_{\MTLayer+1}$ (we sample their children in layer $\MTLayer+1$ at random and hash them to get the values for the vertices at layer $\MTLayer$). The rest of the tree in $\Distribution_{\MTLayer,\MTMessageIndex}$ is obtained as in $\Distribution_{\MTLayer}$.

The only difference between $\Distribution_{\MTLayer,\MTMessageIndex}$ and $\Distribution_{\MTLayer,\MTMessageIndex+1}$ is how the $i$-th vertex of the $\MTLayer$-th layer (and its two children) are sampled. Either the $\MTMessageIndex$-th vertex is random and is inverted via $\ROInverter$, or the two children are random and we forward compute the $\MTMessageIndex$-th vertex. These two cases are $\ROPseudorandomnessError(\ROOutputSize,0,2\ROOutputSize,\ROQueryBound)$-close in statistical distance by \Cref{lemma:ro-pseudorandomness} ($\ROOutputSize$ is the oracle output size, $0$ denotes that there is no target prefix, $2\ROOutputSize$ is the preimage size, and $\ROQueryBound$ is the adversary query bound).

We conclude that
\begin{align*}
\StatisticalDistance{\Distribution_{\MTLayer}}{\Distribution_{\MTLayer+1}}
& \leq
\sum_{\MTMessageIndex \in \{0,1,\ldots,2^{\MTLayer}-1\}}
\StatisticalDistance{\Distribution_{\MTLayer,\MTMessageIndex}}{\Distribution_{\MTLayer,\MTMessageIndex+1}} \\ \\ & \leq
\sum_{\MTMessageIndex \in \{0,1,\ldots,2^{\MTLayer}-1\}}
\ROPseudorandomnessError(\ROOutputSize,0,2\ROOutputSize,\ROQueryBound)
\\ & \leq
2^{\MTLayer} \cdot \ROPseudorandomnessError(\ROOutputSize,0,2\ROOutputSize,\ROQueryBound)
\enspace.
\end{align*}
\end{proof}

Note that $\Distribution_{\SimSymbol}=\Distribution_{0}$. Moreover, $\StatisticalDistance{\Distribution_{\MTDepth}}{\Distribution_{\SimSymbol}} \leq \MTMessageLength \cdot \ROPseudorandomnessError(\ROOutputSize,\log\Cardinality{\MTAlphabet},\MTSaltSize,\ROQueryBound)$. This follows via a similar sequence of $\MTMessageLength$ hybrid distributions as in the above argument; the main difference is that in this case there are different target prefix and salt sizes: each target prefix $\MTMessageVector[\MTMessageIndex]$ has size $\log\Cardinality{\MTAlphabet}$ and each salt $\MTSaltString_{\MTMessageIndex}$ has salt size $\MTSaltSize$; hence the different input parameters to $\ROPseudorandomnessError$.

We conclude that
\begin{align*}
\StatisticalDistance{\Distribution_{\SimSymbol}}{\Distribution_{\SimSymbol}}
& \leq
\sum_{\MTLayer \in \{0,1,\ldots,\MTDepth-1\}}
\StatisticalDistance{\Distribution_{\MTLayer}}{\Distribution_{\MTLayer+1}}
+
\StatisticalDistance{\Distribution_{\MTDepth}}{\Distribution_{\SimSymbol}}
\\ & \leq
\sum_{\MTLayer \in \{0,1,\ldots,\MTDepth-1\}}
2^{\MTLayer} \cdot \ROPseudorandomnessError(\ROOutputSize,0,2\ROOutputSize,\ROQueryBound)
+
\MTMessageLength \cdot \ROPseudorandomnessError(\ROOutputSize,\log\Cardinality{\MTAlphabet},\MTSaltSize,\ROQueryBound)
\\ & \leq
2 \cdot \MTMessageLength \cdot \ROPseudorandomnessError(\ROOutputSize,0,2\ROOutputSize,\ROQueryBound)
+ \MTMessageLength \cdot \ROPseudorandomnessError(\ROOutputSize,\log\Cardinality{\MTAlphabet},\MTSaltSize,\ROQueryBound)
\enspace.
\end{align*}
\end{proof}
